

\chapter{Introducción a las Tuplas, Listas, Diccionarios y Conjuntos, tipo  Colecciones}{\label{Cap4}}\hypertarget{Cap4}{}
%  https://es.wikibooks.org/wiki/Manual_de_LaTeX/Inclusi%C3%B3n_de_gr%C3%A1ficos/Texto_alrededor_de_figuras
% https://realpython.com/python-lists-tuples/#lists-can-contain-arbitrary-objects
%  Tipos de datos complejos,  Tipos de contenedores
%   Eric 

\begin{wrapfigure}{l}{0.63\textwidth} 
	\vspace{-25pt}	\begin{center} \includegraphics[width=0.64\textwidth]{../Logos_de_Python/tuplas_listas}
	\end{center}\vspace{-30.5pt} %\caption{Python 3} 
\end{wrapfigure} En el capítulo anterior estudiamos los tipos de datos simples o básicos de Python, como números, cadenas de texto y los booleanos. En este capítulo trataremos un tipo de dato más complejo,  un tipo  de dato u objeto que se puede usar para  almacenar  secuencias de dato: como  tuplas,  listas y diccionarios. Estos tipos  pueden almacenar diversos objetos, diferenciándose  por su sintaxis y  la forma en que se  manipulan sus elementos.\\

Las listas y tuplas son colecciones ordenadas. Las tuplas usa paréntesis, mientras que la lista usa corchetes. Con una diferencia muy importante, las tuplas son objetos inmutables, mientra que las listas son mutables. 


Las listas funcionan bien para almacenar colecciones de elementos que pueden cambiar a lo largo del programa.
La capacidad de modificar listas es particularmente importante, por ejemplo, cuando
se trabaja con una lista de usuarios en un sitio web o una lista de personajes en un juego. Sin embargo, algunas veces se necesita crear una lista de elementos que no  cambie,  y los objetos o variables  que cumplen con esta características son las tuplas. Python se refiere a valores que no pueden cambiar como
inmutable, y una lista inmutable se llama tupla. 

Por otro lado, las listas  sería equivalente a lo que en otros lenguajes se conoce por arreglos (arrays), o vectores.\\


%  de Bill..   Listas y tuplas La mayoría de los lenguajes informáticos pueden representar una secuencia de elementos indexados por su número entero.  posición: primero, segundo, y así hasta el último. Ya has visto cadenas de Python, que son secuencias de personajes. También has tenido una pequeña vista previa de las listas, que podrás Ahora vemos que son secuencias de cualquier cosa. Python tiene otras dos estructuras de secuencia: tuplas y listas. Estos contienen cero o más elementos. A diferencia de las cadenas, los elementos pueden ser de diferentes tipos. De hecho, cada elemento puede ser cualquier objeto de Python. Esto le permite crear estructuras tan profundas y complejas como desee.  como. ¿Por qué Python contiene listas y tuplas? Las tuplas son inmutables; cuando asignas Los elementos de una tupla están integrados en el pastel y no se pueden cambiar. Las listas son mutables, lo que significa que puedes insertar y eliminar elementos con gran entusiasmo. mostraré muchos ejemplos de cada uno, con énfasis en las listas.



%  Un diccionario de Python puede contener pocos o grandes cantidades de datos (millones). Pueden almacenar información en una variedad de formas en pares,  por lo que existen diferentes enfoques del bucle para recorrer todos los pares clave-valor de un diccionario, puede ser a través de sus claves o a través de sus valores.






comentario de diccionario ...\\






\section{Introducción a las Tuplas}{\label{tuplas01}}	% tomado de  Eric Matthes


\begin{wrapfigure}{l}{0.52\textwidth} 
	\vspace{-25pt}	\begin{center} 
\includegraphics[width=0.53\textwidth]{../Logos_de_Python/tuplas01}
	\end{center}\vspace{-30.5pt} %\caption{Python 3} 
\end{wrapfigure}  
Una tupla  es una colección o sucuencia ordenada; un contenedor (variable) que permite almacenar varios datos de diferentes tipos, pero es   inmutable, significa una vez creada la tupla  no puede ser modificado. Su sintaxis es sencilla (parecida a la listas);  consiste en incluir los valores entre paréntesis y separados por comas. Para definir 
una tupla con un único elemento, es necesario incluir una coma al final, ejemplo: {\ttfamily tupla  = (3,)} generalmente no tiene sentido construir una tupla con un sólo  elemento, pero esto puede suceder cuando las tuplas se generan automáticamente.\\

Veamos algunos ejemplos de tuplas:

\begin{lstlisting}%[frame=single]
	
	In[]:
	
	>>> tupla = (2, "texto", 2.4, True, 2j)
	>>> tupla
	(2, 'texto', 2.4, 2j)
\end{lstlisting}
En este ejemplo la tupla esta formada por un entero, un string, un float, booleano y un número complejo. Una vez definida no es posible hacer cambios en la tupla. Pero se puede acceder a sus elementos individuales mediante
el índice de cada elemento; el operador índice en tuplas usa la misma metodología que usan las cadenas (ver tabla \ref{Tabla02}), y aún más, este operador también se usa en listas. El índice  inicia en cero, seguido por uno, dos y así sucesivamente hasta alcanzar el último. 

Si estamos interesados en conocer las longitud o tamaño de la tupla, podemos usar la función {\ttfamily len()}, que se introdujo en la sección \ref{longitud} dedicada a cadenas, en este caso se usa para calcular el tamaño o longitud  de una tupla, el proceso es análogo al de las cadenas: 
\begin{lstlisting}%[frame=single]
	
	In[]:
	
	>>> len(tupla)
	5 	
\end{lstlisting}
La función indica que el tamaño o longitud de la tupla, que en este caso  es 5; si embargo, debe tener en cuenta que sus elementos son indizados a partir de cero (y no en 1) y como consecuencia en este ejemplo el último queda indizado con 4. 

La sintaxis para acceder a cada uno de los elementos de la tupla es sencilla:  se  escribe el nombre de la lista y se  indica el índice del elemento entre corchetes. Ejemplo:

\begin{lstlisting}%[frame=single]
	
	In[]:
	
	>>> tupla[1]
	'texto'		
\end{lstlisting}
Observe, el índice cero corresponde al entero 2, el índice uno a la cadena 'texto', y el indice cuatro al valor complejo 2j.\\

Las tuplas también utilizan índices negativos, y su sintaxis es igual a las cadenas; iniciando en cero  y finalizando en -1 (posición final), en el ejemplo anterior {\ttfamily  tupla[-4] = tupla[0]= 2 }.
Si se utiliza un número negativo como índice, esto se traduce en que el índice empieza a contar desde el final, hacia la izquierda; es decir, con [-1] accederíamos al último elemento de la lista, con [-2] al penúltimo, con [-3], al antepenúltimo, y así sucesivamente. De esta forma sin conocer la longitud de la tupla se puede extraer su último elemento con -1:     
\begin{lstlisting}%[frame=single]
	
	In[]:
	
	>>> tupla[-1]
	2j	
\end{lstlisting}

Este operador {\ttfamily [i]} de Python se utiliza en string, tuplas, y listas, en todos estos objetos se usa de manera semejante, teniendo en cuenta que string y tuples son inmutables, y las listas son mutables. 

Con operador {\ttfamily []}  también se puede acceder o extraer una porción de la tupla, indicando el  índice de inicio y el índice  final, este procedimiento es análogo al tratamiento que recibe las cadenas.

\begin{lstlisting}%[frame=single]
	
	In[]:
	
	>>> tupla[1:4]
	('texto', 2.4, True)
	
	>>> tupla[3:]
	(True, 2j)
	
	>>> tupla[:2]
	(2, 'texto')	
\end{lstlisting}

Al igual que en  las cadenas, la otra forma de acceder a la tupla de forma inversa (de atrás hacia adelante), es con índices negativos, y el paso o salto positivo o negativo:

\begin{lstlisting}%[frame=single]
	
	In[]:
	
	# Para acceder a la tupla de forma inversa	
	>>> tupla[-1::-1]
	(2j, True, 2.4, 'texto', 2)
	
	# El mismo resultado anterior
	>>> tupla[::-1]
	(2j, True, 2.4, 'texto', 2)	
	
	>>> tupla[-2:]
	(True, 2j)
	
	>>> tupla[:-4]
	(2,)
	
	>>> tupla[::2]	
	(2, 2.4, 2j)
\end{lstlisting}

Observe, la tupla {\ttfamily  tupla[:-4] = (2,)}, es una tupla con un único elemento, longitud igual a uno.
Y la tupla {\ttfamily tupla[::2] } que va desde el inicio de dos en dos, por el salto indicado.


Al igual que en las cadenas con índice positivos el último no se incluye (ver tabla \ref{Tabla02}), debido al indizado que se  inicia en cero:
\begin{lstlisting}%[frame=single]
	
	In[]:
	
	>>> tupla[:4]
	(2, 'texto', 2.4, True)
\end{lstlisting}





\subsection{Ejercicios}

\begin{enumerate}
	
	\item Buffet: Un restaurante estilo buffet ofrece sólo cinco alimentos básicos. Piensa en cinco
	alimentos simples y guárdelos en una tupla.
	
	Utilice un bucle {\ttfamily{for}} para imprimir cada comida que ofrece el restaurante.
	
	Intente modificar uno de los elementos y asegúrese de que Python rechace el	cambiar.
	
	El restaurante cambia su menú, reemplazando dos de los artículos con diferentes
	alimentos. Agregue una línea que reescriba la tupla y luego use un bucle {\ttfamily{for}} para imprimir
	cada uno de los elementos del menú revisado.
	
	
	
\end{enumerate}





\section{Introducción a las Listas}{\label{Listas01}}	% tomado de  Eric Matthes


\begin{wrapfigure}{l}{0.46\textwidth} 
	\vspace{-25pt}	\begin{center} 
\includegraphics[width=0.48\textwidth]{../Logos_de_Python/Listas05}
	\end{center}\vspace{-35pt} %\caption{Python 3} 
\end{wrapfigure}  

En esta sección estudiaremos la variable  listas, objeto que hemos mencionado en secciones anteriores como  información general.  La lista es un tipo de  colección o secuencia ordenada, semejante a las tuplas, pero con la particularidad de ser mutables. Lo que significa que  pueden ser cambiada en cualquier momento del programa. 
Las listas sería la herramienta equivalente a lo que en otros lenguajes se conoce por arreglos (arrays), o vectores. Las listas son una de las funciones más poderosas de Python y con gran facilidad para acceder; pueden contener cualquier tipo de dato: números, cadenas, booleanos, complejos, y inclusive también listas. En Python, los corchetes ({\ttfamily[]}) indican una lista y los elementos de la lista están separados por comas. \\



Las listas un tipo de secuencia ordenada muy buenas para realizar un seguimiento de procesos según su orden, especialmente cuando el orden y el contenido puede cambiar. Permiten almacenar conjuntos de información, con
unos pocos  o millones de datos. A diferencia de las cadenas y tuplas, las listas son mutables. Una lista puede cambiar agregando nuevos elementos, eliminando o sobrescribiendo los elementos existentes; además  un mismo valor puede ocurrir más de una vez en una lista.



Como se señaló antes, crear una lista muy  sencillo, semejante a crear una tupla. En este caso  los valores que queremos incluir en la lista se colocan entre corchetes separados por comas, veamos el siguiente ejemplo:
\begin{lstlisting}%[frame=single]
	
	In[]:
	
	>>> lista = [22, True, 2.35, "una cadena", 2j,  [3, 5, False, "otro"]]
	>>> lista
	[22, True, 2.35, 'una cadena', 2j, [3, 5, False, "otro"]]
\end{lstlisting}

Aquí tenemos una lista formada por un entero, un valor booleano, un float, una cadena, un complejo y una lista de longitud 4.\\ 

Si le pide a Python que imprima una lista, con {\ttfamily print()}, Python devuelve la representación de la
lista, incluidos los corchetes:
\begin{lstlisting}%[frame=single]
	
	In[]:
	>>> prin(lista)
	
	[22, True, 2.35, 'una cadena', 2j, [3, 5, False, "otro"]]
\end{lstlisting}

\subsection{Acceder a los elementos de una lista}

Para acceder a los elementos de una lista,  el procedimiento es análogo al  caso  de las tuplas.  Acceder a cada uno de los elementos de la lista escribiendo el nombre de la lista e indicando el índice del elemento entre corchetes:

\begin{lstlisting}%[frame=single]
	
	In[]:
	
	>>> lista[3]
	"una cadena"
\end{lstlisting}

El  operador {\ttfamily []} funciona para índices en listas  de manera análogo como es tratado en tuplas y cadenas;  se inician en cero, uno, dos y así sucesivamente (ver en tabla \ref{Tabla02}), y recuerde que  el índice del primer elemento de la lista es 0, y no 1.  Y como es de esperar las listas también   soportan los índices negativos, incluyendo el argumento salto o paso. Un  índice negativo indica  que empieza a contar del final hacia la izquierda; es decir, con [-1] accederíamos al último elemento de la lista, con [-2] al penúltimo, con [-3], al antepenúltimo, y así sucesivamente:

\begin{lstlisting}%[frame=single]
	
	In[]:
	
	>>> lista[-1]
	[3, 5, False, 'otro']
	
	>>> lista[5]
	[3, 5, False, 'otro']
	
\end{lstlisting}

Si se necesita conocer el tamaño o longitud de la lista, se puede usar la función {\ttfamily len()}  con la misma metodología usada para cadenas y tuplas. En este caso, el tamaño de la lista es igual a 6, y por supuesto finaliza en -1, el resultado es el mismo en ambos casos:

\begin{lstlisting}%[frame=single]
	
	In[]:
	
	>>> len(lista)
	6
\end{lstlisting}


\noindent La longitud una lista puede ser cero, en esta caso se trata de una lista vacía: 

\begin{lstlisting}%[frame=single]
	
	In[]
	>>> lista_vacia = [ ]
	[]
	>>> len(lista_vacia)
	0
\end{lstlisting}

\noindent Y una lista con un único elemento:  

\begin{lstlisting}%[frame=single]
	
	In[]
	
	>>> l_unico = [3] 
	[3]
	
	>>>len(l_unico) 
	1
\end{lstlisting}

Observe, una tupla con un sólo elemento se escribe el elemento seguido de una coma entre paréntesis {\ttfamily (3,)}, mientras que la lista solo se escribe el elemento entre corchetes.


\subsubsection{Porciones o cortes de lista}\label{rebanada}

También podemos acceder a porciones (rebanadas) o partes  de una lista indicando los índices inicial y  final ({\ttfamily lista[inicio:final]}), o con el argumento de paso ({\ttfamily lista[inicio:fin:paso]}); todo esto siguiendo  la metodología para tuplas y cadenas. Python interpretará que queremos una lista que vaya desde la posición inicial a la posición final, sin incluir este último; y si usa el argumento paso, genera los elementos paso a paso:

\begin{lstlisting}%[frame=single]
	
	In[]
	
	# Tamaño de lista igual a 6
	lista[0:5]
	[22, True, 2.35, 'una cadena', 2j]
	
	>>> lista[1:4]
	[True, 2.35, 'una cadena'] 
	
	>>> lista[-1:]
	[[3, 5, False, 'otro']]
	
\end{lstlisting}

\noindent En caso de escribir los tres argumentos (inicio:fin:salto) en lugar de dos, el tercero se utiliza para determinar cada cuantos pasos se añade un elemento a la lista:

\begin{lstlisting}%[frame=single]
	
	In[]:
	
	>>> lista[1:4:2]
	[True, 'una cadena'] 
	
	# lista invertida
	>>> lista[::-1]
	[[3, 5, False, 'otro'], 2j, 'una cadena', 2.35, True, 22]
	
	# Salto dos en dos
	>>> lista[5::-2]
	[[3, 5, False, 'otro'], 'una cadena', True]
	
\end{lstlisting}

Y una situación más se presenta cuando tenemos una lista como elemento en la lista (una lista incluida dentro de otra lista); si necesitamos  acceder a un elemento de la lista interna, tendremos que utilizar dos veces este operador, primero para indicar en qué posición de la lista se encuentra la lista interna, y luego el índice deseado en la lista interna para acceder a sus elementos: 

\begin{lstlisting}%[frame=single]
	
	In[]:
	
	# El valor en la posición 2 de la lista interna
	>>> lista[5][2]
	False
	
	# Sin conocer la posición final, sabemos que finaliza en -1.
	>>> lista[-1][2]
	False
\end{lstlisting}




\subsubsection{La función split()}   

%Confundido con tupla con lista:
%Puede cambiar el valor de cualquier elemento de una lista, no solo el primer elemento.\\


La función  {\ttfamily split()} se usó en la sección \ref{cadena01} para  cortar o  separar cadenas mediante un separador (algún carácter en la cadena); en el caso de listas, esta función también se usa para separar una cadena y convertirla en una  lista usando un separador. Para ilustrar esta aplicación consideremos la siguiente cadena: {\ttfamily "{}Ya hemos avanzado en este curso, y vienen temas interesantes{}"{}}.

\begin{lstlisting}
	
	In[]:
	
	>>> cadena = "Ya hemos avanzado en este curso, y vienen temas interesantes"
	>>> cadena.split(" ")
	['Ya', 'hemos', 'avanzado', 'en', 'este', 'curso,', 'y', 'vienen', 'temas', 
	'interesantes']
\end{lstlisting} En este caso se usa como separador el espacio en blanco en la cadena, para formar los elementos de la lista. Consideremos un segundo ejemplo, en este caso se usa como separador el {\ttfamily\textbackslash}:

\begin{lstlisting}
	
	In[]:
	
	>>> aniversario = '06/09/2024'
	>>> aniversario.split('/')
	['06', '09', '2024']
\end{lstlisting}


En caso de  no especificar un separador la función usa por defecto espacios en blanco: nuevas líneas, espacios o tabulaciones.








%%%%%%%%%%%%%%%%%%%%%%%%%%%%   realpython  Real Python   %%%%%%%%%%%%%%%%%%%%%%%%%%%%%%%%%%%%%%%%%%%%%%

% https://realpython.com/python-variables/

\subsubsection{Modificar, añadir y eliminar elementos de una lista}   %  Eric 

La mayoría de las listas creadas serán dinámicas; es decir, a medida que el programa sigue su curso se hacen necesario actualizarlas, mediante un proceso de cambios como  modificar, eliminar o agregar elementos. Por ejemplo, podrías crear un juego en el que un jugador tenga que disparar a los extraterrestres. Se inicia una lista con un número  de alienígenas definido y luego cada vez que uno es derribado se  elimina  de la lista; y cada vez que un nuevo extraterrestre aparece en la pantalla, se  agrega a la lista. Esta lista de extraterrestres aumentará y disminuye en tamaño a lo largo del juego.



\subsubsection{Modificar elementos en una lista}


Para modificar   un elemento  de una  lista se usa el operador {\ttfamily[]},  una  operación exclusiva  para listas.  La  sintaxis para cambiar un valor en la lista es la siguiente: {\ttfamily lista[índice] = {\bf{valor}}},  significa que le asigna en la posición índice de la lista un  nuevo  {\bf{valor}}, ejemplo:

\begin{lstlisting}
	
	In[]:
	
	# lista = [22, True, 2.35, 'una cadena', 2j, [3, 5, False, 'otro']]
	>>> lista[0] = False
	>>> lista
	[False, True, 2.35, 'una cadena', 2j, [3, 5, False, 'otro']]
\end{lstlisting}
En este caso se modifico la lista cambiando el entero 22 (elemento en la posición cero) por el valor booleano {\ttfamily False}; la nueva lista modificada resulta: \\

{\ttfamily\scriptsize\scriptsize 
	\hspace{-0.65cm}Out[]:\vspace{-0.3cm}	
	{\color{blue} 	
		\begin{verbatim}
			
			[False, True, 2.35, 'una cadena', 2j, [3, 5, False, 'otro']]
		\end{verbatim}
}}

%El uso de los corchetes para acceder y modificar los elementos de una lista es común en muchos lenguajes, pero Python nos permite  varias operaciones, como aumentar el tamaño de la lista. 

Mediante la siguiente sintaxis se puede cambiar varios elementos de una lista:  {\ttfamily  lista[inicio:fin] = [valor,..., valor]}; significa que los valores de la lista en la posición inicial indicada y hasta la posición final indicada se cambian por los valores de la lista a la derecha de la asignación,  por  ejemplo: 

\begin{lstlisting}
	
	In[]:
	
	>>> lista[1:3] = ["dos", "tres"]
	lista
	[22, 'dos', 'tres', 'una cadena', 2j, [3, 5, False, 'otro']]
	
\end{lstlisting} En este ejemplo se cambian los elementos en la posición 1 y 2 en la lista (el último no se incluye) por los elementos de la lista indicada a la derecha.\\



Además se puede modificar el tamaño de la lista; agregando elementos a la derecha de la lista creando los indices nuevos a partir del último elemento,   mediante una asignación de una lista en la parte izquierda:
\begin{lstlisting}
	
	In[]:
	
	>>> lista
	[22, True, 2.35, "una cadena", 2j, [3, 5, False, "otro"]]
	>>> len(lista)
	6
	
	>>> lista[6:8] = ["seis", "siete"]
	lista
	[22, True, 2.35, 'una cadena', 2j, [3, 5, False, 'otro'], 'seis', 'siete']
	>>> len(lista)
	8
\end{lstlisting}
El tamaño de la lista inicial era igual a 6, y con este procedimiento se aumento a 8.\\


En todo caso las listas ofrecen otros mecanismos más cómodos para ser modificadas a través de las funciones llamadas clase o y  métodos. Por ahora  no trataremos estos mecanismos, pues hace falta conocer  clases y  métodos, y estas definiciones se estudiaran en el Capitulo \ref{Cap9}.\\

% Todo lo que hemos explicado sobre las listas se aplica también a las tuplas, a excepción de la forma de definirla, para lo que se utilizan paréntesis en lugar de corchetes.

\subsubsection{La función list()}  
%  al igual que la funcion {\ttfamily string()} que converte una variable en un a cadena, en cadena, 
La función {\ttfamily list()} permite convertir algunos objetos en lista, como  cadena y  tupla. 
La cadena la convierte en una lista de cadenas de un carácter, veamos el siguiente ejemplo:   

\begin{lstlisting}
	
	In[]:
	
	# list() convierte una cadena, en una lista de cadenas de un carácter:	
	>>> list("una cadena")
	['u', 'n', 'a', ' ', 'c', 'a', 'd', 'e', 'n', 'a']
\end{lstlisting} En este ejemplo la función {\ttfamily list() } convierte una cadena en una lista de cadenas de un carácter; observe que los espacios en blanco igualmente lo representa como una cadena de un espacio en blanco.\\

La función {\ttfamily list()}  también  convierte una tupla en una lista, donde cada elemento de la tupla es un elemento en la lista en el mismo orden de posición. Aquí, podemos decir que una tupla es una lista inmutable.  En el siguiente ejemplo se convierte una tupla en lista:

\begin{lstlisting}
	
	In[]:
	
	>>> tupla = ("string", 2, 0.25, True, 3j)
	>>> list(tupla)
	["string", 2, 0.25, True, 3j]
	
\end{lstlisting}




\subsubsection{Agregar elementos a una lista}

Es posible que se necesite agregar un nuevo elemento a una lista por muchas razones. Por ejemplo, es posible que se necesite  hacer que aparezcan nuevos alienígenas en un juego, agregar nuevos datos a una visualización o agregar nuevos usuarios registrados a un sitio web que haya construido. 

Python proporciona varias formas de agregar nuevos datos a listas existentes.


\subsubsection{Agregar elementos al final de una lista}

La forma más sencilla de añadir  un nuevo elemento a una lista es mediante el método  {\ttfamily .append()}. Cuando usa este método,  se  agrega el nuevo elemento  al final de la lista. Consideremos la {\ttfamily lista = ["{}Luis", "{}Antonio", "{}José", "{}Carlos"{}]}, agregaremos un nuevo elemento {\ttfamily'Jesús'} al final de la lista:

\begin{lstlisting}
	
	In[]:
	
	>>> lista = ["Luis", "Antonio", "José", "Carlos"]
	
	>>> lista
	['Luis', 'Antonio', 'José', 'Carlos']
	
	>>> lista.append('Jesús')
	>>> lista
	['Luis', 'Antonio', 'José', 'Carlos','Jesús']
	
\end{lstlisting} Aquí el método {\ttfamily append()} añade a {\ttfamily 'Jesús'} al final de la lista, sin
afectar  otro elemento de la lista.


El método {\ttfamily append()} facilita la creación de listas dinámicas. Por ejemplo,
puede comenzar con una lista vacía y luego agregar elementos a la lista usando una serie
de llamadas  {\ttfamily append()}. Ilustremos con un ejemplo, usando una lista vacía, y agreguemos los elementos: 

\begin{lstlisting}
	
	In[]:
	
	>>> estudiantes = []
	>>> estudiantes.append('Luis'), 
	>>> estudiantes.append('Antonio')
	>>> estudiantes.append('José')
	>>> print(estudiantes)
	
	['Luis', 'Antonio', 'José']
	
\end{lstlisting}  La lista resultante tiene exactamente el mismo aspecto que las listas de la versión anterior.\\


Una nota importante, el método  {\ttfamily append()} agrega un elemento a la lista, si el elemento es una lista, se agregará una lista. Para ejemplificar consideremos el archivo {\ttfamily agregar.py}: 

\hspace{-2.3cm}{\ttfamily \scriptsize agregar.py}\vspace{-0.5cm}
\begin{lstlisting}		
	otros =["Oswaldo", "manuel"]
	estudiantes.append(otros)   # agrega la lista y no elementos
	print(estudiantes)	
\end{lstlisting}	


{\ttfamily\scriptsize\scriptsize 
	\hspace{-0.65cm}Out[]:\vspace{-0.05cm}	
	{\color{blue} 	
		\begin{verbatim}			
			['Luis', 'Antonio', 'José', ['Oswaldo', 'manuel']]
		\end{verbatim}
}}

Crear listas de esta manera es muy común, porque a menudo no sabrás
los datos que sus usuarios desean almacenar en un programa hasta después de que el programa sea ejecutado. 
Para que sus usuarios tengan el control, comience definiendo una lista vacía que
almacenará  los valores que los usuarios ingresen. Luego agregue cada nuevo valor proporcionado a la lista
creada.


\subsubsection{Dos alternativas más para agregar elementos al final de una lista}  % Bill

Otras dos formas para fusionar una lista con otra, es el método {\ttfamily extend()}  y el operador {\ttfamily +=}. 
Para exponer el método  {\ttfamily extend()}, consideremos la lista anterior y una nueva lista llamada otros:
\begin{lstlisting}
	
	In[]:
	
	>>> lista = ["Luis", "Antonio", "José", "Carlos"]
	>>> otros =["Oswaldo", "manuel"]
	
	>>> lista.extend(otros)  # agrega los elementos a la lista
	>>> print(lista)
	>>> ['Luis', 'Antonio', 'José', 'Carlos', 'Oswaldo', 'manuel']
\end{lstlisting} 



Y la otra  alternativa  para fusionar  listas es el operador {\ttfamily	+=}, en este caso el operador agrega los elementos de la lista a sumar:

\begin{lstlisting}
	
	In[]:
	
	>>> lista = ["Luis", "Antonio", "José", "Carlos"]
	>>> otros =["Oswaldo", "manuel"]
	
	>>> lista += otros    # agrega los elementos a la lista
	>>> print(lista)
	>>> ['Luis', 'Antonio', 'José', 'Carlos', 'Oswaldo', 'manuel']
\end{lstlisting} 



Nota, si agrega una cadena con este procedimiento, el operador agregará los caracteres de  cadena como elementos  
en la lista.


\subsection{Insertar elementos en una lista en cualquier posición}

Puede agregar un nuevo elemento en cualquier posición de su lista se usa el método {\ttfamily insert()}. 
Esto se hace especificando el índice y el valor del nuevo elemento:

\begin{lstlisting}
	In[]:
	
	>>> estudiantes = ["Luis", "Antonio", "José", "Carlos"]
	>>> estudiantes
	['Luis', 'Antonio', 'José', 'Carlos']
	
	#  Observe se agrega un nuevo elemento a la lista
	>>> estudiantes.insert(0, 'Eduardo')
	>>> estudiantes
	['Eduardo', 'Luis', 'Antonio', 'José', 'Carlos']
	
\end{lstlisting} En este ejemplo, se insertó  el valor 'Eduardo' al inicio de la
lista. El  método {\ttfamily insert()} abre un espacio en la posición 0 y almacena el valor
'Eduardo' en esa ubicación.


\subsection{Eliminar elementos de una lista}

Algunas veces  será necesario eliminar un elemento o un conjunto de elementos de una lista. Por ejemplo, en el juego de alienígenas, 
cuando un jugador derriba a un extraterrestre, lo más probable es que debes eliminarlo de la lista de extraterrestre activos, o cuando un usuario decide cancelar su cuenta en una aplicación web que usted creó,  debes  eliminarla de la lista  de los usuarios  activos.

Puedes eliminar un elemento según su posición en la lista o según su valor. Si conoce la posición del elemento en la lista, se puede usar la  declaración {\ttfamily del}  para eliminar el elemento. La sintaxis es la siguiente:

\begin{lstlisting}
	
	In[]:
	
	>>> estudiantes = ["Luis", "Antonio", "José", "Carlos"]
	>>> del estudiantes[0]
	
	>>> estudiantes
	['Antonio', 'José', 'Carlos']
\end{lstlisting} Aquí utilizamos la instrucción {\ttfamily del} para eliminar el primer elemento, "Luis", de la
lista  estudiantes. 

Ahora eliminamos el elemento en la posición  uno de la lista, y sólo queda {}"{}Antonio"{}  y {}
"{}Carlos"{} en la lista,  pues "{}Luis"{} y {}"{}José"{} se eliminaron.


\begin{lstlisting}
	In[]:
	
	>>> del estudiantes[1]
	
	>>> estudiantes
	['Antonio', 'Carlos']
\end{lstlisting}Los elementos eliminados ya  no están en la lista.

\subsubsection{Eliminar un elemento usando el método pop()}

A veces es necesario usar el valor de un elemento  después de eliminarlo de una
lista. Por ejemplo, es posible que desee obtener la posición {\ttfamily x} e {\ttfamily y} de un 
extraterrestre que acaba de ser derribado, para provocar una explosión en esa posición. En una aplicación en la red, es posible que desee eliminar un usuario de una lista de miembros activos, y luego agregarlo  a una lista de miembros inactivos.

El método {\ttfamily pop()} elimina el último elemento de una lista, pero te permite trabajar
con ese elemento después de retirarlo. 
El término pop surge de pensar en enumerado como una pila de elementos y sacar un elemento de la parte superior de la pila.  Por analogía, la parte superior de una pila corresponde al final de una lista.
\begin{lstlisting}
	In[]:
	
	>>> estudiantes = ["Luis", "Antonio", "José", "Carlos"]
	>>> print(estudiantes.pop())
	'Carlos'
	>>> estudiantes
	['Luis', 'Antonio', 'José']
\end{lstlisting} El resultado muestra que el valor {}"{}Carlos"{} se eliminó del final de
la lista, la sentencia {\ttfamily estudiantes.pop()}  registra el valor de {}"{}Carlos"{}. Debe guardarse en una  variable, digamos {\ttfamily pop\_estudiante = estudiantes.pop()}, porque cada vez que se ejecute eliminará el último en la lista resultante.  \\

Un ejemplo de cómo se puede utilizar el método {\ttfamily pop()}; imagínate que el registro de estudiantes en la lista se almacenan en orden de llegada, podemos usar el método {\ttfamily pop()} para imprimir una declaración indicando el último en llegar, veamos con un ejemplo:


\begin{lstlisting}
	In[]:
	
	>>> estudiantes = ["Luis", "Antonio", "José", "Carlos"]
	>>> ultimo\_estudiante = estudiantes.pop()
	>>> print(f"El último estudiante en llegar es {ultimo_estudiante.title()}.")
	El último estudiante en llegar es Carlos.
\end{lstlisting}



\subsubsection{Eliminar un elemento usando el valor}\label{remueve}

A veces no sabrás la posición del valor que deseas eliminar de una lista. Si sólo conoce el valor del artículo que desea eliminar, puede utilizar el  método {\ttfamily remove()}.

Por ejemplo, digamos que queremos eliminar el valor {}"{}Antonio"{} de la lista de
estudiantes:

\begin{lstlisting}
	In[]:
	
	>>> estudiantes = ["Luis", "Antonio", "José", "Carlos"]
	>>> estudiante = estudiantes.remove("Antonio")
	>>> print(estudiantes)
	['Luis', 'José', 'Carlos']
\end{lstlisting} Aquí el método {\ttfamily remove()} le indica  a Python que busque dónde está  el elemento  {}"{}Antonio"{} que aparece en la lista y elimínelo.\\


%También puedes usar el método {\ttfamily remove()} para trabajar con un valor que se está eliminado de una lista. Eliminemos el valor "Antonio" e imprimamos un motivo para eliminándolo de la lista:

\vspace{0.2cm}
\begin{minipage}[t]{0.2\textwidth}
	\centering\raisebox{\dimexpr\topskip-\height}{%
		
\includegraphics[scale=0.20]{../Logos_de_Python/Sugerencia}}  %[width=\textwidth]
	
\end{minipage}%\hfill
\begin{minipage}[t]{0.70\textwidth} 
	{\scriptsize 
		El método {\ttfamily remove()} elimina sólo la primera aparición del valor que se especifica. Si existe la posibilidad de que el valor aparezca más de una vez en la lista, necesitarás usar un bucle para asegurarse que se eliminen
		todas las apariciones del valor. Pruébelo en un archivo.py.
		\begin{lstlisting}	
			student = ["Luis","Antonio","Carlos","José","Carlos"]
			student.remove("Carlos")
			print(student)
			['Luis', 'Antonio', 'José', 'Carlos']
		\end{lstlisting}	
	}
\end{minipage}
\vspace{0.2cm}



\subsection{Extraer elementos de cualquier posición de una lista}

Puede usar el  método {\ttfamily pop()} para eliminar un elemento en cualquier posición en una lista, incluyendo
el índice del elemento que desea eliminar entre paréntesis:

\begin{lstlisting}
	In[]:
	
	>>> estudiantes = ["Luis", "Antonio", "José", "Carlos"]
	>>> primer_estudiante = estudiantes.pop(0)
	>>> print(f"El  primer estudiante en llegar es {primer_estudiante.title()}.")
	El  primer estudiante en llegar es Luis.
\end{lstlisting}

Recuerda que cada vez que usas  el  método {\ttfamily pop()}, el elemento con el que trabajas no se almacenan en la lista.

%%%%%%%%%%%%%%%%%%%%%%%%
\vspace{0.2cm}
\begin{minipage}[t]{0.2\textwidth}
	\centering\raisebox{\dimexpr\topskip-\height}{%
		
\includegraphics[scale=0.20]{../Logos_de_Python/Advertencia}}  %[width=\textwidth]
	
\end{minipage}%\hfill
\begin{minipage}[t]{0.70\textwidth} 
	{\scriptsize 
		Si no está seguro de utilizar la declaración {\ttfamily del} o el método {\ttfamily pop()}, aquí tienes una forma sencilla de decidir: Cuándo quieres eliminar un elemento de una lista y no usar ese elemento de ninguna manera, use la declaración {\ttfamily del}; y si quieres usar un elemento a medida que lo elimina, utilice el método  {\ttfamily pop()}.}
\end{minipage}
\vspace{0.2cm}






\subsection{Organizar una lista}

Con frecuencia, sus listas se crearán en un orden impredecible, porque no siempre podemos controlar el orden en que los usuarios proporcionan sus datos. Aunque esto es inevitable en la mayoría de las circunstancias, algunas veces se necesita un orden particular para presentar su información, y en otras ocasiones se debe preservar el orden original de su lista, o cambiar el orden original. Python proporciona varias formas diferentes para organizar las listas, dependiendo de la situación.

\subsubsection{Ordenar lista de forma permanente con {\ttfamily sort()}} 

El método {\ttfamily sort()}  de Python hace que sea relativamente fácil ordenar una lista. Suponga  que tenemos una lista de autos el letras en minúsculas, y quieres cambiar el orden de la lista para almacenarlos alfabéticamente:

\hspace{-2.1cm}{\ttfamily \scriptsize autos.py}\vspace{-0.56cm}
\begin{lstlisting}	
	autos = ['ford','bmw', 'chevrolet', 'toyota']
	autos.sort()
	print(autos)
\end{lstlisting}	

El método {\ttfamily sort()}  cambia el orden de la lista de forma permanente.

{\ttfamily\scriptsize\scriptsize 
	\hspace{-0.65cm}Out[]:\vspace{-0.3cm}	
	{\color{blue} 	
		\begin{verbatim}
			['bmw', 'chevrolet', 'ford', 'toyota']
		\end{verbatim}
}}

\noindent Los autos ahora están en orden alfabético, y ya no podemos volver al orden original.\\


También se  puede ordenar la lista en orden alfabético inverso, pasando el
argumento {\ttfamily reverse=True} al método {\ttfamily sort()}. El siguiente ejemplo se ordena
la lista autos en orden alfabético inverso:
\begin{lstlisting}	
	autos.sort(reverse=True)
	print(autos)
\end{lstlisting}	

\noindent El método cambia el orden, ahora la lista se muestra en orden alfabético inverso:\\

{\ttfamily\scriptsize\scriptsize 
	\hspace{-0.65cm}Out[]:\vspace{-0.3cm}	
	{\color{blue} 	
		\begin{verbatim}	
			['toyota', 'ford', 'chevrolet', 'bmw']
		\end{verbatim}
}}
\noindent Nuevamente, el orden de la lista cambia de forma  permanente.



\vspace{0.2cm}
\begin{minipage}[t]{0.2\textwidth}
	\centering\raisebox{\dimexpr\topskip-\height}{%
		
\includegraphics[scale=0.20]{../Logos_de_Python/Advertencia}}  %[width=\textwidth]
	
\end{minipage}%\hfill
\begin{minipage}[t]{0.70\textwidth} 
	{\scriptsize 
		Ordenar una lista alfabéticamente es un poco más complicado cuando no todos los valores están en
		minúscula. Hay varias formas de interpretar las letras mayúsculas al determinar un tipo de 	orden, y especificar el orden exacto puede ser más complejo de lo que queremos tratar 	en este momento. Sin embargo, la mayoría de los métodos de clasificación se basarán directamente en lo que usted aprendido en esta sección. 
		En el siguiente ejemplo, el método falla al intentar ordenar, pruébelo en un archivo .py.
		\begin{lstlisting}	
			autos =  ['FOrd','BMw', 'chevrolet', 'TOoyota']
			print(autos.sort())
			['BMw', 'FOrd', 'TOoyota', 'chevrolet']
		\end{lstlisting}	
	}
\end{minipage}
\vspace{0.2cm}



\subsubsection{Ordenar listas de forma temporal con la función {\ttfamily sorted()} }

Para mantener el orden original de una lista y presentarla ordenada, se puede hacer uso de la función  {\ttfamily sorted()}, que permite mostrar la lista en un orden específico, y no afecta el orden original de la lista.
Veamos con el ejemplo de la lista de autos:


\hspace{-2.25cm}{\ttfamily \scriptsize autos.py}  \vspace{-0.5cm}
\begin{lstlisting}	
	# Lista original:
	autos = ['ford','bmw', 'chevrolet', 'toyota']
	
	# Lista ordenada temporal:
	print(sorted(autos))
	['bmw', 'chevrolet', 'ford', 'toyota']
	
	# Lista original:
	print(autos)
	['ford','bmw', 'chevrolet', 'toyota']
	
\end{lstlisting}	

La función  {\ttfamily sorted()}  cambia el orden de la lista de forma temporal. Después de aplicar la función  a la lista,  se muestra en el nuevo orden, y al imprimir la lista se muestra en su orden original.

La función  {\ttfamily sorted()} también permite el orden alfabético inverso con el argumento  {\ttfamily reverse=True}, ejemplo: 
{\ttfamily sorted(autos,reverse=True)}.

\subsubsection{Imprimir una lista en orden inverso con reverse()}

El método {\ttfamily  reverse()}  no ordena alfabéticamente hacia atrás,  simplemente
invierte el orden de la lista:

\hspace{-2.4cm}{\ttfamily \scriptsize invierte.py}\vspace{-0.5cm}
\begin{lstlisting}	
	# Lista original:
	autos = ['ford','bmw', 'chevrolet', 'toyota']
	
	# Lista en reversa:
	autos.reverse()
	print(autos)
\end{lstlisting}	

{\ttfamily\scriptsize\scriptsize 
	\hspace{-0.65cm}Out[]:\vspace{-0.2cm}	
	{\color{blue} 	
		\begin{verbatim}
			['toyota', 'chevrolet', 'bmw', 'ford']
		\end{verbatim}
}}
El método  {\ttfamily  reverse()}  cambia el orden de una lista de manera permanente; pero puede  volver al orden original en cualquier momento aplicando nuevamente el método {\ttfamily  reverse()} sobre la lista.



\subsection{Errores con índice al trabajar con listas}

Hay un tipo de error muy común cuando trabajas con listas por primera vez. Digamos que tienes una lista con tres elementos y preguntas por el cuarto elemento:

\begin{lstlisting}
	
	# Lista original:
	autos = ['ford','bmw', 'chevrolet', 'toyota']
	print(autos[4])
\end{lstlisting}
Este ejemplo da como resultado el error de índice:

{\ttfamily\scriptsize\scriptsize 
	\hspace{-0.65cm}  \vspace{-0.3cm}	
	{\color{blue} 	
		\begin{verbatim}
			Traceback (most recent call last):
			File "C:\Users\Linis\Desktop\IDLE\autos.py", line 21, in <module>
			print(autos[4])
			IndexError: list index out of range
		\end{verbatim}
}}
\noindent Error de índice ({\ttfamily IndexError}): Indica que el índice de lista está fuera de rango. 
Significa que Python no  encontra un elemento en el índice solicitado. Pues el tamaño de la lista es 4, y no hay elemento en la posición 4, recuerde que la indexación en listas de Python inicia en cero.\\

Es importante tener en cuenta que para acceder al último elemento de una lista, es más seguro usar el índice -1. Esto siempre funcionará, incluso si su lista ha cambiado de tamaño desde la última vez que accediste.
La única vez que este enfoque causará  error es cuando solicite el último elemento de una lista vacía:
\begin{lstlisting}
	In[]
	
	>>> auto= []
	>>> print(auto[-1])
\end{lstlisting}

{\ttfamily\scriptsize\scriptsize 
	\hspace{-0.65cm}Out[]\vspace{-0.3cm}	
	{\color{blue} 	
		\begin{verbatim}
			Traceback (most recent call last):
			File "C:\Users\Linis\Desktop\IDLE\Tarea4.py", line 23, in <module>
			print(auto[-1])
			IndexError: list index out of range
		\end{verbatim}
}}
\noindent No hay elementos en la nueva lista auto, por lo que Python devuelve un error de índice.

Si se produce un error de índice y no sabe cómo resolverlo, imprima su lista o simplemente calcule el tamaño
de la lista,  pueden ayudarle a solucionar dichos errores lógicos.



\subsection{Resumen}

En esta sección, se estudió  listas,  cómo  definir lista, tamaño  de una lista, cómo trabajar con los elementos de una lista, agregar y eliminar elementos; también  se estudió cómo ordenar listas de forma permanente y temporal con fines de visualizar;  y se analizó posibles errores de índice cuando trabaja con listas y cómo evitarlos.  En la próxima sección se  trabajara con lista de manera más eficiente. 



\subsection{Ejercicios}\label{Tarea.4.2.8}

Pruebe con estos programas cortos para obtener experiencia de primera mano con las listas de Python.
Es recomendable crear una nueva carpeta para los ejercicios de cada capítulo o por tema, para mantener un orden.

\begin{enumerate}[4.2.8.1]
	
	\item Nombres: Almacene los nombres de algunos de sus amigos en una lista llamada nombres. Imprima el 
	nombre de cada persona accediendo a cada elemento de la lista, uno a uno.
	
	\item Saludos: Comience con la lista que utilizó en el ejercicio 3.3.1, pero en lugar de 
	imprimir solo el nombre de cada persona, imprima un mensaje para ellos. El texto de cada mensaje
	debe ser el mismo, pero cada mensaje debe ser personalizado con el nombre de persona.
	
	\item Su propia lista: Piense en su medio de transporte favorito, como una moto, un coche, y haga una 
	lista que almacene varios ejemplos. Usa su lista para imprimir una serie de declaraciones sobre estos artículos, como {}"{}Me gustaría tener una Motocicleta Honda{}"{}.
\end{enumerate}


Los siguientes ejercicios le brindan la oportunidad de utilizar listas en todas las formas descritas con un  poco más de complejidad:  


\begin{enumerate}[4.2.8.4]
	
	\item Lista de invitados: Si pudiera invitar a cenar a alguien,   ¿qué le invitarías? 
	Haga una lista que incluya al menos tres personas con las que le gustaría invitar a cenar. 
	Luego use su lista para imprimir un mensaje para cada invitando a cenar.
	
	
	\item[4.2.8.5] Cambiar la lista de invitados: Acaba de enterarse  que uno de sus invitados no puede asistir
	cena, por lo que deberás enviar una nueva lista de invitaciones. Tendrás que pensar en alguien más a quien invitar.
	
	Usando el programa del Ejercicio anterior, agregue una llamada {\ttfamily{print()}} al final de
	su programa, indicando el nombre del invitado que no puede asistir.
	
	Modifica tu lista, reemplazando el nombre del invitando que no puede asistir con el
	nombre del  nuevo  invitando.
	
	Imprima un segundo conjunto de mensajes de invitación, uno para cada persona que aún esté en
	tu lista.
	
	
	\item[4.2.8.6] Más invitados: Acabas de encontrar una mesa más grande para la cena, por lo que ahora hay más espacio
	disponible. Piensa en tres invitados más para la cena.
	
	Comience con su programa del Ejercicio 4.2.8.4. Agregue una llamada {\ttfamily{print()}} al
	final de su programa, informando a la gente que encontró una mesa más grande.
	
	Utilice {\ttfamily{insert()}}  para agregar un nuevo invitado al principio de su lista.
	
	Utilice  {\ttfamily{insert()}} para agregar un nuevo invitado al medio de su lista.
	
	Utilice {\ttfamily{append()}} para agregar un nuevo invitado al final de su lista.
	
	Imprima un nuevo conjunto de mensajes de invitación, uno para cada persona de su lista.
	
	\item[4.2.8.7] Lista de invitados cada vez más reducida: Acaba de descubrir que su nueva mesa no 
	llegaría  a tiempo para la cena, y ahora tendrá espacio para solo dos invitados.
	
	Comience con su programa del Ejercicio  4.2.8.4.6. Agregue una nueva línea que imprima un
	mensaje que dice que solo puedes invitar a dos personas a cenar.
	
	Utilice {\ttfamily{pop()}} para eliminar invitados de su lista uno por uno, hasta que solo queden dos.
	Los nombres que permanecen en su lista. Cada vez que saque un nombre de su lista, imprima un
	mensaje a esa persona haciéndole saber que lamentas no poder invitarlo para cenar.
	
	Imprima un mensaje para cada una de las dos personas que aún están en su lista, haciéndole saber de su invitación.
	
	Utilice del para  eliminar los dos últimos nombres de su lista, de modo que tenga un lista vacía.
	Imprima su lista para asegurarse que tiene una lista vacía al final de su programa.
	
\end{enumerate}


\begin{enumerate}[4.2.8.8]
	
	\item   Visitando el mundo: Piensa  al menos cinco lugares del mundo que te gustarían visitar.
	
	Almacenar las ubicaciones en una lista. Asegúrese de que la lista no esté en orden alfabético.
	
	Imprima su lista en su orden original. No se preocupe por imprimir la lista de forma ordenada;
	simplemente imprima como una lista de Python sin formato.
	
	
	Utilice {\ttfamily{sorted()}} para imprimir su lista en orden alfabético sin modificar el lista real.
	
	Demuestre que su lista todavía está en su orden original oprimiéndola nuevamente.
	
	Utilice {\ttfamily{reverse()}} para cambiar el orden de su lista. Imprima la lista para mostrar que el orden ha cambiado.
	
	
	Utilice  nuevamente {\ttfamily{reverse()}} para cambiar el orden de su lista. Imprime la lista para mostrar
	ha vuelto a su orden original.
	
	Utilice {\ttfamily{sort()}} para cambiar su lista y se almacene en orden alfabético. Imprime la lista para mostrar que su orden ha sido cambiado.
	
	Use nuevamente  {\ttfamily{sort()}} para cambiar su lista y se almacene en orden alfabético. Imprime la lista para mostrar
	ha vuelto a su orden original.
	
	
\end{enumerate}



\begin{enumerate}[4.2.8.9]
	
	\item Invitados a la cena: Trabajando con uno de los programas de los ejercicios  4.2.8.4
	a 4.2.8.7, use la función {\ttfamily{len()}} para imprimir un mensaje que indique el número
	de personas invitadas a la cenar.
	
	
	\item [4.2.8.10]Cada función: Piensa en cosas que podrías almacenar en una lista. Por ejemplo podrías hacer una lista de montañas, ríos, países, ciudades, lenguajes o cualquier cosa, más te gustaría. Escriba un programa que cree una lista que contenga estos elementos y
	luego utiliza cada función introducida en este capítulo al menos una vez.
	
	\item[4.2.8.11] Error intencional: Si no ha recibido un error index  en uno de sus programas, intenta que esto suceda.
	Cambiar un índice en uno de sus programas para producir un index error. Asegúrese de corregir el error antes de cerrar el
	programa.
	
\end{enumerate}




\subsection{Trabajando con listas}

En la sección anterior se trabajó con listas en forma simple, se aprendió a construir listas,  agregar y eliminar elementos, se trataron algunos métodos y funciones para lograr cambios permanentes y temporales  en  listas.  

En esta sección se incluirá algunas líneas de código, para tratar listas con instrucciones para recorrer una lista completa,  independientemente del  tamaño de la lista, y tiempo para ejecutar una acción o sentencia. 

Cuando se necesita repetir una misma acción con cada elemento de una lista, es conveniente hacerlo con bucles (Iteraciones o repeticiones).  Esta una de las forma más común en que una computadora automatiza las tareas repetitivas.
Un bucle permite realizar una misma acción o configurar acciones, con cada elemento de una lista. 
Al finalizar esta sección,  estarás en capacidad  de trabajar eficientemente con listas de cualquier tamaño, incluyendo aquellos con miles o incluso millones de elementos.

\subsubsection{Recorriendo una lista completa}

Por lo general es necesario revisar todas las entradas de una lista, realizando la misma
tarea con cada elemento. Por ejemplo, en un juego tal vez es necesario presentar cada elemento en la pantalla con una lista de números, es posible que se requiera la misma operación estadística en cada elemento; o tal vez se necesite mostrar cada título de una lista de artículos en un sitio web.  Cuando se realiza la misma acción con cada elemento de una lista, se debe hacer con ayuda de un bucle.\\

Consideremos un caso práctico, supongamos que tenemos una lista de nombres de los personas participantes en un curso online de Python, y necesitamos imprimir cada nombre. Este trabajo lo podríamos hacer recuperando cada nombre de la lista individualmente, sin embargo este enfoque puede causar algunos problemas; por un lado, sería muy repetitivo hacer esto con una larga lista de nombres. Además, cada vez que cambie la longitud de la lista, tendríamos que cambiar el código. El uso de un bucle {\ttfamily for} evita estos problemas, al permitir que Python gestione internamente esta situación. En Python los objetos como listas, cadenas, tuplas, diccionarios y conjuntos son iterables. 

Ilustremos la aplicación del bucle {\ttfamily for} con una lista de participantes en el curso:
Una vez tengamos la lista, tal como se hizo en la sección anterior, se define
un bucle {\ttfamily for}; en esta línea se le indica a Python que extraiga un nombre de la lista de {\ttfamily participantes}, y se asocia con la variable {\ttfamily participante}, y se le indica a Python que imprima el nombre del participante que  acaba de    

\hspace{-2.3cm}{\ttfamily \scriptsize curso.py} \vspace{-0.58cm}
\begin{lstlisting}		
	participantes = ["antonio", "Oswaldo", "manuel", "carlos", "maría"]
	for participante in participantes:
	print(participante)	
\end{lstlisting}	

\noindent ser asignado. Python luego repite estas últimas dos líneas, una para cada nombre de la lista. Podría ser útil leer este código como:  "Para cada participante en la lista de participantes, imprima el nombre del participante". El resultado es una impresión simple de cada nombre en la lista:

{\ttfamily\scriptsize\scriptsize 
	\noindent Out[]: \vspace{-0.3cm}	
	{\color{blue} 	
		\begin{verbatim}		
			antonio
			Oswaldo
			manuel
			carlos
			maría
		\end{verbatim}
}}
\noindent Observe, para cada elemento dentro de un bucle {\ttfamily for} se usó el nombre de participante en singular y plural para la lista,  estas convención de nomenclatura que pueden ayudarle a seguir la acción que se está realizando,
identificando si una sección de código está funcionando con un solo elemento de la lista o la lista completa.


\subsection{Iterando con el bucle {\ttfamily for ... in} }


Puedes hacer casi cualquier cosa con cada elemento de una lista en un bucle {\ttfamily for}. Incluso puede iterar sobre datos creados sobre la marcha, lo que permite el procesamiento de flujos de datos, que de otro modo no tal vez generaría problemas en la memoria del computadora .


Consideremos el ejemplo anterior,  podemos escribir en el bucle un mensaje a cada participante, felicitándolos por la decisión de tomar el curso:

\hspace{-2.3cm}{\ttfamily \scriptsize curso.py}  \vspace{-0.58cm}
\begin{lstlisting}		
	participantes = ["antonio", "Oswaldo", "manuel", "carlos", "maría"]
	for participante in participantes:
	print(f"Felicitaciones {participante.title()}, fue su mejor decisión.")	  
\end{lstlisting}	

{\ttfamily\scriptsize\scriptsize 
	\hspace{-0.65cm} \vspace{-0.3cm}	
	{\color{blue} 	
		\begin{verbatim}
			Felicitaciones Antonio, fue su mejor decisión.
			Felicitaciones Oswaldo, fue su mejor decisión.
			Felicitaciones Manuel, fue su mejor decisión.
			Felicitaciones Carlos, fue su mejor decisión.
			Felicitaciones María, fue su mejor decisión.
		\end{verbatim}
}}

Puede escribir tantas líneas de código como quieras en el bucle  {\ttfamily for}. Cada  
línea  debajo del  {\ttfamily for} con la indentación\footnote{En Python, la indentación son los espacios que se dejan al lado izquierdo en las líneas del código para identificar a qué bloque de código pertenece cada línea. Por convención, se usan cuatro espacios como indentación  (PEP 8).} o sangrado de 4 espacios en blanco, con lo que se le indica a Python que esta dentro el bucle (bloque),  cada línea indentada se ejecuta una vez para cada valor de la lista. Por lo tanto, puedes hacer todo el trabajo que quieras con cada valor de la lista.
Si en las  líneas posteriores no introducimos la indentación, significa bloque de código ha finalizado, en este caso es el final  del {\ttfamily for}. Para finalizar un bloque de código simplemente se dejar de introducir la indentación, no se requiere usar alguna llave u otro símbolo.

\begin{lstlisting}		
	participantes = ["antonio", "Oswaldo", "manuel", "carlos", "maría"]
	for participante in participantes:
	print(f"Felicitaciones {participante.title()},"
	" fue su mejor decisión.")
	print("Ya vienen mas sorpresas para seguir avanzando,"
	f"Preparate {participante.title()}.\n")	
\end{lstlisting}

Observe, la sentencia {\ttfamily print()} en los dos llamados cumplen con la indentación;  de esta forma,  cada línea se ejecutará una vez por cada participante en la lista. También que se agregó una nueva sentencia ("{}{ \scriptsize {\bf {\textbackslash}}}n{}") al final de la segunda línea, esta sentencia inserta una línea en blanco después de cada paso del bucle. Creando un conjunto de mensajes que están cuidadosamente agrupados para cada persona en la lista:

{\ttfamily\scriptsize\scriptsize 
	\hspace{-0.65cm} \vspace{-0.3cm}	
	{\color{blue} 	
		\begin{verbatim}
			Felicitaciones Antonio, fue su mejor decisión.
			Ya vienen mas sorpresas para seguir avanzando, Preparate Antonio.
			
			Felicitaciones Oswaldo, fue su mejor decisión.
			Ya vienen mas sorpresas para seguir avanzando, Preparate Oswaldo.
			
			Felicitaciones Mauel, fue su mejor decisión.
			Ya vienen mas sorpresas para seguir avanzando, Preparate Manuel.
			
			Felicitaciones Carlos, fue su mejor decisión.
			Ya vienen mas sorpresas para seguir avanzando, Preparate Carlos.
			
			Felicitaciones María, fue su mejor decisión.
			Ya vienen mas sorpresas para seguir avanzando, Preparate María.
		\end{verbatim}
}}

Después del bloque del bucle {\ttfamily for}, podemos escribir una despedida general para todos, sin indentación: 

\begin{lstlisting}
	for participante in participantes:
	print(f"Felicitaciones {participante.title()},"
	" fue su mejor decisión.")
	print("Ya vienen mas sorpresas para seguir avanzando,"
	f"Preparate {participante.title()}.\n")
	print("Hasta la próxima clase!!!")         	
\end{lstlisting}

\noindent Pruébelo y analice el resultado.


\subsection{Evitar errores de Indentación}


Python utiliza la indentación  para determinar cómo una línea, o grupo de líneas, se relaciona con el resto del programa. En los ejemplos anteriores, las líneas que imprimían mensajes a los participantes eran parte del bucle {\ttfamily for}, porque estaban dentro de la indentación de buble.
\begin{figure}[h!]
	\centering  
	%
\includegraphics[scale=0.25]{../Logos_de_Python/Advertencia}
	\includegraphics[width=6cm, height=3.6cm]{../Logos_de_Python/Bloques} %scale=0.40,
	\caption{Diferentes bloques}
	\label{Indentacion}
\end{figure} El uso de indentación de Python hace que el código sea muy fácil de leer. Básicamente, utiliza espacios en blanco para forzarte a escribir código con un formato con una estructura visual clara. En programas más largos, notarás bloques de código con indentación en diferentes niveles (ver figura \ref{Indentacion}). Estos niveles de indentación ayudan a visualizar un sentido de la organización general del programa. 

A medida que comienzas a escribir código basado en la indentación adecuada, tendrás que tener en cuenta algunos errores de identación muy comunes: Por ejemplo,  algunas veces indentará líneas de código que no necesitan ser identadas o tal vez olvidará identar líneas que necesitan ser identadas. 

Ilustremos con algunos ejemplos de estos errores que le  ayudará a evitarlos en el futuro y corregirlos cuando aparezcan en tus propios programas. Después de la declaración {\ttfamily for} en un bucle. Si lo olvidas identar, Python  te lo recordará:

\begin{lstlisting}
	participantes = ["antonio", "Oswaldo", "manuel", "carlos", "maría"]
	for participante in participantes:
	print(f"Felicitaciones {participante.upper().title()}, fue su mejor decisión.")	
\end{lstlisting} 

\noindent Al ejecutar el código,  Python genera el mensaje:\\

{\ttfamily\scriptsize\scriptsize 
	\hspace{-0.65cm} SyntaxError \vspace{-0.3cm}	
	{\color{blue} 	
		\begin{verbatim}
			IndentationError: expected an indented block after 'for' statement on line 3
		\end{verbatim}
}}



Generalmente puedes resolver este tipo de error  aplicando la identación inmediatamente después de la declaración {\ttfamily for}.

Otras veces puede hacer una identacion innecesaria, o tal vez olvidas identar en un bucle, en este caso el código puede  correr sin errores pero no producirá el resultado esperado. Por ejemplo, si olvidamos identar la segunda una línea en el bucle, que le dice a cada participante que vendrán nuevas sorpresas. Este es un error lógico, y la sintaxis es código válido en Python, pero el código no producirá el resultado deseado porque hay un problema en su lógica. Si espere una determinada acción repetida en cada uno de los elementos de una lista y es ejecutado solo una vez, debe  asegúrese  identar bien las líneas o el  grupo de líneas.\\


Los dos puntos al final de una declaración {\ttfamily for} le indica a Python que interprete la siguiente
línea como inicio del bucle. Si accidentalmente olvida los dos puntos al final de la declaración {\ttfamily for} 
obtendrá un error de sintaxis, porque Python no sabe exactamente lo que estás intentando hacer. 

Python no sabe si olvidaste los dos puntos o si  quiere escribir 
un código adicional para configurar un bucle más complejo. Si el intérprete 
identifica una posible solución   le dará la  sugerencia,  que puede ser agregar dos puntos al final
de una línea. Algunos errores tienen soluciones fáciles y obvias, gracias a las sugerencias 
en los rastreos de Python, y otros errores son un poco más difíciles de resolver, 
incluso cuando la solución final sólo implica un solo carácter. 


\subsection{Hacer listas numéricas}

Existen muchas razones para almacenar un conjunto de números. Por ejemplo, necesitarás
un registro de las posiciones de cada personaje en un juego, y  es posible que también  necesite 
llevar un seguimiento de las puntuaciones más altas de un jugador. En general al tratar datos numéricos, se presentan conjuntos de números, como temperaturas, distancias, tamaños de población, o valores de latitud y longitud, entre otros tipos de conjuntos numéricos.\\

Las listas son ideales para almacenar conjuntos de números y Python proporciona una variedad
de herramientas que le ayudarán a trabajar eficientemente con listas de números. Una vez que entiendas cómo utilizar estas herramientas de manera efectiva, su código funcionará bien incluso cuando
sus listas contengan  millones de elementos.


Pero antes de continuar con listas haremos una introducción a la función {\ttfamily range()},  una función integrada de Python que generar secuencias numéricas iterativas.


\section{La función range()}\label{rango}   % Al igual que zip(), 
%  {\ttfamily range()}

La función  {\ttfamily range()} devuelve una salida de números dentro de un rango específico,  sin tener que crear un estructura de datos grande para almacenar, como una lista o tupla. Esto  permite crear rangos enormes sin usar toda la 
memoria de su computadora, y sin generar problemas su programa. Esta función devuelve un objeto iterable, y para recorrer los valores es  necesario usar un  {\ttfamily for ... in} , o convertir el objeto en una secuencia como una lista, tupla o cualquier otro objeto.\\  


En Python la función {\ttfamily range()}  necesita argumentos para generar la secuencia;  su uso es semejante a  tuplas y listas:
{\ttfamily range(inicio, final, paso)}. Si se omite el paso, la función usa un paso por defecto, y si se omite el inicio comienza en 0 (valor por defecto, también puede iniciar en valores negativos). El único valor requerido es final ({\ttfamily range(3)}, inicia en cero y finaliza en la posición dos, {\ttfamily range(0,3)}). Siguiendo la metodología para porciones en lista y tuplas, el último valor de la secuencia será justo antes del final. Y por supuesto también funciona con paso negativos para ir hacia atrás, por ejemplo con -1 será el orden inverso. Ilustremos con un ejemplo el uso de la función {\ttfamily range()}:

\hspace{-2.5cm}{\ttfamily \scriptsize rango01.py}
{\vspace{-0.5cm}
	\begin{lstlisting}
		for valor in range(1, 4):
		print(valor):
\end{lstlisting}}

{\ttfamily\scriptsize\scriptsize\noindent Out[]:\vspace{-0.3cm}{\color{blue}
		\begin{verbatim}
			1
			2
			3
		\end{verbatim}
	}
}\vspace{-0.1cm}
\noindent En esta ilustración, la función  imprime los números del 1 al 3;  Python comienza a contar desde el primer valor inicial dado,  hasta un valor antes del valor final dado.  El resultado inicia en uno y llega hasta tres; pues el número final dado es 4, y el paso es uno por defecto.


\subsection{Usando range() para hacer listas de números}

Los resultados de {\ttfamily range()} se pueden convertir directamente en una lista mediante la función  {\ttfamily list()}.  En el ejemplo anterior (sección \ref{rango}),  imprimimos una serie de
números. Ahora  usaremos la función  {\ttfamily list()} para convertir ese mismo conjunto de números en una lista:

\hspace{-2.4cm}{\ttfamily \scriptsize rango02.py}  \vspace{-0.53cm}
\begin{lstlisting}
	números = list(rango(1, 4))
	print(números) # El resultado una lista:
\end{lstlisting}

{\ttfamily\scriptsize\scriptsize\noindent Out[]: \vspace{-0.3cm}	%      	\hspace{-0.65cm}	
	{\color{blue} 	
		\begin{verbatim}
			[1, 2, 3]
		\end{verbatim}
}}


También podemos usar la función  {\ttfamily range()} para indicarle a Python que omita números en un
rango dado, con el argumento paso. Por ejemplo, enumerar los números pares entre -2 y 10:

\hspace{-2.2cm}{\ttfamily \scriptsize pares.py}  \vspace{-0.52cm}
\begin{lstlisting}
	números\_pares = list(range(-2, 11, 2))
	print(números\_pares) 
\end{lstlisting}	
En este ejemplo, la función {\ttfamily range()} comienza con el valor -2 y luego
suma 2 a ese valor. Agrega 2 repetidamente hasta llegar uno antes del final, y produce el siguiente resultado:

\vspace{0.09cm}
{\ttfamily\scriptsize\scriptsize\noindent Out[]: \vspace{-0.3cm}	
	{\color{blue} 	
		\begin{verbatim}
			[-2, 0, 2, 4, 6, 8, 10]
		\end{verbatim}
}} 


Puedes crear casi cualquier conjunto de números que quieras usando la función  {\ttfamily range()}. Por ejemplo, considere cómo podría hacer una lista los primeros cinco números al cuadrados (es decir, el cuadrado de cada número entero del 1 al 5). Recuerde que en Python dos asteriscos (**) representan exponentes.

Como ejemplo pongamos los primeros 5 números cuadrados en una lista, comenzamos con una lista vacía que llamaremos cuadrados. 
Luego un bucle con el {\ttfamily for ...in} con  la función {\ttfamily range()}  del 1 al 5; dentro del bucle, el valor actual se eleva al exponente 2 y se asigna a la variable cuadrado. Cada nuevo valor de cuadrado se agrega a la lista de cuadrados.
Finalmente, cuando el bucle ha terminado de ejecutarse, se imprime la lista de cuadrados:

\hspace{-3cm}{\ttfamily \scriptsize cuadrados01.py}  \vspace{-0.5cm}
\begin{lstlisting}
	cuadrados = []
	for valor in range(1, 6):
	cuadrado = valor ** 2 # variable temporal
	cuadrados.append(cuadrado)
	print(cuadrados)
\end{lstlisting}

{\ttfamily\scriptsize\scriptsize 
	\noindent Out[]:\vspace{-0.3cm}	
	{\color{blue} 	
		\begin{verbatim}
			[1, 4, 9, 16, 25]
		\end{verbatim}
}}
\noindent Para escribir este código de manera más concisa, omita la variable temporal cuadrado
y agregue cada nuevo valor directamente a la lista:

\hspace{-3cm}{\ttfamily \scriptsize cuadrados02.py}  \vspace{-0.5cm}
\begin{lstlisting}
	cuadrados = []
	for valor in range(1, 6):
	cuadrados.append(valor ** 2)
	print(cuadrados)
\end{lstlisting}

{\ttfamily\scriptsize\scriptsize 
	\noindent Out[]:\vspace{-0.3cm}	
	{\color{blue} 	
		\begin{verbatim}
			[1, 4, 9, 16, 25]
		\end{verbatim}
}}

Estas líneas de código  hace el mismo trabajo que las líneas en el anterior ejemplo anterior.
Cada valor en el bucle se eleva a la segunda potencia y luego inmediatamente añadido a la lista de cuadrados.
Puedes utilizar cualquiera de estos enfoques cuando hagas tareas más complejas. A veces, el uso de una variable temporal hace que su código sea más fácil de leer; otras veces hace que el código sea innecesariamente largo. Concéntrate en escribir el código que entienda claramente, y luego busque enfoques más eficientes mientras revisa su código.

\subsection{Estadísticas simples con una lista de números}

Algunas funciones de Python son útiles cuando se trabaja con listas de números. 
Por ejemplo, puede encontrar fácilmente el mínimo, el máximo y la suma de los elementos  de una lista de
números:
\begin{lstlisting}
	>>> dígitos = [1, 2, 3, 4, 5, 6, 7, 8, 9, 0]
	>>> min(dígitos)
	0
	>>> máx(dígitos)
	9
	>>> suma(dígitos)
	45
\end{lstlisting}


\vspace{0.2cm}
\begin{minipage}[t]{0.2\textwidth}
	\centering\raisebox{\dimexpr\topskip-\height}{%
		
\includegraphics[scale=0.20]{../Logos_de_Python/Busquedas_R}}  %[width=\textwidth]
	
\end{minipage}%\hfill
\begin{minipage}[t]{0.70\textwidth} 
	{\scriptsize 
		Los ejemplos de esta sección utilizan listas cortas de números, pero igualmente funcionan  bien  para listas de un millón o más.
	}
\end{minipage}
\vspace{0.2cm}

%%%%%%%%%%%%%%%%%%%%%%%%%%%%%%%%%%%%%%%%%%%%%%%%%%%%%%%



\subsection{Ejercicios}\label{Tarea3.2}

Ejercicios con listas....

\begin{enumerate}[4.2.16.1]
	
	\item  Contando hasta veinte: use un bucle {\ttfamily{for}}  para imprimir los números del 1 al 20,
	inclusivo.
	
	\item [4.2.16.2]Un millón: Haga una lista de los números del uno al un millón y luego
	use un bucle {\ttfamily{for}}  para imprimir los números. (Si la salida tarda demasiado, deténgala
	presionando CTRL-C o cerrando la ventana de salida).
	
	\item [4.2.16.3]Números impares: Use el tercer argumento de la función {{ttfamily{range()}} para hacer una lista
		de los números impares del 1 al 20. Utilice un bucle {\ttfamily{for}}  para imprimir cada número.
		
		\item De tres en tres: Haga una lista de los múltiplos de 3, del 3 al 30. Usa un bucle {\ttfamily{for}} para
		imprime los números en la lista.
		
		\item  Cubos: Un número elevado a la tercera potencia se llama cubo. Por ejemplo,
		el cubo de 2 se escribe como 2**3 en Python. Haga una lista de los primeros 10 cubos (es decir, el cubo de cada número entero del 1 al 10) y utilice un bucle {\ttfamily{for}} para imprimir el valor de cada cubo.		
		
		\item  Animales: Piensa en al menos tres animales diferentes que tengan una característica común. Almacene los nombres de estos animales en una lista y luego use un bucle {\ttfamily{for}} para imprime el nombre de cada animal.
		
		Modifique su programa para imprimir una declaración sobre cada animal, como el perro sería una gran mascota.
		
		Añada una línea al final de tu programa, indicando qué tienen estos animales en común. Podrías imprimir una oración, como cualquiera de estos animales puede  ¡Ser una gran mascota!
		
		\item[4.2.16.7] Pizzas: Piensa en al menos tres tipos de sus pizzas favoritas. Guarde estos
		nombres de pizza en una lista y luego use un bucle {\ttfamily{for}} para imprimir el nombre de cada pizza.
		
		Modifica tu bucle {\ttfamily{for}} para imprimir una oración usando el nombre de la pizza, 	en lugar de imprimir solo el nombre de la pizza. Para cada pizza, construya una línea de salida que contenga una declaración simple, tal vez {}"{}Me gustan las pizza  con  pepperoni{}"{}}.
	
	Agregue una línea al final de su programa, fuera del bucle {\ttfamily{for}}, que indique cuánto te gusta la pizza.
	La salida debe tener tres o más líneas, sobre los tipos de pizza que te gustan y luego una oración adicional, como ¡Realmente amo la pizza!
	
	\item Sumar un millón: Haga una lista de los números del uno al un millón y
	luego use {\ttfamily{min()}} y {\ttfamily{max()}} para asegurarse de que su lista realmente comience en uno y termine
	en un millón. 
	Además, use la función {\ttfamily{sum()}} para ver qué tan rápido Python puede sumar un millón de números.
	
	
	\item  Mis pizzas: Comienza con su programa del ejercicio 4.2.16.7. Haga una copia de la lista de pizzas y llámela {\ttfamily{amigo\_pizzas}}. Entonces, haga lo siguiente:
	
	Agregue una nueva pizza a la lista original.
	
	Agregue una pizza diferente a la lista {\ttfamily {amigo\_pizzas}}.
	
	Demuestre que tiene dos listas separadas. Imprimir el mensaje Mis pizzas favoritas 	son:, y luego use un bucle {\ttfamily{for}} para imprimir la primera lista. Imprimir el mensaje: Las pizzas favoritas de mis amigo son: y luego usa un bucle {\ttfamily{for}} para imprimir la segunda lista. 	Asegúrese de que cada pizza nueva esté almacenada en la lista adecuada.
	
\end{enumerate}




\subsection{Trabajar con porciones de lista}

En la secciones anteriores se ha trabajado con elementos individuales en una lista o con todos los elementos de una lista.
También se trabajó porciones o partes específica de elementos en una lista; por ejemplo, si necesitas recorrer un subconjunto de elementos en una lista, se puede usar una parte (rebanada) de la lista en con el bucle {\ttfamily for}. En el siguiente ejemplo, recorremos los primeros tres participantes, e imprimimos sus nombres como parte de una lista simple:

\hspace{-2.3cm}{\ttfamily \scriptsize porcion.py}\vspace{-0.5cm}
\begin{lstlisting}
	participantes = ["Antonio", "oswaldo", "Manuel", "Carlos", "María"]
	print("Aquí están los tres primeros participantes en el curso de Python:\n")
	for participante in participantes[:3]:
	print(participante.title())
\end{lstlisting}
{\ttfamily\scriptsize\scriptsize\noindent Out[]:\vspace{-0.3cm}	  %  \hspace{0cm} 
	{\color{blue} 	
		\begin{verbatim}
			Aquí están los tres primeros participantes en el curso de Python:
			
			Antonio
			Oswaldo
			Manuel
		\end{verbatim}
}}

\noindent Sin embargo, esta tarea se puede desarrollar con una comprensión,  haciendo el código mas conciso. Veamos este fragmento de código en una comprensión:

\hspace{-2.9cm}{\ttfamily \scriptsize compresion.py}\vspace{-0.5cm}
\begin{lstlisting}
	print("Aquí están los tres primeros participantes en el curso de Python:")
	print([participante.title() for participante in participantes[:3]] )  
\end{lstlisting} 
{\ttfamily\scriptsize\scriptsize\noindent Out[]:\vspace{-0.25cm}	  %  \hspace{0cm} 
	{\color{blue} 
		\begin{verbatim}
			Aquí están los tres primeros participantes en el curso de Python:			
			['Antonio', 'Oswaldo', 'Manuel']
		\end{verbatim}
}}Observe, se usó una rebanada de la lista {\ttfamily participantes[:3]}, inicio por defecto y final tres.\\

Las porciones de listas son muy útiles en varias situaciones. Por ejemplo, cuando
estás creando un juego, puedes agregar la puntuación final de un jugador a una lista en cada
momento en que el jugador finaliza el juego; cuando trabajas con datos, puedes usar sectores para procesar sus datos en fragmentos de un tamaño específico; cuando estás construyendo una aplicación web, puede utilizar sectores para mostrar información en una serie de páginas con una cantidad adecuada de información en cada página.



%   Revisar los títulos de listas de Bill desde pag 42, para finalizar con listas. 


\subsection{Resumen de listas}

En esta sección primera parte, aprendiste cómo trabajar eficientemente con los elementos en un
lista. Aprendiste cómo trabajar con una lista usando un bucle   {\ttfamily for}, cómo Python
utiliza indentación  para estructurar un programa y cómo evitar algunos problemas comunes
errores de indentación. Aprendiste a hacer listas numéricas simples, así como algunas  operaciones se pueden realizar en listas numéricas. Aprendiste a cortar una lista para trabajar con un subconjunto de elementos rebanadas o porciones, y cómo copiar listas correctamente usando porciones. Y en la sección \ref{tuplas01} aprendió sobre las tuplas, que brindan cierto grado de protección a un conjunto de valores que no deberían cambiar, y cómo diseñar su cada vez más código complejo para que sea fácil de leer.\\


En la próxima sección  iniciamos el aprendizaje de una estructura de datos muy importante, los diccionarios de Python, utilizados básicamente para almacenar grandes cantidad de datos de forma ordenada con un par de valores, clave y valor. 


%%%%%%%%%%%%%%%%%%%%                %%%%%%%%%%%%%%%%%%%%%%%%%%
%%                      Diccionarios                           %
%%%%%%%%%%%%%%%%%%%%%%%                %%%%%%%%%%%%%%%%%%%%%%%%%

\section{Introducción a los Diccionarios}\label{diccionario}\hypertarget{diccionario}{}
% Los diccionarios se ordenan (según Python 3,7) y son mutables . pero el orden de los elementos no importa  


\begin{wrapfigure}{l}{0.5\textwidth} 
	\vspace{-25pt}	\begin{center} 
\includegraphics[ width=0.48\textwidth]{../Logos_de_Python/Diccionario1} % angle=45,
	\end{center} \vspace{-35pt}%\caption{Python 3} 
\end{wrapfigure} 
Los diccionarios en Python son estructura de datos para trabajar con colecciones almacenados en pares  claves-valores (key-value). Es decir, los elemento del diccionario son pares clave-valor,  denominado item. 
Se trata de una estructura de datos similar a las listas, las tuplas o conjuntos, con la particularidad de que el orden de las claves es arbitrario.  Son utilizados básicamente para almacenar grandes cantidad de datos de forma ordenada, con el objetivo de utilizarlos para algún propósito en concreto.  Son objetos mutables, por lo que puede soportar cambios después de su creación, como agregar, eliminar y cambiar sus elementos clave-valor; pero son estrictos respecto a entradas duplicadas, cada elemento es único, la clave es única.  Esta restricción aporta la ventaja que puedes acceder a todos los valores almacenados usando simplemente las claves.\\

La clave generalmente es una cadena, pero puede ser cualquiera tipo inmutables de Python, como booleano, entero, flotante, tupla; y como valor un diccionario, puedes usar cualquier objeto que puedas crear en Python.

Cuando de proporciona una clave, Python devuelve el valor asociado con esa clave. Cada clave es conectado a su valor mediante dos puntos; los elementos (pares clave-valor)  están separados por comas. Puede almacenar tantos elementos como desee en un diccionario,  un número ilimitado de pares clave-valor en un diccionario.\\



En otros lenguajes, los diccionarios podrían denominarse matrices asociativas,  hashes o mapas hash; si ha trabajado con lenguajes que sólo admiten matrices o listas, le encantarán los diccionarios.


\subsection{Crear un diccionario con \{\},  o con la función dict()}  % Convertir usando dict() 


En Python, un diccionario consiste en una serie de elementos en pares clave-valor entre llaves \{\}, separados por comas {\ttfamily \{key1: valor1, key2: valor2, ...\}}.  El diccionario más simple es el vacío, que no contiene claves ni valores:
\begin{lstlisting}
	In[]:		
	>>> dictado_vacío = {}
	>>> dictado_vacío
	{}	
\end{lstlisting}	
%  comentarios del uso del dic vacio

En el siguiente ejemplo definimos un diccionario con tres claves: nombre, apellido y la edad, y sus respectivos valores: Antonio, López, 25:

\begin{lstlisting}	
	In[]:		
	>>> diccionario = {"nombre": "Antonio", "apellido": "López", "eda": 45}
	>>> diccionario
	{'nombre': 'Antonio', 'apellido': 'López', 'eda': 45}
\end{lstlisting}


También se puede crear diccionarios usando la función integrada de Python {\ttfamily dict()},  a partir de  secuencias de dos valores. El primer elemento de cada secuencia se utiliza como clave y el
segundo como valor. Veamos algunos ejemplos: \\

Una lista de tuplas con dos elementos:

\begin{lstlisting}
	In[]:
	>>> lista_tuplas = [ ('a', 'b'), ('c', 'd'), ('e', 'f') ]
	>>> dict(lista_tuplas)
	{'a': 'b', 'c': 'd', 'e': 'f'}
\end{lstlisting}



Una tupla de listas con dos elementos:
\begin{lstlisting}
	In[]:
	>>> tupla_lista = ( ['a', 'b'], ['c', 'd'], ['e', 'f'] )
	>>> dict(tupla_lista)
	{'c': 'd', 'a': 'b', 'e': 'f'}
\end{lstlisting}

Una lista de cadenas con dos caracteres:
\begin{lstlisting}
	In[]:
	>>> lista_cadenas = [ 'ab', 'cd', 'ef' ]
	>>> dict(lista_cadenas)
	{'c': 'd', 'a': 'b', 'e': 'f'}
\end{lstlisting}

Una tupla de cadenas con dos caracteres:
\begin{lstlisting}
	In[]:
	>>> tupla_cadenas = ('ab', 'cd', 'ef' )
	>>> dictar(tupla_cadenas)
	{'c': 'd', 'a': 'b', 'e': 'f'}
\end{lstlisting}

Y una lista de listas con dos elementos:
\begin{lstlisting}
	In[]:
	>>> Lista_listas= [ ['a', 'b'], ['c', 'd'], ['e', 'f'] ]
	>>> dict(Lista_listas)
	{'c': 'd', 'a': 'b', 'e': 'f'}
\end{lstlisting}
La función {\ttfamily dict()} no afecta el objeto original, crea un nuevo objeto tipo diccionario.

Los objetos a convertir deben tener siempre elementos en pares. Recuerde que el orden de las claves en un diccionario es arbitrario y puede diferir dependiendo de cómo agregue elementos.



\subsection{Acceder a valores en un Diccionario}

Un diccionario en Python es una colección de pares clave-valor. Cada clave  está conectada a
un valor, se puede usar la clave para acceder al valor asociado con esa clave.   


Para obtener el valor asociado con una clave, se escribe el nombre del diccionario seguido de la clave dentro de un corchetes; al  proporcionar la clave, Python devuelve el valor asociado con esa clave. Para ilustrar consideremos un diccionario de prueba:

\begin{lstlisting}
	In[]:	
	# diccionario de prueba	
	>>> diccionario = {"nombre": "Antonio", "apellido": "López", "edad": 45}
	>>> diccionario["nombre"]
	'Antonio'
\end{lstlisting}

\noindent El valor asociado con la clave {}"{}nombre"{} es {}"{}Antonio"{}.\\


También puede usa como variable para incluir  el valor de una clave   en una {\ttfamily "f-cadena"{}}, por ejemplo:
\begin{lstlisting}
	In[]:	
	>>> variable = diccionario["nombre"]
	>>> print(f"Felicidades {variable} ¡Acabas de iniciar un nuevo curso!")
	Felicidades Antonio ¡Acabas de iniciar un nuevo curso!
\end{lstlisting}

De esta forma podemos extraemos el valor asociado con cualquier clave del diccionario, {\ttfamily 
	diccionario[{}"{}apellido{}"{}], diccionario[{}"{}edad{}"{}]}. El único inconveniente es que debe conocer la clave.


\subsection{Agregar elementos al diccionario}


Los diccionarios son estructuras dinámicas y se le pueden agregar nuevos elementos o item (par clave-valor)  en cualquier momento del programa.

Para agregar un nuevo par clave-valor, se escribe el nombre del diccionario seguido de la nueva clave entre corchetes, y se le asigna el nuevo valor: {\ttfamily diccionario[key] = {\bf{valor}}};  significa agrega la clave y su valor asignado.\\     

Agreguemos un nuevo elemento al diccionario de prueba: 
\begin{lstlisting}
	In[]:	
	>>> diccionario["dirección"] = "Mérida"
	>>> diccionario
	{'nombre': 'Antonio', 'apellido': 'López', 'eda': 45, 'dirección': 'Mérida'}
\end{lstlisting}

\noindent Aquí se observa  el nuevo elemento agregado al final del diccionario de prueba.


\subsection{Cambiar elementos en un diccionario}


Esta operación es semejante al procedimiento seguido con lista; para cambiar el valor de una clave en   diccionario, se coloca la clave en entre corchetes y se le asigna el nuevo valor, visualicemos con un ejemplo:

\begin{lstlisting}
	In[]:	
	>>> diccionario ={"nombre": "Antonio", "apellido": "López", "edad": 45, 
		"dirección": "Mérida"}
	# Cambio de dirección
	>>> diccionario["dirección"] = "Merida, Venezuela"
	>>> diccionario
	{'nombre': 'Antonio', 'apellido': 'López', 'edad': 45, 
		'dirección': 'Merida, Venezuela'}
\end{lstlisting}

En este ejemplo se cambió el valor de la clave {\ttfamily dirección}.\\


Ahora construimos un diccionario con la mayoría de los miembros de Monty Python{\footnote{Monty Python (también conocidos colectivamente como Pythons ) fue un grupo de comedia británico formado en 1969 formado por Graham Chapman , John Cleese , Terry Gilliam , Eric Idle , Terry Jones y Michael Palin . El grupo saltó a la fama por la serie de comedia Monty Python's Flying Circus , que se emitió en la BBC de 1969 a 1974}}, usando sus apellidos
como claves y nombres como valores:

\begin{lstlisting}
	>>> grupo = {'Chapman': 'Graham', 
		'Cleese': 'Juan', 
		'Inactivo': 'Eric', 
		'Jones': 'Terry', 
		'Palin': 'Michael'}
	>>> grupo
	{'Chapman': 'Graham', 'Cleese': 'Juan', 'Inactivo': 'Eric', 'Jones': 'Terry',
		'Palin': 'Michael'}
\end{lstlisting}

En este diccionario falta un miembro del grupo, Terry Gilliam nacido en Estados Unidos. 
Se procede agregar el nombre:

\begin{lstlisting}
	In[]:
	# nombre equivocado, error de transcripción! 
	>>> grupo["Gilliam"] = "Gerry"
	>>> grupo
	{'Chapman': 'Graham', 'Cleese': 'Juan', 'Inactivo': 'Eric', 'Jones': 'Terry', 
		'Palin': 'Michael', 'Gilliam': 'Gerry'}
\end{lstlisting}

El programador se equivoca con el nombre, por un error de transcripción; y procese a corregir el error:

\begin{lstlisting} 
	>>> grupo[el programador se equivoca con el nombre] = "Terry"
	>>> grupo
	{'Chapman': 'Graham', 'Cleese': 'Juan', 'Idle': 'Eric', 'Jones': 'Terry',
		'Palin': 'Michael', 'Gilliam': 'Terry'}
\end{lstlisting}

Usando la misma clave ({}"{}Gilliam{}"{}), reemplazamos el valor inicial {}"{}Gerry{}"{} con {}"{}Terry{}"{}.

Recuerde que las claves del diccionario deben ser únicas. Por eso usamos apellidos para las clave,
en lugar de nombres aquí: dos miembros de Monty Python tienen el mismo nombre ¡Terry!. Si usa una clave más de una vez, Python se queda con el último valor; ejemplo: 

\begin{lstlisting}
	In[]:
	>>> grupo["Gilliam"]
	'Gerry'
\end{lstlisting}
Python se usa el último elemento con esa clave, es decir actualiza.

\subsection{Unir diccionarios con update()}

La función {\ttfamily update()} se usa para copiar elementos (clave-valore) de un diccionario a 
otro. Veamos con un ejemplo, consideremos nuevamente de miembros de Monty Python, esta vez divididos en dos  diccionarios,   grupo01 y grupo02:


\begin{lstlisting}
	In[]:
	>>> grupo02 = {'Jones': 'Terry', 'Palin': 'Michael', 'Gilliam': 'Terry'}
	>>> grupo01 = {'Chapman': 'Graham', 'Cleese': 'Juan', 'Idle': 'Eric',}
\end{lstlisting}

Unimos estos dos diccionarios con  {\ttfamily update()}: 		
\begin{lstlisting}
	In[]:			
	>>> grupo01.update(grupo02)
	>>> grupo01
	{'Chapman': 'Graham', 'Cleese': 'Juan', 'Idle': 'Eric', 'Jones': 'Terry',
		'Palin': 'Michael', 'Gilliam': 'Terry'
	\end{lstlisting}
	Al diccionario  grupo01 se añadió los elementos del diccionario grupo02.
	
	\subsection{Diccionarios con claves iguales, y update()}
	
	
	Si el diccionario que se usa como argumento del método {\ttfamily update()} tiene algunas claves iguales al diccionario  principal, el valor del diccionario argumento es el que toma Python, veamos en el siguiente ejemplo:
	\begin{lstlisting}
		In[]:		
		>>> principal = {"A": 1, "B": "primero"}
		>>> argumento = {"C": 3, "B": "segundo"}
		>>> principal.update(argumento)
		>>> principal
		{'A': 1, 'B': 'segundo', 'C': 3}
	\end{lstlisting}
	El diccionario principal se actualiza con el  valor de la clave del segundo. 
	
	\subsection{Eliminar un elemento del diccionario con {\ttfamily del}}
	
	Con la sentencia {\ttfamily del}  se puede eliminar un elemento del diccionario usando la clave. 
	Consideremos el diccionario {\ttfamily datos}, para ejemplificar:
	
	\begin{lstlisting}
		In[]:	
		datos =  {"nombre": "Antonio", "apellido": "López", "edad": 45, 
			"dirección":"Mérida, Venezuela"}
		>>> del datos["edad"]
		>>> datos
		{'nombre': 'Antonio', 'apellido': 'López', 'dirección': 'Mérida, Venezuela'}
	\end{lstlisting}
	En el resultado ya no aparece el elemento {\ttfamily "{}edad{}"{}: 45}.
	
	
	\subsection{Eliminar todos los elementos del diccionario con {\ttfamily clear()}}
	
	El método {\ttfamily clear()} se puede usa para eliminar todos los elementos (clave-valor) de un diccionario, resultando un diccionario vacío: 
	
	\begin{lstlisting}
		In[]:	
		>>> diccionario = {"A": 1, "B":2,"C": 3} 
		>>> diccionario.clear()
		>>> diccionario
		{}
	\end{lstlisting}
	
	\noindent También se pueden eliminar todos le elementos del diccionario  asignándole  un diccionario vacío:
	\begin{lstlisting}
		In[]:	
		>>> diccionario = {}
		>>> diccionarios 
		{}
	\end{lstlisting}
	
	\noindent También se obtiene un diccionario vacío. 
	
	
	\subsection{Verificando la existencia de una clave con {\ttfamily in}}
	
	Para verificar si una clave existe en un diccionario, se puede usar la instrucción  {\ttfamily in}. Consideremos  nuevamente el diccionario con los miembros de Monty Python, esta vez dejamos por fuera a  "Gilliam Terry", y verifiquemos si se encuentra en el diccionario:
	
	\begin{lstlisting}
		In[]:	
		>>> grupo = {"Chapman": "Graham", "Cleese": "Juan", "Idle": "Eric", 
			"Jones": "Terry", "Palin":"Michael"}
		>>> "Chapman" in grupo
		True
		>>> "Gilliam" in grupo
		False
	\end{lstlisting}
	Con esta instrucción se prueba si una clave esta o no en el diccionario.  
	
	\subsection{Obtener un elemento con  {\ttfamily []} o  {\ttfamily get()}}
	
	Este es una de las instrucciones más usadas del diccionario. Para obtener el valor correspondiente se especifica el nombre del diccionario, seguido de la clave entre paréntesis:
	\begin{lstlisting}
		In[]:	
		>>> grupo = {"Chapman": "Graham", "Cleese": "Juan", "Idle": "Eric", 
			"Jones": "Terry", "Palin":"Michael","Gilliam": "Terry"}
		>>> grupo["Chapman"]
		'Graham'
	\end{lstlisting}
	
	
	Si la clave no existe en el diccionario, obtendrá una excepción, ejemplo:
	
	\begin{lstlisting}
		In[]:
		>>> grupo["Antonio"]
	\end{lstlisting}
	{\ttfamily\scriptsize\scriptsize\noindent Out[]: \vspace{-0.3cm}	
		{\color{red} 	
			\begin{verbatim}
				Traceback (most recent call last):
				File "<pyshell#46>", line 1, in <module>
				grupo["Antonio"]
				KeyError: 'Antonio'
			\end{verbatim}
	}}
	
	
	Hay dos buenas maneras de evitar este error. La primera es probar si la clave está en el diccionario con {\ttfamily in} ({\ttfamily key in diccionario}); y la segunda con el método {\ttfamily get()}, en este caso
	se recomienda dos argumentos, el primer argumento  la clave y el segundo argumento una opción incluida por el programador como valor, que se devolverá si la clave no existe. Visualicemos con algunos ejemplos:\\ 
	
	Se proporciona la clave, si existe en el diccionario, devuelve el valor correspondiente:
	\begin{lstlisting}
		In[]:	
		>>> grupo.get("Cleese")
		'John'
	\end{lstlisting}
	
	Ahora, especificando la clave y el valor opcional:\label{valor_predet}
	
	\begin{lstlisting}
		In[]:	
		>>> pythons.get("Manuel", "No es está")
		'No es está'
	\end{lstlisting}
	
	Si no especifica  el valor opcional, el resultado será nada, el intérprete interactivo no muestra nada:
	\begin{lstlisting}
		In[]:	
		>>> pythons.get("Manuel")
		>>> 
	\end{lstlisting}
	
	Significa que no está en el diccionario. Pero lo mejor e usar los dos argumentos, y evitar confusión.
	
	\subsection{Obtener claves por defecto}\label{predeterminado}
	
	Antes de explorar los diferentes enfoques de bucle, consideremos un diccionario diseñado para almacenar información sobre un usuario en un sitio web. El siguiente diccionario {\ttfamily usuario\_01}\label{usuario.01} almacenaría el nombre de usuario, el primer nombre y el apellido de una persona. Para acceder a las claves del diccionario, lo haremos con ayuda del bucle {\ttfamily for ... in},  el procedimiento es semejante al empleado en listas y tuplas:
	
	%\hspace{-2.5cm}{\ttfamily\scriptsize usuario.py}\vspace{-0.5cm}
	\begin{lstlisting}
		In[]:
		>>> usuario_01 = {'nombre de usuario': 'caRLos28', 
			'primero_nombre': 'cArlos',
			'apellido': 'DáVilA',}
		
		>>> for x in usuario_01:
		print(x)	          
	\end{lstlisting}
	
	Y se obtienen las claves:
	
	{\ttfamily\scriptsize\scriptsize\noindent Out[]:\vspace{-0.35cm}
		{\color{blue}			
			\begin{verbatim}
				nombre de usuario
				primer_nombre
				apellido
			\end{verbatim}
		}
	}
	
	
	
	Esta instrucción devuelve sólo las claves del diccionario, en realidad puede considerarse como un método
	predeterminado del diccionario.
	
	A continuación se presentan otros métodos más específico.
	
	
	\subsection{Obtener todas las claves con keys()}
	
	Con el método {\ttfamily  keys()}  se pueden obtener todas las claves de un diccionario en un objeto iterable tipo  {\ttfamily{}'{}dict\_keys{}'{}}   ({\ttfamily <class {}'{}dict\_keys{}'{}>{}}). Para ilustrar el uso de este método  consideremos el diccionario {\ttfamily transito}\label{transito}:
	
	\begin{lstlisting}
		In[]:	
		>>> transito = {"verde": "pasar", "amarillo": "prevenido", "rojo": "alto"}
		>>> transito.keys()
		dict_keys(['verde', 'amarillo', 'rojo'])
	\end{lstlisting} 
	
	Se obtienen  todas las claves del diccionario en un objeto iterable tipo  {\ttfamily dict\_keys()}; y para extraer las claves se puede usar un bucle {\ttfamily for ... in}. Ejemplo:
	
	\begin{lstlisting}
		In[]:
		>>> for clave in transito.keys():
		print(f"Clave: {clave.title()}")
	\end{lstlisting}
	
	{\ttfamily\scriptsize\scriptsize \noindent Out[]: \vspace{-0.3cm}	
		{\color{blue} 	
			\begin{verbatim}
				Clave: Verde
				Clave: Amarillo
				Clave: Rojo
			\end{verbatim}
	}}
	
	
	\subsubsection{Obtener las claves con un orden particular}
	
	Al extraer los elementos de un diccionario con cualquier método, estos se obtienen en el mismo orden en que
	fueron insertados. Sin embargo, algunas veces se necesitaran en un orden particular. 
	
	Par ordenar las claves se puede usar la función  {\ttfamily sorted()}; ilustremos el procedimiento con el diccionario {\ttfamily usuario.01} de  la página \pageref{usuario.01}
	\begin{lstlisting}
		In[]:
		>>> print("\n",sorted(usuario_01.keys()))
		['apellido', 'nombre de usuario', 'primer_nombre']
	\end{lstlisting}
	
	El resultado muestra las claves ordenadas en una lista. Para extraer la claves ordenadas se puede usar un {\ttfamily for .. in}.\\
	
	
	Consideremos un segundo ejemplo, para ilustrar el uso de este método.  Supongamos una lista de estudiantes encuestados por una universidad que ofrece ciertas carreras. Los estudiantes encuestados eligieron una carrera universitaria; y supongamos que estamos interesados saber los nombres de los estudiante que participaron en la encuesta,  qué carreras que eligieron, y generar una lista o diccionario con los que estudiantes que eligieron la opción Medicina.\\
	
	Con el siguiente código se extraen los nombres de las claves ordenados, y se presentan en formato de titulo:
	% sorted(encuesta)
	% ['Antonio', 'Erika', 'José', 'Juan', 'carlos', 'luis', 'manuela', 'maría']
	
	\hspace{-2.5cm}{\ttfamily \scriptsize encuesta.py}\label{encuesta} \vspace{-0.5cm}
	\begin{lstlisting}
		encuesta = {"luis":"Informática", "Antonio":"Estadística",
			"carlos":"Informática", "manuela":"Medicina",
			"Juan": "administración","maría":"Medicina", "Erika":"Medicina", 
			"José":"Matemática",}
		
		for nombre in sorted(encuesta):
		print(f"Nombre: {nombre.title()}")
	\end{lstlisting}	
	{\ttfamily\scriptsize\scriptsize 
		\hspace{-0.0cm}Out[]: \vspace{-0.3cm}	
		{\color{blue} 	
			\begin{verbatim}
				Nombre: Antonio
				Nombre: Erika
				Nombre: José
				Nombre: Juan
				Nombre: María
				Nombre: Carlos
				Nombre: Luis
				Nombre: Manuela
			\end{verbatim}
	}}
	
	\noindent El resultado,  los nombres de los participantes en la encuesta, ordenados y en formato de título.
	
	
	\subsection{Obteniendo todos los valores con {\ttfamily values()}} 
	
	Algunas veces sólo se necesitan los valores de un diccionario, con el método {\ttfamily values()} se obtienen unicamente los valores; este método  devuelve un objeto iterable tipo {\ttfamily{}'{}dict\_values{}'{}} ({\ttfamily <class {}'{}dict\_values{}'{}>{}}), como en el caso del método {\ttfamily keys}.\\
	
	Ejemplifiquemos con el diccionario de prueba {\ttfamily encuesta}, ahora estamos interesado es saber qué carreras eligieron: 
	
	\hspace{-2.5cm}{\ttfamily\scriptsize carreras.py}\vspace{-0.5cm}
	\begin{lstlisting}
		# encuesta.values() es un objeto tipo <class 'dict\_items'>
		for opcion in encuesta.values():
		print(opcion.title())
	\end{lstlisting}
	{\ttfamily\scriptsize\scriptsize 
		\noindent Out[]: \vspace{-0.3cm}	
		{\color{blue} 	
			\begin{verbatim}
				Informática
				Estadística
				Informática
				Medicina
				Administración
				Medicina
				Medicina
				Matemática
			\end{verbatim}
	}}
	
	Este enfoque extrae todos los valores del diccionario sin verificar repeticiones. Esto podría funcionar para una pequeña cantidad de valores, pero en una encuesta con un gran número de encuestados, el resultado es poco útil. Para ver cada opción sin repetición, podemos utilizar un conjunto. En Python, existe la función conjunto  {\ttfamily set()},  es una estructura de datos, como una colección desordenada de elementos que no admite duplicados, esto significa que los elementos de un conjunto son todos distintos; y además el objeto en iterable. En la sección \ref{conjuntos} se estudia esta función en detalle, por ahora sólo se aplica para lograr el resultado deseado. 
	
	Con la función {\ttfamily set()} aplicada al método {\ttfamily values()} se obtienen la colección de  valores sin duplicados:
	
	\begin{lstlisting}
		In[]:
		>>> for opcion in set(encuesta.values()):
		print(opcion.title())
	\end{lstlisting}
	{\ttfamily\scriptsize\scriptsize\noindent Out[]:{\color{blue}{\vspace{-0.5cm}}
			\begin{verbatim}	
				Informática
				Matemática
				Estadística
				Medicina
				Administración 	
			\end{verbatim}
	} }El resultado una colección de valores sin duplicados. \\
	
	
	%{'Informática', 'Matemática', 'Estadística', 'Medicina', 'administración'}
	
	
	Python identifica los elementos únicos en la colección y construye un conjunto con estos elementos. Aquí usamos la función  {\ttfamily set()}  para obtener los valores sin duplicados  en {\ttfamily encuesta.valores()}.
	
	
	
	\vspace{0.2cm}
	\begin{minipage}[t]{0.2\textwidth}
		\centering\raisebox{\dimexpr\topskip-\height}{%
			
\includegraphics[scale=0.20]{../Logos_de_Python/Busquedas_R}}  %[width=\textwidth]
		
	\end{minipage}%\hfill
	\begin{minipage}[t]{0.70\textwidth} 
		{\scriptsize 
			En Python 2, el método {\ttfamily keys()}  devuelve sólo una lista. Ahora, Python 3 devuelve un objeto iterable tipo {\ttfamily dict\_keys()}, que es una vista  de las claves. Esto es útil con diccionarios grandes, porque no utiliza ni
			tiempo, ni memoria para crear y almacenar una lista, que quizás no utilice más. Si necesita una lista en Python 3, 
			puede usar la función {\ttfamily list()} para convertir el objeto {\ttfamily dict\_keys()} en lista de claves, ejemplo:
			\begin{lstlisting}
				In[]:	
				>>> list(transito.keys())
				['verde', 'rojo', 'amarillo']
			\end{lstlisting}
			
			En Python 3, también necesitas usar la función {\ttfamily list()} para  convertir  el
			resultado de {\ttfamily values()}  y {\ttfamily items()} en listas de Python.
		}
	\end{minipage}
	\vspace{0.2cm}
	
	\subsection{Obtener todos los pares con {\ttfamily items()}}
	
	Cuando se necesita  obtener todos los elementos (clave-valor) del diccionario, se utiliza método {\ttfamily  items()}, 
	que devuelve un objeto iterable tipo {\ttfamily dict\_items} ({\ttfamily <class {}'{}dict\_items{}'{}>{}}). 
	Veamos con el diccionario {\ttfamily transito} de la página \pageref{transito}:
	
	\begin{lstlisting}
		In[]
		>>> transito.items()
		dict_items([('verde', 'pasar'), ('amarillo', 'prevenido'), ('rojo', 'alto'), 
		('azul', 'confundir a todos')])
		
		>>> list(transito.items())
		[('verde', 'pasar'), ('amarillo', 'prevenido'), ('rojo', 'alto')]
	\end{lstlisting}
	
	
	
	
	En este caso,  la función {\ttfamily list()} convierte el objeto {\ttfamily dict\_items} en una lista de tuplas con dos elementos,  el primer elemento de la tupla es la clave y el segundo elemento es el valor. Ahora con el  bucle {\ttfamily for ... in},  se le asigna una variable a las claves y otra a los valores para extraerlos.\\  %     
	
	
	
	
	
	% dict_items([('nombre de usuario', 'cArlos28'), ('primer_nombre', 'caRloS'), ('apellido', 'DáVila')])
	% ({\ttfamily diccionario.items()})
	
	Ilustremos un segundo ejemplo, ahora con el diccionario {\ttfamily usuario\_01} de la página \pageref{usuario.01}): Primero se crean nombres para las dos variables que contendrá la clave y el valor de cada item:  
	
	%\hspace{-2.5cm}{\ttfamily \scriptsize curso.py}  \vspace{-0.5cm}
	\begin{lstlisting}		
		In[]:
		>>> for clave, valor in usuario_01.items():
		print(f"\nClave: {clave}")
		print(f"Valor: {valor}")
	\end{lstlisting}	
	
	\noindent El código funciona  con cualquier nombre que asigne a estas dos variables; por ejemplo,  si usas abreviaturas para los nombres de las variables, como  {\ttfamily k} para la clave, y  {\ttfamily v}  para el valor.\\
	
	
	Se definen las variables clave y valor,  y se  imprimen en una {\ttfamily{f-cadena}}. El carácter  {\ttfamily "{}\textbackslash n{}"{}} indica de salto de línea en Python, para marcar el fin de una línea y el comienzo de una nueva línea.  En este ejemplo se insertó una línea antes de cada par clave-valor en la salida:
	
	{\ttfamily\scriptsize\scriptsize 
		\hspace{-0.65cm}Out[]: \vspace{-0.3cm}	
		{\color{blue} 	
			\begin{verbatim}
				
				Clave: nombre de usuario
				Valor: cArlos28
				
				Clave: primer_nombre
				Valor: caRloS
				
				Clave: apellido
				Valor: DáVila	
			\end{verbatim}
	}}
	
	En vista que los datos ingresado no son confiable respecto al formato de título, le pedimos la salida con
	formato de título para cada clave y valor en una {\ttfamily f-cadena}:  
	
	\begin{lstlisting}
		In[]:  
		>>> for clave, valor in usuario_01.items():
		print(f"\nClave: {clave.title()}")
		print(f"Valor: {valor.title()}")
	\end{lstlisting}
	{\ttfamily\scriptsize\scriptsize 
		\hspace{-0.15cm}Out[]: \vspace{-0.3cm}	
		{\color{blue} 	
			\begin{verbatim}
				
				Clave: Nombre De Usuario
				Valor: Carlos28
				
				Clave: Primer_Nombre
				Valor: Carlos
				
				Clave: Apellido
				Valor: Dávila
			\end{verbatim}
	}}
	
	
	Para obtener un diccionario con los estudiantes que eligieron medicina en la {\ttfamily encuesta} prueba de la página \pageref{encuesta}; primero se crea un diccionario vacío con el nombre de {\ttfamily medicina}, luego con un {\ttfamily for} y un {\ttfamily if}, se construye el diccionario {\ttfamily medicina}  con los estudiantes que eligieron  Medicina,  el código se expone a continuación:
	
	\hspace{-2.5cm}{\ttfamily \scriptsize medicina.py}  \vspace{-0.5cm}
	\begin{lstlisting}
		medicina = {}
		for x,y in encuesta.items():
		x1 = x.title()
		if y =="Medicina":
		medicina[x1] = "Medicina"
		print(medicina)  
	\end{lstlisting}	
	{\ttfamily\scriptsize\scriptsize 
		\hspace{-0.0cm} Out[]:\vspace{-0.3cm}	
		{\color{blue} 	
			\begin{verbatim}
				{'Manuela': 'Medicina', 'María': 'Medicina', 'Erika': 'Medicina'}
			\end{verbatim}
	}}
	Para obtener la lista puede usar el código semejante o simplemente usar la función {\ttfamily list()}, obtendrá una lista de los que seccionaron Medicina. 
	
	
	
	
	\subsection{Asignar con =, y Copiar con copy()}
	
	Cuando se hacen cambios en un diccionario, estos se reflejará en todos los diccionarios que se formaron a partir de este, veamos:
	
	\begin{lstlisting}
		In[]:	
		>>> transito = {"verde": "pasar", "amarillo": "prevenido", "rojo": "alto"}
		>>> save_transito = transito
		# Generamos un cambio agregando un elemento:
		>>> transito["azul"] = "confundir a todos"
		>>> save_transito
		{'verde': 'pasar', 'amarillo': 'prevenido', 'rojo': 'alto', 
			'azul': 'confundir a todos'}
	\end{lstlisting}
	
	\noindent El diccionario {\ttfamily save\_transito} sufrió los cambios, del diccionario  {\ttfamily transito}.  Para evitar este inconveniente, puede usar el método {\ttfamily copy()} que hace  copias del diccionario a otro diccionario independiente:
	
	\begin{lstlisting}
		In[]:	
		>>> transito = {"verde": "pasar", "amarillo": "prevenido", "rojo": "alto"}
		>>> transito_original = transito.copy()
		>>> transito["azul"] = "confundir a todos"
		>>> transito
		{'verde': 'pasar', 'amarillo': 'prevenido', 'rojo': 'alto',
			'azul': 'confundir a todos'}
		>>> transito_original
		{'verde': 'pasar', 'amarillo': 'prevenido', 'rojo': 'alto'}
	\end{lstlisting}
	
	Con esta instrucción se mantiene el diccionario original en la copia sin ser afectado por cambios en el futuro.  
	
	
	
	
	
	
	
	\vspace{0.2cm}
	\begin{minipage}[t]{0.2\textwidth}
		\centering\raisebox{\dimexpr\topskip-\height}{%
			
\includegraphics[scale=0.20]{../Logos_de_Python/Advertencia}}  %[width=\textwidth]
		
	\end{minipage}%\hfill
	\begin{minipage}[t]{0.70\textwidth} 
		{\scriptsize Es importante no confundir conjuntos con diccionarios. Puede prestar confusión debido a que ambos objetos usan las llaves; sin embargo, los item o elementos en diccionarios son pares clave-valor, mientras que los elementos de los conjuntos son individuales, y además no obedecen a un orden específico, simplemente son conjuntos.  
		}
	\end{minipage}
	\vspace{0.2cm}
	
	
	
	
	\section{Ejercicios}\label{Tarea Dic }
	
	Ejercicios Diccionarios
	
	\begin{enumerate}%[6.3.1]
		
		\item  Persona: Use un diccionario para almacenar información sobre una persona que conoce.
		
		Almacene su nombre, apellido, edad y la ciudad en la que vive. Debe tener claves como nombre, apellido, edad y ciudad. 
		Imprime cada clave de información almacenada en su diccionario.
		
		
		\item Números favoritos: Utilice un diccionario para almacenar los números favoritos de las personas.
		Piensa en cinco nombres y úsalos como claves en tu diccionario. Piensa en un número favorito  para cada persona y 
		almacene cada uno como un valor en su diccionario. Imprimir el nombre de cada persona y su número favorito. 
		Para divertirse aún más, encuesta algunos amigos y obtenga algunos datos reales para su programa.
		
		
		\item Glosario 1: Se puede utilizar un diccionario de Python para modelar un diccionario real.
		Sin embargo, para evitar confusiones, llámelo glosario.
		
		Piensa en cinco palabras de programación que aprendiste. Utilice estas palabras como
		claves en su glosario y almacene sus significados como valores.
		
		Imprima cada palabra y su significado como resultado cuidadosamente formateado.
		Podrías imprimir la palabra seguida de dos puntos y luego su significado, o imprima la palabra
		en una línea y luego imprima su significado con sangría en una segunda línea. Utilizar un nueva línea ({\bf\ttfamily{\textbackslash}}n)  para insertar una línea en blanco entre cada salida de significado de palabra.
		
		
		\item  Glosario 2: Ya sabes cómo recorrer un diccionario en bucle, actualice el código del Ejercicio Glosario 1,
		reemplazando la serie de llamadas {\ttfamily{print()}} con un bucle que recorre las claves y valores del diccionario.
		Cuando estés seguro de que el bucle funciona, agrega cinco términos más de Python a tu glosario. 	Cuando ejecute su programa nuevamente, estas nuevas palabras y significados deberían incluirse automáticamente en la salida.
		
		\item Ríos: Haga un diccionario que contenga tres ríos principales y el país por donde pasa cada río. 
		Un par {\ttfamily{clave-valor}} podría ser 'nilo': 'egipto'.
		
		Utilice un bucle para imprimir una oración sobre cada río, como: El Nilo corre 	a través de Egipto.
		
		Utilice un bucle para imprimir el nombre de cada río incluido en el diccionario.
		
		Utilice un bucle para imprimir el nombre de cada país incluido en el diccionario.
		
		
		\item Encuesta: Use el código en {\ttfamily{encuesta.py}}. 
		Haga una lista de personas que deberían realizar la encuesta de estudios universitarios. Incluya  algunos nombres que ya están en el diccionario y otros que no. 
		
		Recorra la lista de personas que realizaran la encuesta. Si alguno ya participó en la encuesta, imprima un mensaje agradeciéndoles por responder. Si aún no han realizado la encuesta, imprima un mensaje invitándolos a participar.
		
		\item Personas: Comience con el programa que escribió  en ejercicio Persona:.
		Haga dos nuevos diccionarios que representan a personas diferentes, y almacena los tres diccionarios
		en una lista llamada personas. Recorre tu lista de personas. 
		
		Mientras recorre la lista con un bucle, imprima todo lo que conoce sobre cada persona.
		
		\item Mascotas: Haga varios diccionarios, donde cada diccionario represente una mascota diferente.
		En cada diccionario incluye la clase de animal y el nombre del dueño.
		Guarde estos diccionarios en una lista llamada mascotas. 
		A continuación, recorra su lista con un bucle, e imprime todo lo que conoce de  cada mascota.
		
		\item  Lugares favoritos: Haga un diccionario llamado {\ttfamily{lugares\_favoritos}}. Piensa en tres
		nombres para usar como claves en el diccionario y almacenar de uno a tres lugares favoritos para
		cada persona. Para hacer este ejercicio un poco más interesante, pídele a algunos amigos que 
		nombre algunos de sus lugares favoritos. Recorra el   diccionario con un bucle e imprime el nombre de cada persona y sus lugares favoritos.
		
		\item Números favoritos: Modifique su programa del Ejercicio 6.2 (página 98) para que 
		cada persona tenga  más de un número favorito. Luego imprima el nombre de cada persona  junto con sus números favoritos.
		
		
		\item  Ciudades: Haga un diccionario llamado ciudades. Usa los nombres de tres ciudades como
		claves en tu diccionario. Crear un diccionario de información sobre cada ciudad, 
		incluya el país en el que se encuentra la ciudad, su población aproximada y un dato
		sobre esa ciudad. Las claves para el diccionario de cada ciudad deben ser algo como
		país, población y hechos. Imprime el nombre de cada ciudad y toda la información almacenada.
		
		
		\item  Extensiones: Ahora estamos trabajando con ejemplos que son lo suficientemente complejos, que se pueden ampliar de muchas maneras. Utilice uno de esos programas ejemplares de este capítulo y amplíelo  agregando nuevas claves y valores, cambiando el contexto del programa, o mejorar el formato de la salida.
		
		
		
	\end{enumerate}
	
	
	
	
	
	\section{Conjuntos}\label{conjuntos}
	
	\begin{wrapfigure}{l}{0.48\textwidth} %  lado derecho, izquierdo, L y l R y r
		\vspace{-25pt} \begin{center} \includegraphics[width=0.48\textwidth]{../Logos_de_Python/conjuntoA1}
			% angle=45, conjuntos05
		\end{center} \vspace{-29pt}	 %\caption{Python 3} 
	\end{wrapfigure} La sintaxis de un conjunto es como un diccionario pero sin valores, sólo con claves, y como en un diccionario cada clave es única; significa que no se repiten elementos (clave-valor); el orden  no es importante, es decir, en los diccionarios no hay un orden específico en sus elementos.
	
	Generalmente la teoría básica de conjuntos se enseña en la escuela primaria,  junto con las matemáticas elementales, y forma parte del conocimiento común. En esta sección se comenta las definiciones y operaciones entre conjunto, no como un tema sino como ideas aclaratorias por si no recuerda esas nociones. En la Figura \ref{union} muestra las ideas gráficas e  intuitivas  de unión e intersección; si requiere de más información puede consultar en  el texto Teoría de conjuntos de la Serie Schaum\footnote{Teoría De Conjuntos Y Temas Afines. Seymur Lipschutz. Schaum Mc Graw Hill}
	\begin{figure}[h!]
		\centering
		\includegraphics[scale=0.55]{../Logos_de_Python/union03}
		\caption{Operaciones con conjuntos}
		\label{union}
	\end{figure} 
	%  https://www.programaenpython.com/fundamentos/sets-en-python/#Diferencia
	
	En Python un conjunto es tratado bajo la misma idea de conjunto en matemáticas, colecciones de elementos sin un orden establecido, y sin elementos duplicados; con sus operaciones de unión intersección , diferencia de cojuntos, diferencia simétrica (figura \ref{union01}). Y por supuesto la existencia del conjunto nulo o conjunto vacío, que es un conjunto con cero elementos. 
	
	Supongamos que debemos unir dos conjuntos que tienen algunas elementos en común, en Python la unión de los dos conjuntos contendrá los elemento de los dos conjunto sin duplicados; como se establece en un conjunto en  matemática. Las demás operaciones con conjunto en Python siguen las reglas de matemáticas. \\
	
	En las próximas secciones trataremos la construcción de conjuntos en Python y las operaciones. 
	
	% Pendiente En la Figura \ref{vacio}, un ejemplo de un conjunto nulo serían ejemplo,  ... :
	%
	
	\subsection{Crear conjuntos con  \{ ...\} o con set()}
	
	%  Python Built-in Functions
	Para crear un conjunto, los elementos  o valores se encierran entre llaves separados por comas, como se muestra a continuación, con algunos ejemplos de conjuntos: 
	\begin{lstlisting}
		>>> Números_pares = {0, 8, 2, 4, 6}
		>>> Números_pares
		{0, 2, 4, 6, 8}
		
		>>> Números_impares = {9, 3, 1, 5, 7}
		>>> Números_impares
		{1, 3, 5, 7, 9}
		
	\end{lstlisting}
	
	Al igual que las claves en un diccionario que no siguen un orden;  en los conjuntos los elementos no tienen un orden, se pueden construir sin orden.\\
	
	
	En Python el conjunto vacío se codifica como {\ttfamily set()}, y no como es de esperar, llaves sin elementos, esa codificación Python la reserva para el diccionario vacío: 
	
	\begin{lstlisting}
		>>> conjunto_vacío = set()
		>>> conjunto_vacío
		set()
		>>> objeto = { }
		>>> type(objeto)
		<class 'dict'>
	\end{lstlisting}
	
	\vspace{0.2cm}
	\begin{minipage}[t]{0.2\textwidth}
		\centering\raisebox{\dimexpr\topskip-\height}{%
			
\includegraphics[scale=0.20]{../Logos_de_Python/Busquedas_R}}  %[width=\textwidth]
		
	\end{minipage}%\hfill
	\begin{minipage}[t]{0.70\textwidth} 
		{\scriptsize 	
			Naturalmente se puede intuir que {\ttfamily  \{  \}} representa un conjunto vacío, pero Python usa esta notación para un diccionario vacío, y no para un conjunto; para el conjunto vacío usa {\ttfamily set()}. El intérprete imprime el conjunto vacío como {\ttfamily set()} en lugar de {\ttfamily  \{  \}}.
			El motivo se debe a que los diccionarios ya estaban definidos en Python (primero que los conjuntos) y tomaron posesión de las llaves.
		}
	\end{minipage}
	\vspace{0.2cm}
	
	
	La función {\ttfamily set()} también se  utiliza para construir conjuntos y con mayor frecuencia por ser automática. La próxima sección es dedica a esta función.
	
	\subsection{Convertir un tipo objeto en conjunto con {\ttfamily set()}}
	
	Se pueden crear conjuntos a partir de una cadena, tupla, lista,  o diccionario, con la función {\ttfamily set()};
	automáticamente la función retorna un conjunto (sin repetir elementos). \\
	
	
	Veamos un ejemplo con una cadena con espacios y algunas letras repetidas:
	\begin{lstlisting}
		>>> conjunto = set("letras y algo más")
		>>> print(conjunto)
		{'s', 'á', 't', 'o', 'e', 'g', 'y', 'r', 'a', 'm', 'l', ' '}
	\end{lstlisting}
	Observe, la letra {}"{}a{}"{} aparece tres veces en la cadena, pero en el conjunto sólo una vez; igual para los espacios en blanco, en la cadena hay tres, en el conjunto sólo uno y al final. \\
	
	
	Construyamos un conjunto a partir de una lista:
	\begin{lstlisting}
		>>> lista = ["José", "María", "Vehículo", "moto", "Bicicleta", "Carlos", 
		"María"]
		>>> set(lista)
		{'moto', 'María', 'Bicicleta', 'Vehículo', 'Carlos', 'José'}
	\end{lstlisting}
	Observe, "María{}"{} aparece dos veces en la lista, mientras que en el conjunto una sola vez.\\ 
	
	Ahora, un conjunto a partir de una tupla:
	\begin{lstlisting}
		>>> tupla = ("Vehículo", "moto", "Bicicleta", "José","María", "Carlos",
		"María")
		>>> set(tupla)
		{'moto', 'María', 'Bicicleta', 'Vehículo', 'Carlos', 'José'}
	\end{lstlisting}
	Igualmente  "María{}"{} aparece dos veces en la tupla, mientras que en el conjunto una sola.\\ 
	
	
	Esta vez, construimos un conjunto a partir un diccionario; en este caso la función {\ttfamily set()} solo toma las claves del diccionario para formar el conjunto:
	\begin{lstlisting}
		>>> diccionario = {"Manzana":'verde', "Naranja":'naranja', "Cereza":'rojo'}
		>>> set(diccionario)
		{'Manzana', 'Cereza', 'Naranja'}
	\end{lstlisting}
	\noindent Resultado, un conjunto con las claves del diccionario.
	
	\subsection{Operaciones con conjuntos, intersección, unión, diferencia}
	
	Para implementar estos conceptos en Python, disponemos de los operadores de comparación habituales:\vspace{0.2cm}
	
	{\bf Operadores con conjunto en Python}
	\begin{lstlisting}
		== igual
		!= distinto
		>  mayor
		>= mayor o igual
		<  menor
		<= menor o igual
	\end{lstlisting}
	
	
	% Objetivo: implementar los conceptos de subconjuntos y conjuntos disjuntos en Python.
	
	
	
	Para ejemplificar las operaciones de conjunto en Python, consideremos los conjuntos de prueba {\ttfamily $A$ } y {\ttfamily $B$ }. Por el momento abusaré del uso mayúsculas para representar los conjuntos, sin embargo, para seguir la PEP 8, debe ser minúsculas:
	
	
	\hspace{-2.7cm}{\ttfamily \scriptsize conjuntos.py}  \vspace{-0.5cm}
	\begin{lstlisting}		
		A, B = {1, 3, 5, 7, 9}, {2, 3, 5, 7}
	\end{lstlisting}
	
	Para interceptar conjuntos en Python se puede usar el operador   \& 
	o  el método {\ttfamily intersection()} de conjuntos; tenga presente que al interceptar conjuntos el resultado son los elementos comunes en ambos conjuntos.
	\begin{figure}[h!]
		\centering
		\includegraphics[scale=0.35]{../Logos_de_Python/union01}
		\caption{Operaciones con conjuntos}
		\label{union01}
	\end{figure}Ilustremos los ejemplos de prueba: 
	
	
	\hspace{-3cm}{\ttfamily \scriptsize operaciones.py}  \vspace{-0.5cm}
	\begin{lstlisting}
		# intersección con & 
		>>> A & B
		{3, 5, 7}
		
		# Con el método:
		>>> A.intersection(B)
		{3, 5, 7}
	\end{lstlisting}
	
	\noindent  Resultado de esta operación, un conjunto con los elementos comunes en ambos conjuntos(figura \ref{union01}).\\
	
	
	Para unir conjuntos se puede usar el operador {\ttfamily |} o el método {\ttfamily union()} de conjuntos. La operación unión de conjuntos devuelve todos los elementos entre los conjuntos (sin repetir),  nuevamente ilustremos con los ejemplos de prueba:
	
	\begin{lstlisting}
		# unión con el operrador | 	
		>>> A | B
		{1, 2, 3, 5, 7, 9}
		# Union con el método
		>>> A.union(B)
		{1, 2, 3, 5, 7, 9}
	\end{lstlisting}
	
	
	Para la operación diferencia  entre conjuntos (ver gráfico \ref{union01}, los que están el primer conjunto y no en el segundo)  se puede   usar símbolo de resta {\ttfamily -}  o el método {\ttfamily difference()} de conjuntos. Ilustramos estas operaciones con los conjuntos de prueba:
	
	\begin{lstlisting}
		>>> A-B
		{1, 9}
		>>> A.difference(B)
		{1, 9}
	\end{lstlisting}
	
	La diferencia entre el conjunto {\ttfamily A} y  el conjunto {\ttfamily  B}, es el conjunto {\ttfamily \{1, 9\}}, como se muestra en el gráfico (\ref{union01}). \\
	
	
	
	Las operaciones de conjuntos más comunes son unión, intersección y diferencia; sin embargo, se incluyen la  diferencia simétrica (exclusive or) por si en algún momento lo necesita, y haremos un ejemplo de su aplicación en la sección \ref{aplicaciones}.\\
	
	Para la diferencia simétrica (ver gráfico \ref{union01}, los elementos del primer conjunto que no están en el segundo, y los elementos del segundo conjunto que no están en el primero conjunto)  se usa como operador el símbolo {\ttfamily \^{}}, o  el método {\ttfamily symmetric\_difference()} de conjuntos, ilustremos:
	
	\begin{lstlisting}
		>>> A^B
		{1, 3}
		>>> A.symmetric_difference()(B)
		{1, 3}
	\end{lstlisting}
	Resultado el conjunto \{ 1,3 \}, como se muestra en la figura (\ref{union01}).
	
	
	\subsection{Subconjuntos}
	
	Algunas veces es necesario trabajar con subconjuntos ( {\ttfamily A} está contenido en  {\ttfamily B}, todos los elementos del primer conjunto {\ttfamily A}  están en el segundo conjunto {\ttfamily B},  figura \ref{subconjunto})
	\begin{figure}[h!]
		\centering
		\includegraphics[width=0.45\linewidth, height=0.225\textheight]{../Logos_de_Python/subconjunto01}
		\caption{Operaciones con subconjuntos}
		\label{subconjunto}
	\end{figure} por ejemplo; podemos necesitar comprobar si un conjunto {\ttfamily A} es subconjunto de otro {\ttfamily B ($A\subseteq B$)}. El operador {\ttfamily <=} se usa para verificar si el conjunto  {\ttfamily A} es  subconjunto {\ttfamily B}, el operador devuelve el valor {\ttfamily True} en caso de ser cierto,  o  {\ttfamily False}  {\ttfamily A} si no está. Python también tiene el método {\ttfamily issubset()}, que hace el mismo trabajo, y devuelve el mismo resultado. Veamos con los conjuntos de prueba:
	\begin{lstlisting}
		>>> A <= B
		False
		>>> A.issubset(B)
		False
	\end{lstlisting}
	Decir que {\ttfamily A} está contenido en  {\ttfamily B}  ($A\subseteq B$), es lo mismo que decir  {\ttfamily B} contiene a {\ttfamily A}  ($B\supseteq A$), y Python lo codifica como un superconjunto (superset) o conjunto propio de B:  % $ \supseteq$
	\begin{lstlisting}
		>>> B >= A
		False
		>>> B.issuperset(A)
		False
	\end{lstlisting}
	
	En matemáticas se sabe que todo conjunto es subconjunto de sí mismo, y Python también lo codifica con  {\ttfamily A.issubset(A)}  y  {\ttfamily A <= A }, y ambos códigos devuelve el valor booleano True.
	
	\subsubsection{Inclusión estricta}
	
	Un conjunto A está estrictamente incluído\footnote{Si A es un subconjunto de B tal que A $\neq$ B, entonces se dice que A es un {\bf subconjunto propio} de B, y se denota por A $\subset$ B. A su vez, se dice que B es un {\bf superconjunto propio} de A.} en B, si y sólo si todo elemento de A pertenece a B, pero existe por lo menos un elemento de B que no pertenece a A; es decir, $A \subseteq B$, pero A $\neq$ B,  en ese caso escribimos A $\subset$ B, de aquí se concluye que {\ttfamily A} no está estrictamente incluído en {\ttfamily A}.
	
	Si necesita la inclusión estricta {\ttfamily A $\subset $ B}, Python tiene el operador {\ttfamily <{}} como inclusión estricta,  veamos con algunos ejemplos:
	\begin{lstlisting}
		>>> A < A
		False
	\end{lstlisting}
	
	Decir que  A $\subset$ B, es lo mismo que decir B $ \supset$ A.  Ejemplifiquemos estas operaciones.
	\begin{lstlisting}
		>>> A < B
		False
		>>> A > B
		False
	\end{lstlisting}
	
	Ahora, definimos un nuevo conjunto de prueba {\ttfamily C}:
	
	\begin{lstlisting}  
		>>> C = {1,3,5}	
		>>> C < A
		True
		>>> A > C
		True 
		>>> B < C
		False
		>>> C > B
		False
	\end{lstlisting}
	
	
	
	\subsection{Conjuntos inmutables Frozenset}
	
	%  A \ B diferncia conjunto  Sintaxis frozenset(iterable)
	
	En Python los conjuntos son mutables, pero existe la posibilidad de crear conjuntos inmutables, con la  función {\ttfamily frozenset()}, una función interna  de Python. Su sintaxis es sencilla:  {\ttfamily frozenset(argumento)},  admitiendo como argumento cualquier objeto iterable. 
	
	La diferencia fundamental con respecto a los conjuntos ({\ttfamily set()}) es que son objetos inmutables; lo que significa que una vez creados no permiten cambios;  otra diferencia con los conjuntos es que ocupan menos espacio;  el resto de sus características son similares. El conjunto vacío  {\ttfamily frozenset} se codifica como {\ttfamily frozenset()};  la función sin argumento, crea el conjunto  {\ttfamily frozenset}  vacío. Ilustremos con algunos ejemplos:
	
	Primero consideremos una lista de números enteros:
	%\hspace{-2.5cm}{\ttfamily \scriptsize curso.py}  \vspace{-0.5cm}
	\begin{lstlisting}	
		In[]:	
		>>> lista_de_números = [0, 3, 2, 3, 1, 4, 0, 4, 3]		
		>>> frozenset(lista_de_números)
		frozenset({0, 1, 2, 3, 4})
		>>> type(frozenset({0, 1, 2, 3, 4}))
		<class 'frozenset'>
	\end{lstlisting}	
	El resultado un conjunto tipo {\ttfamily frozenset <class {}'{}frozenset{}'{}>{}}.\\
	
	
	Ahora consideremos una tupla de números 
	\begin{lstlisting}		
		In[]:
		>>> tupla_de_números = (0, 3, 2, 3, 1, 4, 0, 4, 3)	
		>>> frozenset(tupla_de_números)
		frozenset({0, 1, 2, 3, 4})
	\end{lstlisting}	
	
	El resultado un conjunto  {\ttfamily frozenset}, como se ve, es el mismo creado para lista y tuplas. \\
	
	
	
	
	A partir de una cadena de caracteres también se puede generar un conjunto frozenset:
	\begin{lstlisting}	
		In[]:	
		>>> cadena = "Conjunto de caracteres"	
		>>> frozenset(cadena)
		frozenset({'u', 'r', 'a', 'd', 'e', 'o', 'C', ' ', 'j', 't', 'c', 's', 'n'})
	\end{lstlisting}	
	
	\noindent Observe, un conjunto frozenset con los caracteres de la cadena, incluido los espacios en blanco.\\
	
	Los frozenset son inmutables no aceptan métodos que traten de alterar su composición, si intenta usar un método que que genere cambios  provocará una excepción en Python, indicando que  el objeto frozenset no tiene método ese método. Pero los conjuntos {\ttfamily frozenset()}>  heredan todas las operaciones de los  conjuntos {\ttfamily set()}.  Consideremos los siguientes los conjuntos  de prueba forzenset {\ttfamily A} y {\ttfamily B}, para ejemplificar:
	
	\hspace{-2.7cm}{\ttfamily \scriptsize congelado.py}  \vspace{-0.5cm}
	\begin{lstlisting}
		A, B = frozenset({1, 3, 5, 7, 9}), frozenset({2, 3, 5, 7})
		C = frozenset({3,5,1})
		# intersección 
		print("\nintersección: ",A&B)
		print("intersección: ",A.intersection(B))
		
		# Union 
		print("\nUnión: ", A|B)
		print("Unión: ", A.union(B))
		
		# Diferencia de conjuntos
		print("\nDiferencia: ", A-B)
		print("Diferencia: ", A.difference(B))
		
		# Diferencia simétrica
		print("\nDiferencia simétrica: ",A ^ B)
		print("Diferencia simétrica: ",A.symmetric_difference(B))
		
		# Comparaciones 
		print("\nA está contenido en B?:", A<=B)
		print("C está  en B?:", C<=A)      
	\end{lstlisting}
	{\ttfamily\scriptsize\scriptsize 
		\hspace{-0.65cm} Out[]:\vspace{-0.3cm}	
		{\color{blue} 	
			\begin{verbatim}		
				
				intersección:  frozenset({3, 5, 7})
				intersección:  frozenset({3, 5, 7})
				
				Unión:  frozenset({1, 2, 3, 5, 7, 9})
				Unión:  frozenset({1, 2, 3, 5, 7, 9})
				
				Diferencia:  frozenset({1, 9})
				Diferencia:  frozenset({1, 9})
				
				Diferencia simétrica:  frozenset({1, 2, 9})
				Diferencia simétrica:  frozenset({1, 2, 9})
				
				A está contenido en B?: False
				C está  en B?: True
			\end{verbatim}
	}}
	
	\vspace{0.2cm}
	\begin{minipage}[t]{0.2\textwidth}
		\centering\raisebox{\dimexpr\topskip-\height}{%
			
\includegraphics[scale=0.20]{../Logos_de_Python/Busquedas_R}}  %[width=\textwidth]
		
	\end{minipage}%\hfill
	\begin{minipage}[t]{0.70\textwidth} 
		{\scriptsize Dado que los conjuntos tanto {\ttfamily set} como {\ttfamily frozenset} no están ordenados, no admiten operaciones de indexación ni  división.\vspace{0.16cm}
			
			
			Observe que los conjuntos {\ttfamily frozenset} son a los conjuntos {\ttfamily sets} como las tuplas a las listas: una forma de «congelar los elementos en objetos tuplas o conjuntos {\ttfamily frozenset}, conservando los valores sin la opción a modificar.
			
		}
	\end{minipage}
	\vspace{0.2cm}
	
	
	
	
	
	
	
	
	%La sección “Iterar múltiples secuencias con zip()” en la página 79 presenta una función llamado zip() que facilita la creación de estas secuencias de dos elementos.
	
	
	%  Create with [] or list() o compresion
	%%
	
	
	
	
	\subsection{La función {\ttfamily zip()} }
	
	
	%  https://ellibrodepython.com/zip-python#zip-con-diferentes-longitudes
	
	La función {\ttfamily zip()}, es una función integrada o incorporada de Python (no necesita importarse). Se usa para construir un objeto Iterable tipo {\ttfamily {}<{}class {}'{}zip{}'{}>{} }, cuyos elementos son tuplas  que están formados por pares, triples, cuádruples, etc; es decir, secuencias múltiples dependiendo del numero de iterables que se ingresen como argumento.   Si la entrada son dos iterables, el objeto resultante sera una secuencia iterable de tuplas con dos elementos (tuplas de pares), y cada tupla de  estas contienen un elemento de cada uno de los iteradores originales; si la entrada son tres iterables, el objeto resultante sera una secuencia iterable de tuplas con tres elementos; y  así sucesivamente.\\ 
	
	La sintaxis es sencilla, el argumento son dos o más iterables que la función comprime en un objeto iterable tipo {\ttfamily zip} (secuencias múltiples), con la particularidad de que sus elementos son tuplas que contienen un elemento de cada uno de los iteradores originales (el primero con el primero, el segundo con el segundo, y así sucesivamente).  Para extraer los elementos del objeto {\ttfamily zip} puede usar un  {\ttfamily for ... in}, puede extraer las tuplas o los elementos individuales de las tuplas.
	
	Para ejemplificar, consideremos el siguiente código en el archivo {\ttfamily zip\_01.py}:
	
	
	\hspace{-2.5cm}{\ttfamily \scriptsize zip\_01.py}  \vspace{-0.5cm}
	\begin{lstlisting}		
		lista1 = [0,3,5,9]
		lista2 = [2,4,7,1]
		lista3 = ["a","b","c","d"]
		
		zip(lista1,lista2)	
		print(type(zip("Tipo de objeto:", lista1,lista2))) # Tipo de objeto.
		
		zip02 =  zip(lista1,lista2,lista3)	
		
		for t in zip02:
		print(t)  
	\end{lstlisting}	
	
	{\ttfamily\scriptsize\scriptsize 
		\hspace{-0.65cm}Out[]: \vspace{-0.3cm}	
		{\color{blue} 	
			\begin{verbatim}
				Tipo de objeto: <class 'zip'> 
				
				(0, 2, 'a')
				(3, 4, 'b')
				(5, 7, 'c')
				(9, 1, 'd')
			\end{verbatim}
	}}
	
	En este archivo se construye dos iterables tipo {\ttfamily zip}, el primero con dos listas, e imprimimos el typo de objeto; el segundo con tres listas, y se extraen los elementos del objeto (las tuplas) construido  {\ttfamily zip(lista1,lista2,lista3)}.\\
	
	Consideremos un segundo ejemplo, ahora una tupla y una lista: 
	
	
	\hspace{-2.5cm}{\ttfamily \scriptsize zip\_01.py}  \vspace{-0.5cm}
	\begin{lstlisting}
		# Una tupla
		paises = ("Venezuela", "china", "India", "Estados Unidos")
		# Una lista
		poblaciones = [30,1391, 1364, 327]
		
		for pais, poblacion in zip(paises, poblaciones):
		print("{}: {} millones de habitantes.".format(pais, poblacion))
	\end{lstlisting}
	{\ttfamily\scriptsize\scriptsize \noindent Out[]: \vspace{-0.3cm}	
		{\color{blue} 	
			\begin{verbatim}
				Venezuela: 30 millones de habitantes.
				china: 1391 millones de habitantes.
				India: 1364 millones de habitantes.
				Estados Unidos: 327 millones de habitantes.
			\end{verbatim}
	}}
	
	
	
	En este código se usaron dos iterables diferentes, una tupla y una lista, para mostrar que el argumento de la función trabaja bien con cualquier iterable, incluyendo combinaciones.  Se extrajeron los elementos individuales de las tuplas en 
	{\ttfamily zip(paises, poblaciones) },  y se imprimen con una oración en formato {\ttfamily  str.format()} ({\ttfamily cadena.format()}). También  puede dar formato con una {\ttfamily f{}"{}cadena{}"{}} e imprimir como: {\ttfamily print(f{}"{}\{pais.title()\}:\{poblacion\} millones de habitantes.{}"{})},  obtiene el mismo resultado, pero con menos esfuerzo manual.\\
	
	
	Por último también puede usar la función {\ttfamily  zip()}  con argumentos de iterables de diferentes longitudes.
	En este caso  el iterador {\ttfamily  zip()} finaliza con el iterador  más pequeña que ingre en en argumento. Usted puede hacer pruebas y comprobar.
	
	
	Y como ejercicio,   use la función {\ttfamily list()} para convertir el objeto {\ttfamily zip} en una lista (lista de tuplas) como: {\ttfamily list(zip(argumento))}. \\ 
	
	
	
	\subsubsection{Crear Diccionarios con {\ttfamily zip()} como argumento} 
	
	La función {\ttfamily zip()}  proporciona un medio sencillo para crear diccionarios, usando esta función, conjuntamente con la función {\ttfamily dict()}, veamos un ejemplo para ilustrar este procedimiento. 
	
	Considere los iterables de la sección anterior {\ttfamily paises} (que se corresponderá con la clave de cada elemento del diccionario) y {\ttfamily poblaciones}    (que se corresponderá con el valor de cada elemento del diccionario). 
	
	Ahora usamos la función {\ttfamily dict()} que tomará como argumento la función {\ttfamily zip()}, que a su vez, tomará las listas {\ttfamily paises} y {\ttfamily  poblaciones}  para generar los pares clave-valor que conformará cada elemento el diccionario:
	
	
	
	
	\hspace{-2.5cm}{\ttfamily \scriptsize dicc\_zip.py}  \vspace{-0.5cm}
	\begin{lstlisting}
		diccionario = dict{zip(paises,poblaciones)}
		print(diccionario)
	\end{lstlisting}
	{\ttfamily\scriptsize\scriptsize \noindent Out[]: \vspace{-0.3cm}	
		{\color{blue} 	
			\begin{verbatim}
				{'Venezuela': 30, 'china': 1391, 'India': 1364, 'Estados Unidos': 327}
			\end{verbatim}
	}}
	
	
	\section{Ejercicios}     
	
	Crear un diccionario con un objeto {\ttfamily dict(zip(a, b))}.
	
	
	
	% https://jarroba.com/resumibles-hashables-en-python/   “indexando” 
	
	% Como se ve, la función «zip()» nos puede ser de gran ayuda, sobre todo, en aquellos casos en los que necesitamos manipular varios iterables dentro de un solo ciclo o realizar acciones y operaciones con sus elementos al mismo tiempo.
	
	
	
	
	
	%%%%%%%%%%%%%%%%%%%%%%%%%%%%%%%%%%%%  Capitulo           %%%%%%%%%%%%%%%%%%%%%%%%%%%%%%%%%%%%%%%%%
	
	%                     Entrada del Usuario y el Bucle While
	
	%%%%%%%%%%%%%%%%%%%%%%%%%%%%%%%%%%%%             %%%%%%%%%%%%%%%%%%%%%%%%%%%%%%%%%%%%%%%%%
	
	
	\chapter{Entrada del Usuario y el Bucle While}\label{Cap5}\hypertarget{Cap5}{}  % input = Cap5
	
	
	
	
	\begin{wrapfigure}{l}{0.48\textwidth} %  lado derecho, izquierdo, L y l R y r
		\vspace{-25pt}	\begin{center} 
\includegraphics[width=0.48\textwidth]{../Logos_de_Python/Entrada} % angle=45,
		\end{center} \vspace{-30pt}%\caption{Python 3} 
	\end{wrapfigure}
	La mayoría de los programas se escriben para resolver algún  problema  del usuario,  por lo que se necesita recibir información del usuario, para tratar el problema. Por ejemplo, digamos que alguien quiere saber si tienen edad suficiente para votar. Si escribes un programa para responde esta pregunta, necesitas saber la edad del usuario, para dar una respuesta. El programa debe pedirle al usuario que ingrese su edad (o algún dato o información que permita hacer este cálculo), una vez que el programa tiene esta información, puede comparar a la edad  para votar, y
	determinar si el usuario tiene la edad suficiente, luego informar el resultado, que es la salida.
	
	En este capítulo aprenderá cómo codificar la entrada del usuario para que su programa
	puede trabajar con los datos ingresados. Cuando su programa necesite un nombre o una lista de nombres, podrá
	para solicitar al usuario el nombre, o la lista de nombres. Para hacer esto, se usará
	la función {\ttfamily input()}.
	
	También aprenderá cómo mantener los programas ejecutándose el tiempo que los usuarios requieran  para que puedan ingresar tanta información como sea necesaria, y su programa puede usar esa información. Se usará el bucle {\ttfamily while} de Python para mantener los programas en ejecución mientras ciertas condiciones sigan siendo verdaderas.
	
	Podrá escribir programas totalmente interactivos mientras se ejecuten sus programas, con el dominio de la entrada ({\ttfamily input}) del usuario y el control del bucle {\ttfamily while}. 
	
	
	\section{La función input()}
	
	La función {\ttfamily input()} pausa su programa y espera a que el usuario ingrese
	algún texto. Una vez que Python reciba la información  del usuario  (dato que ingresa el usuario, entrada), le asigna  a esa entrada una
	variable que le resulte más cómodo trabajar con ella.
	
	Por ejemplo, el siguiente programa le pide al usuario que ingrese su nombre (puede ser cualquier texto),
	luego imprime el  mensaje, mostrándolo en pantalla al usuario:
	
	\hspace{-2.5cm}{\ttfamily \scriptsize usuario.py}\vspace{-0.5cm}
	\begin{lstlisting}		
		mensaje = input("Escriba su nombre y te lo repetiré: ")	
		print("Hola", mensaje)
	\end{lstlisting}	
	
	La función {\ttfamily input(argumento)} toma el argumento: el mensaje ({\ttfamily prompt}) que queremos
	mostrar al usuario, para que  conozca  qué tipo de información debe ingresar. 
	
	Por ejemplo, cuando Python ejecuta la primera línea, el usuario ve el mensaje en pantalla:
	{\ttfamily Escriba su nombre y te lo repetiré}. Y el programa espera mientras el usuario ingresa su respuesta y presiona {\bf ENTER}, luego continúa. La respuesta del usuario se asigna a la variable {\ttfamily mensaje}, que luego lo imprime en pantalla de vuelta al usuario:
	
	{\ttfamily\scriptsize\scriptsize 
		\hspace{-0.65cm} Out[]:\vspace{-0.3cm}	
		{\color{blue} 	
			\begin{verbatim}
				¡Carlos Dávila!
			\end{verbatim}
	}}
	
	
	
	\vspace{0.2cm}
	\begin{minipage}[t]{0.2\textwidth}
		\centering\raisebox{\dimexpr\topskip-\height}{%
			
\includegraphics[scale=0.20]{../Logos_de_Python/Busquedas_R}}  %[width=\textwidth]
		
	\end{minipage}%\hfill
	\begin{minipage}[t]{0.70\textwidth} 
		{\scriptsize 
			Algunos editores de texto no ejecutan programas que solicitan (prompt) información al usuario. Puedes usar estos editores para escribir programas que soliciten (prompt) entradas, pero necesitará ejecutar estos programas 	desde una terminal. Consulte “Ejecución de programas Python desde una terminal” en la página \ref{terminal}.}
	\end{minipage}
	\vspace{0.2cm}
	
	
	
	\subsection{Escribir claro el mensaje (Prompts)}
	
	Cada vez que utilice la función {\ttfamily input()}, debe incluir un mensaje claro y fácil de seguir. Un 
	mensaje que le diga al usuario exactamente qué tipo de información debe ingresar. 
	De tal forma que la declaración que el usuario ingrese funcione bien. Por ejemplo: 
	
	\hspace{-2.3cm}{\ttfamily \scriptsize saludo.py}\vspace{-0.6cm}
	\begin{lstlisting}		
		nombre = input("Ingrese su nombre: ")
		print(f"\n¡Hola, {nombre}!")	
	\end{lstlisting}	
	
	En esta instrucción se agregó un espacio al final de las indicaciones en el {\ttfamily input()} (después de los dos puntos de la cadena {}"{}Ingrese su nombre{}"{}) para separar el mensaje ("¡Hola") de la respuesta del usuario y dejar claro al usuario dónde introducir su texto. 
	
	{\ttfamily\scriptsize\scriptsize 
		\hspace{-0.65cm}Out[]: \vspace{-0.3cm}	
		{\color{blue} 	
			\begin{verbatim}
				¡Hola Carlos!		
			\end{verbatim}
	}}
	Algunas  veces es necesario   escribir un mensaje (un prompt) que cubre más de una línea. Por ejemplo, es posible que deba indicarle al usuario por qué se solicitan ciertas entradas.  Este tipo de presentación se puede hacer asignando su mensaje a una variable, y pasar esa variable a la función {\ttfamily input()}. Esto permite construir su mensaje en varias líneas, luego se escribe la declaración {\ttfamily input()} clara y elegante. Ejemplo:
	
	\hspace{-2.3cm}{\ttfamily \scriptsize saludo.py}\vspace{-0.5cm}
	\begin{lstlisting}		
		prompt = "Si comparte su nombre, podemos personalizar los mensajes."
		prompt += "\n¿Cuál es tu nombre? "
		nombre = input(prompt)
		print(f"\nHola, {nombre}!")	
	\end{lstlisting}	
	
	Este ejemplo muestra una forma de crear una cadena multilínea. En la primera linea el mensaje
	es asignado  la variable prompt, que será la primera parte del mensaje. En la segunda
	línea, el operador  {\ttfamily += } toma la cadena que fue asignada al mensaje y agrega
	la nueva cadena al final. El mensaje ahora abarca dos líneas, nuevamente de deja un espacio después de la pregunta para mayor claridad del usuario:
	
	{\ttfamily\scriptsize\scriptsize 
		\hspace{-0.65cm}Out[]: \vspace{-0.3cm}	
		{\color{blue} 	
			\begin{verbatim}
				Si comparte su nombre, podemos personalizar los mensajes.
				¿Cuál es tu nombre? Carlos
				
				Hola, Carlos!
			\end{verbatim}
	}}
	
	\subsection{Usando int(), float() para aceptar entradas numéricas}
	
	Cuando usas la función {{\ttfamily input()}, Python interpreta todo lo que el usuario
		ingresa como una cadena.  Consideremos  la  siguiente sesión con el intérprete, en la que se solicita la edad del usuario:
		
		%\hspace{-2.5cm}{\ttfamily \scriptsize curso.py}  \vspace{-0.5cm}
		\begin{lstlisting}
			In[]:
			>>> edad  = input("¿Cuantos años tienes? ")
			¿Cuántos años tiene? 21
			>>> edad
			'21'
		\end{lstlisting}		
		
		El usuario ingresa el número 21, pero cuando le pedimos  a Python el valor de
		edad, devuelve '21', que es la representación en cadena del valor numérico ingresado.
		
		Sabemos que Python interpretó la entrada como una cadena porque el número está entre comillas. Este resultado funciona si lo que quieres  es imprimir la entrada.  Pero si intentas utilizar la entrada como un número, obtendrás un error:
		
		\begin{lstlisting}
			In[]:
			>>> edad = entrada("¿Cuantos años tienes? ")
			¿Cuántos años tiene? 21
			>>> edad >= 18
		\end{lstlisting}
		{\ttfamily\scriptsize\scriptsize \noindent Out[]:\vspace{-0.3cm}	
			{\color{red} 	
				\begin{verbatim}
					Traceback (most recent call last):
					File "C:/Users/Linis/Desktop/IDLE/Tarea9 Input.py", line 11, in <module>
					edad = entrada("¿Cuantos años tienes? ")
					NameError: name 'entrada' is not defined
				\end{verbatim}
		}}
		
		Cuando intentas usar la entrada para hacer una comparación numérica, Python
		produce un error porque no puede comparar una cadena con un número entero: la cadena
		'21' que se asigna a la edad no se puede comparar con el valor numérico 18.
		
		Podemos resolver este problema usando la función {\ttfamily int()}, que convierte la cadena de entrada en un valor numérico. Esto permite que el cálculo se ejecute correctamente:
		\begin{lstlisting}
			>>> edad = input("¿Cuantos años tienes? ")
			¿Cuántos años tiene? 21
			>>> edad = int(edad)
			>>> edad >= 18
			Verdadero
		\end{lstlisting}
		
		
		
		En este ejemplo, cuando ingresamos 21 en el {\ttfamily prompt}, Python interpreta el número como una cadena, pero función {\ttfamily int()} o  {\ttfamily float()}  convierte el valor a una representación numérica. Ahora Python puede ejecutar el cálculo, y compara la edad (pues 21 ahora representa el valor numérico) con 18, y prueba  si la edad es mayor o igual a 18. El valor de esta prueba es booleano, {\ttfamily True} o {\ttfamily False}.
		
		Una aplicación de  las funciones {\ttfamily int()} o {\ttfamily float)} en un programa más real:
		Considere un programa que determina si las personas son lo suficientemente altas como para subirse a un pico en los Andes:
		
		\hspace{-3.0cm}{\ttfamily \scriptsize pico\_andino.py}\vspace{-0.5cm}
		\begin{lstlisting}		
			altura = input("¿Qué altura tienes, en metros? ")
			altura = float(altura)
			if altura >= 1.60:
			print("\n¡Eres lo suficientemente alto para subir!")
			esle:
			print("\nDisculpe, podrás subir cuando seas  mayor.")
		\end{lstlisting}	
		
		{\ttfamily\scriptsize\scriptsize 
			\hspace{-0.65cm} \vspace{-0.3cm}	
			{\color{blue} 	
				\begin{verbatim}
					
				\end{verbatim}
		}}
		
		
		El programa puede comparar el valor de la {\ttfamily altura} ingresado con 1.60, porque {\ttfamily float(altura)}
		lo convierte en un flotante antes de realizar la comparación. Si el número ingresado es mayor o igual a 1.60, le decimos al usuario que  ¡Eres lo suficientemente alto para montar!, en caso contrario se le informa: Disculpe, podrás subir cuando seas  mayor.\\
		
		\vspace{0.2cm}
		\begin{minipage}[t]{0.2\textwidth}
			\centering\raisebox{\dimexpr\topskip-\height}{%
				
\includegraphics[scale=0.20]{../Logos_de_Python/Sugerencia}}  %[width=\textwidth]
			
		\end{minipage}%\hfill
		\begin{minipage}[t]{0.70\textwidth} 
			{\scriptsize Se recomienda  que  utilice estas funciones cuando trabajes con entradas numéricas, para  
				asegurar  la representación  numérica, y  permita realizar los cálculos.}
		\end{minipage}
		\vspace{0.2cm}
		
		
		
		
		
		
		\subsection{El operador módulo {{ }\ttfamily \%} }
		
		Una herramienta útil para trabajar con información numérica es el operador de módulo ({\ttfamily \%}), que divide un número entre otro y devuelve el resto de la división, por ejemplo:
		%\hspace{-2.5cm}{\ttfamily \scriptsize curso.py}  \vspace{-0.5cm}
		\begin{lstlisting}		
			>>> 7 % 3
			1
			>>> 5 % 3
			2
			>>> 6 % 3
			0
		\end{lstlisting}	
		
		
		El operador de módulo  devuelve el resto de dividir un número  en otro. Cuando un número es divisible por otro (división exacta) el resto es 0, por lo que el operador de módulo siempre devuelve cero. Puede utilizar este hecho para determinar si un número es par o impar:
		
		\hspace{-3.1cm}{\ttfamily \scriptsize par\_impar.py.py}\vspace{-0.5cm}
		\begin{lstlisting}
			número = input("Ingresa un número y te diré si es par o impar: ")
			número = int(número)
			if número %2 == 0:
			print(f"\nEl número {número} es par.")
			else:
			print(f"\nEl número {número} es impar.")
		\end{lstlisting} 
		
		Los números pares siempre son divisibles por dos, por lo que si el módulo 2 de un número
		es cero (número \% 2 == 0) el número es par,  de lo contrario, es impar.
		
		{\ttfamily\scriptsize\scriptsize \noindent Out[]:\vspace{-0.3cm}	
			{\color{blue} 	
				\begin{verbatim}
					Ingresa un número y te diré si es par o impar: 42
					El número 42 es par.			
				\end{verbatim}
		}}
		
		
		
		\subsection{Introducción al  bucle While}
		
		El bucle {\ttfamily for} toma una colección de elementos y  por cada elemento (item) de la colección ejecuta un bloque de código. Por el contrario, el bucle  {\ttfamily while} se ejecuta mientras una  cierta condición sea verdadera (true). 
		
		\subsubsection{El bucle while en acción}
		
		El bucle {\ttfamily while} se puede utilizar para contar una serie de números, 
		como en el siguiente ejemplo, el  {\ttfamily while} cuenta de 1 hasta el 5:
		
		\hspace{-2.5cm}{\ttfamily \scriptsize contando.py}\vspace{-0.55cm}
		\begin{lstlisting}		
			
		\end{lstlisting}	
		%lista.append(número_actual)
		%print(número_actual, lista)
		
		En la primera línea, comenzamos asignándole a la variable  {\ttfamily número\_actual} el valor de 1. Luego, el bucle {\ttfamily while} se configura para ejecutarse mientras el valor de {\ttfamily número\_actual}  sea menor o igual a 5. El código dentro del bucle imprime el valor de {\ttfamily número\_actual}, y suma 1 a ese valor {\ttfamily número\_actual} con el operador {\ttfamily +=} (El operador {\ttfamily +=} es una abreviatura de {\ttfamily x = x + 1}). Python repite este ciclo siempre que la condición {\ttfamily número\_actual <= 5} sea verdadera. En el primer ciclo la variable  {\ttfamily número\_actual} es igual a  1, que es menor que 5, Python imprime 1 y  suma 1 a {\ttfamily número\_actual}, ahora {\ttfamily número\_actual} tiene valor igual a 2. Y se inicia nuevamente el proceso;  como 2 es menor que 5, Python imprime 2 y nuevamente  suma 1  a {\ttfamily número\_actual},  ahora {\ttfamily número\_actual}  es 3. Se inicia nuevamente el proceso y así sucesivamente, hasta  que el valor de la variable {\ttfamily número\_actual} es mayor que 5, en este momento el bucle deja de ejecutarse y el programa finaliza:
		
		{\ttfamily\scriptsize\scriptsize \noindent Out[]:\vspace{-0.3cm}	
			{\color{blue} 	
				\begin{verbatim}	
					1
					2
					3
					4
					5	
				\end{verbatim}
		}} 
		
		En los programas que se  utilizan todos los días probablemente tengan un bucle {\ttfamily while}. Por ejemplo, un juego necesita un bucle {\ttfamily while} para seguir ejecutándose mientras el usuario quiera seguir jugando, por lo que puede dejar de ejecutarse tan pronto como le pidas que salga. Los programas no serían agradable  si dejaran de ejecutarse antes de que se lo indiquemos o si siguieran ejecutándose incluso después de que quisiéramos salir,  es un buen motivo del por qué los bucles {\ttfamily while} son tan importante y útiles.
		
		\subsubsection{Permitir que el usuario elija cuándo salir}
		
		Podemos hacer que un programa se mantenga en ejecución  hasta que el usuario decida salir; esto se hace insertando la mayor parte del programa dentro de un bucle {\ttfamily while},  y definiendo  un valor para salir, de tal forma que al ejecutar el programa, este permanecerá en ejecución hasta que el usuario ingrese  el valor de salida. Ilustremos con un ejemplo este problema: 
		
		\hspace{-2.2cm}{\ttfamily \scriptsize activo.py} \vspace{-0.55cm}
		\begin{lstlisting}		
			prompt = "\nIngrese 'salir' para finalizar el programa."
			prompt +="\nIngrese cualquier texto y te lo repetiré: "
			
			mensaje = ""
			while mensaje != 'salir':
			mensaje = input(prompt)
			print(mensaje)
		\end{lstlisting}	
		
		Primero definimos un {\ttfamily prompt} que le da al usuario dos opciones: Ingresar el valor para salir del programa (en este caso, es "salir") o ingresar un texto cualquiera. Luego configuramos un mensaje variable para realizar un seguimiento a cualquier valor que ingrese el usuario. 
		
		El mensaje lo definimos como una cadena vacía (\ttfamily {}"{}"{}}), por lo que Python ya tiene algo para comparar en la primera vez que llega a la línea del {\ttfamily while}.  La primera vez que se ejecuta el programa y Python llega a la declaración del {\ttfamily while},  necesita comparar el valor del mensaje con "salir",  { }por ahora es la cadena vacía definida, porque el usuario  aún no ha ingresado ninguna entrada.  Si Python no tiene algo para hacer la comparación, no podrá continuar ejecutando el programa y generara un error. Para evitar este problema, debemos asegurarnos de darle un valor inicial al mensaje, aunque es sólo un cadena vacía, tendrá sentido para Python y le permitirá realizar la comparación, iniciando el proceso del ciclo {\ttfamily while}. El bucle {\ttfamily while} se ejecutará mientras el valor del mensaje sea diferente de {\ttfamily "{}salir{}"{}}.
	
	La primera vez que recorre el bucle, el mensaje es sólo una cadena vacía, por lo que Python
	entra en el bucle. 
	
	Python muestra el mensaje {\ttfamily mensaje = input(prompt)}, y espera a que el usuario ingrese su entrada. Todo lo que 
	ingresa se asigna al mensaje y se imprime; luego, Python reeeevalúa la condición en el {\ttfamily while}. 
	Mientras el usuario no  ingrese la palabra "salir", el mensaje se muestra nuevamente y Python espera más entradas. 
	Cuando el usuario finalmente ingresa 'salir', Python deja de ejecutar el bucle {\ttfamily while} y el programa finaliza:
	
	
	{\ttfamily\scriptsize\scriptsize \noindent Out[]: \vspace{-0.3cm}	
		{\color{blue} 	
			\begin{verbatim}
				
				Ingrese 'salir' para finalizar el programa.
				Ingrese cualquier texto y te lo repetiré: Carlos
				Carlos
				
				Ingrese 'salir' para finalizar el programa.
				Ingrese cuarquie texto y te lo repetiré: Carlos
				Carlos
				
				Ingrese 'salir' para finalizar el programa.
				Ingrese cuarquie texto y te lo repetiré: salir
				salir	
			\end{verbatim}
	}}
	
	
	Este programa funciona bien, pero imprime la palabra "salir",  como si se tratara de un mensaje programado, para soluciona este problema, usamos una comparación con {\ttfamily if}, y resuelto:
	
	\hspace{-2.2cm}{\ttfamily \scriptsize activo.py} \vspace{-0.55cm}
	\begin{lstlisting}		
		prompt = "\nIngrese 'salir' para finalizar el programa."
		prompt +="\nIngrese cuarquie texto y te lo repetiré: "
		
		mensaje = ""
		while mensaje != 'salir':
		mensaje = input(prompt)
		if mensaje != 'salir':
		print(mensaje)
	\end{lstlisting}	
	
	
	Ahora el programa hace un chequeo rápido antes de mostrar el mensaje y sólo imprime el mensaje si no coincide con el valor "salir":
	
	{\ttfamily\scriptsize\scriptsize \noindent Out[]:\vspace{-0.3cm}	
		{\color{blue} 	
			\begin{verbatim}
				
				Ingrese 'salir' para finalizar el programa.
				Ingrese cuarquie texto y te lo repetiré: Carlos
				Carlos
				
				Ingrese 'salir' para finalizar el programa.
				Ingrese cuarquie texto y te lo repetiré: Prueba
				Prueba
				
				Ingrese 'salir' para finalizar el programa.
				Ingrese cuarquie texto y te lo repetiré: salir			
			\end{verbatim}
	}}
	
	
	\subsubsection{Usando una bandera como señal indicadora}
	
	En el ejemplo anterior, hicimos que el programa realizara ciertas tareas mientras una condición sea verdadera. Pero ¿qué sucede en programas más complicados en los que varios eventos diferentes pueden provocar la salida del programa?  Por ejemplo, en un juego, varios eventos diferentes pueden finalizar el juego; cuando el jugador se queda sin barcos, cuando se le acaba el tiempo o cuando las ciudades que se suponía que debían proteger son destruidas, el juego debe terminar si ocurre alguno de estos eventos. El programa se detiene si al menos ocurre un evento de los que forma la declaración condicional,  probar todas estas condiciones un único bule  while  puede  ser complicado y difícil.  
	
	Para escribir un programa que permanezca en ejecución  mientras se cumplan varias condiciones, se puede definirse una variable que determine si todas las condiciones están activas o no. Esta variable es llamada  bandera (flag), el conjunto de evento que forma la declaración condicional,  actúa como una señal indicadora para el programa.  El programa permanece  en ejecuten mientras el indicador sea  Verdadero, y será interrumpido cuando el  indicador sea Falso. Como consecuencia, la declaración  {\ttfamily while} sólo necesita verificar una condición: si el indicador es Verdadero. Luego, las demás condiciones de prueba (para verificar si un evento en el indicador ha ocurrido y lo configura como Falso) se pueden organizar claramente en el otra parte  del programa.
	
	Agreguemos una variable indicadora al programa anterior, que llamamos  {\ttfamily activo}; este indicador   permite  monitorear el programa,  para mantenerlo en ejecución o proceder a finalizarlo:
	
	\begin{lstlisting}
		prompt = "\nIngrese 'salir' para finalizar el programa."
		prompt +="\nIngrese cualquier texto y te lo repetiré: "
		
		active = True
		while active:
		mensaje = input(prompt)
		if mensaje == 'salir':
		active = False
		else:
		print(mensaje)            
	\end{lstlisting}	
	
	
	A la variable  indicadora {\ttfamily active} se le puede dar cualquier nombre de variable, y funcionará igual. Se configura como  {\ttfamily True},  para iniciar el programa en un estado activo. 
	
	Este fragmento de código simplifica la declaración {\ttfamily while},  porque  en sí no se realiza comparación alguna en la declaración del {\ttfamily  while}; por lógica se toma en cuenta las demás partes del programa. Mientras la variable {\ttfamily active}  siga siendo Verdadera, el ciclo continuará ejecutándose.
	
	
	La declaración {\ttfamily if} dentro del bucle {\ttfamily while}, verifica el valor del mensaje una vez que el usuario ingresa su entrada. Si el usuario ingresa "salir",  se configura la variable {\ttfamily active} en {\ttfamily False} y el ciclo {\ttfamily while} se detiene. Si el usuario ingresa algún texto distinto a "salir",  { }se imprime la entrada como mensaje y  el ciclo continua.
	
	Este programa tiene el mismo resultado que el ejemplo anterior donde colocamos la prueba condicional directamente en la declaración {\ttfamily while}. Pero ahora con una variable indicadora  que  verifica si el programa se encuentra en estado {\ttfamily activo}. Además será  fácil agregar otros  test (declaraciones como {\ttfamily  elif}) de eventos que convierta en {\ttfamily False} la variable {\ttfamily activo}.
	
	
	Esta idea es usada en programas complejos, como en juegos, en los que pueden haber varios eventos y cada uno de ellos hace que el programa deja de ejecutarse. Cuando cualquiera de estos eventos hace que el indicador  {\ttfamily active} cambie de {\ttfamily True} a  {\ttfamily False}, el bucle principal del juego se cerrará y se  mostrara un mensaje al jugador que  puede ser darle la opción de empezar un nuevo juego.
	
	
	\subsubsection{Usar un {\ttfamily break} para salir de un bucle} 
	
	
	Para salir de un bucle {\ttfamily while} inmediatamente sin ejecutar códigos restante en
	en el bucle, y sin considerar resultados de cualquier prueba o test condicional, se puede usar la  instrucción {\ttfamily break}; esta declaración  dirige el flujo de su programa; puedes usarlo
	para controlar qué líneas de código se ejecutan y cuáles no, por lo que el programa
	solo ejecuta el código que usted desea, cuando lo desea.
	
	Por ejemplo, considere un programa que pregunta al usuario sobre los lugares que ha visitado.
	Podemos detener el ciclo {\ttfamily while} en este programa llamando a {\ttfamily break} tan pronto como
	cuando el usuario ingresa el valor 'salir':
	
	\hspace{-2.6cm}{\ttfamily \scriptsize ciudades.py}\vspace{-0.5cm}
	\begin{lstlisting}		
		prompt = "\nIngrese el nombre de una ciudad que haya visitado:"
		prompt += "\n(Ingrese 'salir' cuando haya terminado.) "
		
		while True:
		ciudad = input(prompt)
		if ciudad == 'salir':
		break
		else:
		print(f"¡Me encantaría ir a {ciudad.title()}!")
	\end{lstlisting}	
	
	El bucle que comienza con {\ttfamily while True} estará  siempre en ejecución  a menos que alcance la
	declaración {\ttfamily break}. En este programa el bucle continúa pidiendo al usuario que ingrese
	el nombre de una ciudad que haya visitado, hasta que ingrese "salir". Cuando ingrese 'salir', se ejecuta la instrucción {\ttfamily break},haciendo que Python salga del bucle:\vspace{0.08cm}
	
	{\ttfamily\scriptsize\scriptsize\noindent Out[]:\vspace{-0.3cm}	
		{\color{blue} 	
			\begin{verbatim}
				Ingrese el nombre de una ciudad que haya visitado:
				(Ingrese 'salir' cuando haya terminado.) Mérida
				¡Me encantaría ir a Mérida!
				
				Ingrese el nombre de una ciudad que haya visitado:
				(Ingrese 'salir' cuando haya terminado.) Caracas
				¡Me encantaría ir a Caracas!
				
				Ingrese el nombre de una ciudad que haya visitado:
				(Ingrese 'salir' cuando haya terminado.) salir
			\end{verbatim}
	}}
	
	
	
	\vspace{0.2cm}
	\begin{minipage}[t]{0.2\textwidth}
		\centering\raisebox{\dimexpr\topskip-\height}{%
			
\includegraphics[scale=0.20]{../Logos_de_Python/Busquedas_R}}  %[width=\textwidth]
		
	\end{minipage}%\hfill
	\begin{minipage}[t]{0.70\textwidth} 
		{\scriptsize Puede utilizar la instrucción {\ttfamily break} en cualquiera de los bucles de Python. Por ejemplo, podrías usarla
			para salir de un bucle {\ttfamily for} que está trabajando en una lista o un diccionario.
		}
	\end{minipage}
	\vspace{0.2cm}
	
	
	
	\subsubsection{Usando {\ttfamily continue} en un bucle}
	
	En lugar de salir de un bucle por completo sin ejecutar el resto del código, puede utilizar la instrucción {\ttfamily continue} para volver al principio del bucle, según el resultado del test condicional. Por ejemplo, considere un bucle que cuenta del 1 al 10 pero imprime solo los números impares en ese rango:
	
	\hspace{-2.5cm}{\ttfamily \scriptsize contador.py} \vspace{-0.55cm}
	\begin{lstlisting}		
		contador = 0
		while  contador < 10:
		contador += 1
		if contador % 2 == 0:
		continue    
		print(contador)	
	\end{lstlisting}	
	
	
	Primero, configuramos la variable {\ttfamily contador} igual a cero 0; como es menor que 10, Python ingresa al ciclo {\ttfamily while}, una vez dentro del bucle, incrementamos {\ttfamily contador} en una unidad ({\ttfamily += 1}).  Y luego la instrucción {\ttfamily if}  verifica la  operación {\ttfamily contador} módulo de 2 ({\ttfamily contador \% 2}). Si el módulo es 0 (significa {\ttfamily contador} es divisible por 2), la instrucción {\ttfamily continue} le dice a Python que ignore el resto del ciclo y regrese al principio. Si el  {\ttfamily contador} no es divisible por 2, se ejecuta el resto del ciclo y Python imprime el  {\ttfamily contador}:\vspace{0.09cm}
	
	{\ttfamily\scriptsize\scriptsize \noindent Out[]:\vspace{-0.3cm}	
		{\color{blue} 	
			\begin{verbatim}
				1
				3
				5
				7
				9		
			\end{verbatim}
	}}
	
	
	\subsection{Evitar bucles infinitos}
	
	Todo  bucle  {\ttfamily while} necesita una forma para detener su ejecución para que no continúe ejecutándose por siempre. Por ejemplo, este ciclo de conteo cuenta  del 1 al 5:
	
	\hspace{-2.5cm}{\ttfamily \scriptsize cuenta5.py} \vspace{-0.55cm}
	\begin{lstlisting}		
		x = 1
		while x <= 5:
		print(x)
		x += 1
	\end{lstlisting}	
	
	
	Sin embargo, si accidentalmente omites la línea {\ttfamily x += 1}, el bucle se ejecutará para siempre:
	
	\begin{lstlisting}	
		# ¡Este bucle dura para siempre!
		x = 1
		while x <= 5:
		print(x)
	\end{lstlisting}
	
	Ahora el valor de {\ttfamily x} comenzará en 1 pero nunca cambiará. Como resultado, el condicional
	test {\ttfamily x <= 5} siempre se evaluará como {\ttfamily True} y se ejecutará el bucle {\ttfamily while}
	para siempre, imprimiendo una serie infinita de 1, como la siguiente:
	{\ttfamily\scriptsize\scriptsize\noindent  Out[]:\vspace{-0.3cm}	
		{\color{blue} 	
			\begin{verbatim}
				1
				1
				1
				1
				--recorte--	
			\end{verbatim}
	}}
	
	
	Todo programador escribe accidentalmente un bucle {\ttfamily while} infinito de vez en cuando, especialmente cuando los bucles de un programa tienen condiciones de salida sutiles. En un caso de bucle infinito, presione {\ttfamily {\bf CTRL-C}} o simplemente cierre el ventana de terminal que muestra la salida de su programa.
	
	\vspace{0.2cm}
	\begin{minipage}[t]{0.2\textwidth}
		\centering\raisebox{\dimexpr\topskip-\height}{%
			
\includegraphics[scale=0.20]{../Logos_de_Python/Busquedas_R}}  %[width=\textwidth]
		
	\end{minipage}%\hfill
	\begin{minipage}[t]{0.70\textwidth} 
		{\scriptsize En VS Code, como muchos editores, muestra el resultado en una ventana de terminal integrada. Al
			cancelar un bucle infinito, asegúrese de hacer clic en el área de salida del editor antes presionando CTRL-C.}
	\end{minipage}
	\vspace{0.2cm}
	
	Para evitar escribir bucles infinitos, pruebe cada bucle {\ttfamily while} por separado y asegúrese que
	el bucle se detiene cuando se le indica. Si desea que su programa finalice cuando el usuario
	ingresa un valor específico de entrada, ejecute el programa e ingresa ese valor. Si el programa no finaliza, revise la forma en que su programa maneja el valor Booleano (condicional) que finaliza el bucle. Asegúrese de que al menos una parte del programa active la condición {\ttfamily Falsa} del  bucle, o  que alcance una declaración de interrupción.
	
	
	
	
	
	
	\subsection{Usando {\ttfamily while} con listas y diccionarios}
	
	Hasta ahora, hemos trabajado solo con una información del usuario, recibimos la entrada del usuario y luego imprimimos la entrada o una respuesta. Sin embargo, cuando hay varios usuarios, y se debe hacer seguimiento a los usuarios  y sus  datos de entrada, esta tarea se puede hacer con listas o diccionarios en nuestros bucles {\ttfamily while}.
	
	
	Un bucle {\ttfamily for} es eficaz para recorrer una lista, pero no debe modificar una lista dentro de un bucle {\ttfamily for}  porque Python tendrá problemas para realizar un seguimiento de los elementos de la lista. Para modificar una lista a medida que trabaja en ella, utilice un bucle {\ttfamily while}. El uso de bucles {\ttfamily while} con listas y diccionarios le permite recopilar, almacenar y organizar una gran cantidad de entradas para examinarlas y generar reportes.
	
	\subsubsection{Mover elementos de una lista a otra}
	
	Considere una lista de usuarios recién registrados pero no verificados de un sitio web. Después de verificar a estos usuarios, debemos  moverlos a una lista de usuarios verificados; una forma de hacer esta tarea, sería usando un bucle {\ttfamily while} para extraer usuarios de la lista de usuarios no verificados,  y  agregarlos a otra lista de usuarios verificados, a medida que son verificados. Ilustremos la idea con un ejemplo:
	
	
	\hspace{-2.5cm}{\ttfamily \scriptsize usuarios.py}  \vspace{-0.55cm}
	\begin{lstlisting}		
		# Comience con los usuarios que necesitan ser verificados,
		# y una lista vacía para guardar usuarios verificados.
		
		no_verificados = ['carlos', 'Julian', 'oswaldo']
		verificados = []
		# Verificar cada usuario hasta que no queden más usuarios sin verificar.
		# Mueva cada usuario verificado a la lista de usuarios verifcado.
		
		while no_verificados:
		usuario_actual = no_verificados.pop()
		
		print(f"Verificando usuario: {usuario_actual.title()}")    
		verificados.append(usuario_actual)
		
		# Mostrar todos los usuarios confirmados.
		print("\nLos siguientes usuarios han sido verificados:")
		for  verificado in  verificados:
		print(verificado.title())
		
	\end{lstlisting}	
	
	
	Comenzamos con una lista de usuarios no verificados (carlos, Julian, oswaldo) y una lista vacía para guardar  a los usuarios verificados. El bucle {\ttfamily while} se ejecuta siempre que la lista usuarios no verificados  no esté vacía. Dentro de este bucle, el método {\ttfamily pop()} elimina el último usuario no verificado de la lista.  Debido a que oswaldo es el último en la lista de usuarios no verificados, su nombre sera el primero en ser eliminado, y lo agregamos a la lista de usuarios verificados.  El siguiente es Julian, y luego Carlos.
	Simulamos una  verificación  de cada usuario imprimiendo un mensaje de verificación y luego lo agregamos  a la lista de usuarios verificados. A medida que la lista de usuarios no verificados se reduce, la lista de usuarios verificados crece. Cuando la lista de usuarios no verificada queda vacía, el bucle se detiene y la lista de usuarios verificados se imprime:\vspace{0.08cm}
	
	{\ttfamily\scriptsize\scriptsize \noindent Out[]: \vspace{-0.3cm}	
		{\color{blue} 	
			\begin{verbatim}
				Verificando usuario: Oswaldo
				Verificando usuario: Julian
				Verificando usuario: Carlos
				
				Los siguientes usuarios han sido verificados:
				Oswaldo
				Julian
				Carlos		
			\end{verbatim}
	}}
	
	
	
	\subsubsection{Eliminar valores específicos de una lista}
	
	En la sección \ref{remueve}, usamos el método {\ttfamily remove()} para eliminar un valor específico de una lista. El método {\ttfamily remove()} funciona bien  para  eliminar un valor en la lista, pero si el valor se repite  más de una vez en la lista,  sólo se elimina el primero que aparece. Para eliminar todas las instancias en la lista debe hacerlo mediante un bucle.
	
	Supongamos que tenemos  una lista de mascotas con el valor {}"{}gato{}"{} repetido varias veces. Para eliminar todas las instancias con  ese valor, se puede emplear un bucle {\ttfamily while} hasta que {}"{}gato{}"{}  quede eliminado de la lista, como se muestra en el siguiente código:
	
	
	\hspace{-2.5cm}{\ttfamily \scriptsize mascotas.py}  \vspace{-0.5cm}
	\begin{lstlisting}		
		mascotas = ['perro', 'gato', 'perro', 'loro', 'gato', 'conejo', 'gato']
		print(mascotas)
		while 'gato' in mascotas:
		mascotas.remove('gato')
		print(mascotas)
	\end{lstlisting}	
	
	Comenzamos con una lista que contiene múltiples instancias de 'gato'. Después de imprimir
	la lista, Python ingresa al bucle {\ttfamily while} porque encuentra el valor 'gato' en la lista
	al menos una vez. Una vez dentro del bucle, Python elimina la primera instancia de 'gato',
	regresa a la línea {\ttfamily while} y luego vuelve a ingresar al bucle cuando descubre que 'gato' es
	todavía en la lista. Elimina cada instancia de 'gato' hasta que el valor ya no esté en
	la lista, momento en que Python sale del bucle e imprime la lista nuevamente: \vspace{0.08cm}
	
	{\ttfamily\scriptsize\scriptsize \noindent Out[]: \vspace{-0.3cm}	
		{\color{blue} 		
			\begin{verbatim}
				['perro', 'gato', 'perro', 'loro', 'gato', 'conejo', 'gato']
				['perro', 'perro', 'loro', 'conejo'] 			
			\end{verbatim}
	}}
	
	
	\subsubsection{Llenar un diccionario con entradas del usuario}
	
	Puede solicitar toda la información que necesite en cada paso con el bucle 
	{\ttfamily while}. Hagamos un programa de sondeo (encuestas) en el que cada  paso por el bucle
	solicita el nombre del participante y su respuesta. Almacenaremos los datos recolectados en un diccionario, porque necesitamos  conectar cada respuesta con su correspondiente usuario:
	
	
	\hspace{-2.3cm}{\ttfamily \scriptsize sondeo.py}  \vspace{-0.5cm}
	\begin{lstlisting}		
		respuestas = {}  # Diccionario vacío
		sondeo_activo = True	
		while sondeo_activo:
		# Solicite el nombre de la persona y su respuesta.
		nombre = input("\n¿Cuál es su nombre? ")
		respuesta = input("¿Qué montaña te gustaría escalar algún día? ")
		# Almacenar la respuesta en el diccionario.
		respuestas[nombre] = respuesta
		# Descubra si alguien más va a realizar la encuesta.
		repetir = input("¿Le gustaría dejar que otra persona responda? (sí/no) ")
		if repetir == 'no':
		sondeo_activo = False
		# El sondeo está completo. Muestra los resultados.
		print("\n--- Resultados de la encuesta ---\n")
		for nombre, respuesta in respuestas.items():
		print(f"{nombre.title()} le gustaría subir el pico {respuesta.title()}.")
	\end{lstlisting}	
	
	%  Las variables “bandera” nos sirven para indicar y almacenar que algún evento o condición se cumplió.
	
	El programa primero define un diccionario vacío (respuestas) y establece una bandera ({\ttfamily sondeo\_activo}) para indicar que el sondeo está activo. Mientras {\ttfamily sondeo\_activo} sea True, Python ejecutará el código en el bucle {\ttfamily while}.  Dentro del bucle, se le solicita al usuario que ingrese su nombre y la montaña que le gustaría escalar. Esa información se almacena en el diccionario de respuestas, y se le pregunta al usuario si desea mantener o no la encuesta en ejecución. Si ingrese sí, el programa ingresa nuevamente al ciclo {\ttfamily while}. Si ingresa no, el indicador {\ttfamily sondeo\_activo} su valor a {\ttfamily False}, y el bucle {\ttfamily while} deja de ejecutarse,  el bloque de código finaliza y muestra los resultados de la encuesta.\vspace{0.08cm}
	
	{\ttfamily\scriptsize\scriptsize \noindent Out[]: \vspace{-0.3cm}	
		{\color{blue} 	 	
			\begin{verbatim}
				¿Cuál es su nombre? Oswaldo
				¿Qué montaña te gustaría escalar algún día? BolivAr
				¿Le gustaría dejar que otra persona responda? (sí/no) si
				
				¿Cuál es su nombre? CArlos
				¿Qué montaña te gustaría escalar algún día? EspEjo
				¿Le gustaría dejar que otra persona responda? (sí/no) si
				
				¿Cuál es su nombre? José
				¿Qué montaña te gustaría escalar algún día? Indio
				¿Le gustaría dejar que otra persona responda? (sí/no) no
				
				--- Resultados de la encuesta ---
				
				Oswaldo le gustaría subir el pico Bolivar.
				CArlos le gustaría subir el pico Espejo.
				José le gustaría subir el pico Indio. 			
			\end{verbatim}
	}}
	
	
	
	
	\section{Resumen}
	
	En este capítulo, aprendió cómo usar la función  {\ttfamily input()} para permitir a los usuarios ingresar
	su información en sus programas. Aprendiste a trabajar  entrada de texto y numérica, y cómo usar bucles  {\ttfamily while}  para hacer sus programas se mantengan en ejecución en tiempo que sus usuarios quieran. 
	
	Se estudió  varias formas para controlar el  flujo de un bucle  {\ttfamily while}, mediante un indicador activo, usando la declaración  {\ttfamily break}, y  la declaración {\ttfamily continue}. Aprendiste a usar un bucle  {\ttfamily while} para mover  elementos de una lista a otra y cómo eliminar todas las instancias de un valor en una lista. También aprendió cómo se pueden usar los bucles  {\ttfamily while} con diccionarios.
	
	
	
	\section{Ejercicios}
	
	
	Sin revisar....
	
	\begin{enumerate}[3.5.1]
		
		\item  Coche de alquiler: Escriba un programa que pregunte al usuario qué tipo de coche gustaría para alquilar.
		Imprime un mensaje sobre ese auto, como "Déjame ver si puedo encontrarte un Toyota”.
		
		\item Asientos en restaurante: Escriba un programa que le pregunte al usuario cuántas personas
		están en su grupo de cena. Si la respuesta es más de ocho, imprima un mensaje que diga:
		Tendrán que esperar por una mesa. En caso contrario, informe que su mesa está lista.
		
		\item  Múltiplos de diez: Pida al usuario un número y luego informe si el número es múltiplo de 10 o no.
		
		\item Ingredientes para pizza: Escriba un bucle que solicite al usuario que introduzca una serie de
		ingredientes para pizza hasta que ingresen un valor de "salir". A medida que ingresan a cada ingrediente, imprima
		un mensaje que diga: agregarás ese aderezo a su pizza.
		
		\item Entradas de cine: Una sala de cine cobra diferentes precios de entrada dependiendo de
		la edad de una persona. Si es menor de 3 años, la entrada es gratuita; si tienen entre las 3 y  12 años, el billete cuesta \$ 10;  si son mayores de 12 años, el billete es  {\scriptsize\$} 15. Escribe un bucle en el que preguntes a los usuarios su edad y luego les digas el precio de su entrada al cine.
		
		\item Tres salidas: Escriba diferentes versiones del ejercicio 7.4 o 7.5 en cada uno de los siguientes casos, al menos una vez:
		
		Utilice una prueba condicional en la sentencia {\ttfamily{while}} para detener el ciclo.
		
		Utilice una variable activa para controlar cuánto tiempo se ejecuta el bucle.
		
		Utilice una instrucción {\ttfamily{break}} para salir del bucle cuando el usuario introduzca un valor de {\ttfamily{'salir'}}  ('quit').
		
		\item  Infinito: Escribe un bucle que nunca termine y ejecútalo. (Para finalizar el bucle, presione
		{\ttfamily{CTRL-C}} o cierre la ventana que muestra el resultado).
		
	\end{enumerate}
	
	
	
	
	\chapter{Estructuras de control de flujo: if-elif-else}\label{Cap6}\hypertarget{Cap6}{}
	
	%	Estructurar código Python, declaraciones  if}


\begin{wrapfigure}{l}{0.45\textwidth} 
	\vspace{-25pt}	\begin{center} 
\includegraphics[width=0.48\textwidth]{../Logos_de_Python/elif01}  % if_elif
	\end{center}\vspace{-40pt} %\caption{Python 3} 
\end{wrapfigure}En los Capítulos anteriores se han visto muchos ejemplos con datos simples y estructuras de datos con algunos de  sus métodos, pero no ha hecho mucho con ellos. La mayoría de los ejemplos de código son cortos, utilizados en el intérprete para ilustrar su aplicación.  Ahora trabajaremos la idea general de cómo estructurar el código Python, no solo con los datos, sino con algunas aplicaciones sencillas.

Muchos lenguajes de programación  utilizan caracteres como llaves (\{ \}) o palabras clave, 
para marcar secciones de código al inicio y al final. En esos lenguajes, es una buena práctica
utilizar sangrías consistentes para hacer que su programa sea más legibles tanto para el programador como para futuros lectores. 
Incluso hay herramientas para hacer que su código se alinee bien.\\


Cuando Guido van Rossum estaba diseñando el lenguaje que se convirtió en Python, decidió que la sangría en sí era suficiente para definir la estructura de un programa, y evitó escribir todos esos paréntesis y llaves. Esto es uno de los primeros aspectos que notan los recién llegados al lenguaje Python, y puede parecer extraño para quienes tienen  experiencia con otros lenguajes. Resulta que después de escribir Python durante un cierto tiempo, se hace natural y dejas de notarlo. Incluso te acostumbras a hacer más mientras escribes menos.\\
%  Recuerde, print() es la función incorporada de Python para imprimir cosas, normalmente a tu pantalla.


Hasta ahora hemos hablado casi exclusivamente de estructuras de datos, pero ya estamos preparados para dar los primeros pasos con las estructuras de código que entrelazan los datos en los programas. Generalmente la programación informática implica examinar un conjunto de condiciones y decidir qué acciones tomar en base a esas condiciones. \\


Los condicionales nos permiten comprobar si se cumple o no cierta condición y hacer que nuestro programa se comporte de una forma u otra; que ejecute un fragmento de código u otro, dependiendo de esa condición. Algunas veces esta condición es llamada test condicional, en este caso podemos decir, que un test condicional es una comparación de dos o más situaciones que genera un resultado booleano verdadero o falso ({\ttfamily True}  o {\ttfamily False}).   Aquí es donde cobran su importancia el tipo booleano, los operadores lógicos y relacionales que se estudian en la sección \ref{logico}. 



\vspace{0.2cm}
\begin{minipage}[t]{0.2\textwidth}
	\centering\raisebox{\dimexpr\topskip-\height}{%
		
\includegraphics[scale=0.20]{../Logos_de_Python/Advertencia}}  %[width=\textwidth]	
\end{minipage}%\hfill
\begin{minipage}[t]{0.70\textwidth} 
	{\scriptsize Debe tener presente  que los números enteros con valor cero o flotantes, cadenas vacías {(\ttfamily ''}), listas ({\ttfamily []}), tuplas ({\ttfamily(,)}), diccionarios ({\ttfamily\{\}}) y conjuntos ({\ttfamily set()}) son todos falsos, pero no son iguales a ninguno. }
\end{minipage}
\vspace{0.2cm}




En Python, la sentencia {\ttfamily if } (del inglés si) se utiliza para ejecutar un bloque de código si y sólo si se cumple una determinada condición o test condicional. En tal sentido, podemos decir que  {\ttfamily if} se usa en la toma de decisiones.  La estructura básica de esta sentencia es la siguiente:

\begin{lstlisting}
	if test conditional:
	Ejecutar alguna accción
\end{lstlisting} 


En este capítulo aprenderá a escribir un test o prueba condicional, que le permitirá comprobar cualquier condición de interés. 
Empezando con sentencias muy simples, hasta crear una serie de declaraciones {\ttfamily if}  más compleja, y necesarias para identificar las condiciones exactas requeridas, combinadas con las estructuras de datos como listas y diccionarios. 


\section{Sentencia if, elif y else, y test de comparación}  % Comparar condiciones con

% Recuerde, print() es la función incorporada de Python para imprimir cosas, normalmente a tu pantalla.

En general las sentencias condicionales  permiten controlar el flujo lógico de los programas, de manera clara y compacta. La sentencia {\ttfamily if} también se conoce como una sentencia de toma de decisiones, ya que toma una decisión en base a una condición bien definida. El bloque de código dentro de la instrucción {\ttfamily if} se ejecuta si  la condición es evaluada como verdadera, y es ignorado  si la condición es false.\\

En cada sentencia condicional ({\ttfamily if, elif, else}) hay una prueba comparativa, que al ser evaluada resulta dos posibles resultados  {\ttfamily True} o {\ttfamily False}.  Generalmente una la expresión condicional es una comparación matemática  o expresión booleana.  A continuación se presentan algunas pruebas matemáticas más comunes y expresiones booleana que se emplean en los condicionales.


\subsection{Test condicionales}

En el corazón de cada declaración {\ttfamily if}, hay una expresión que puede evaluarse como
Verdadero o Falsa y se denomina prueba condicional o test condicional. Python usa los valores {\ttfamily True} y {\ttfamily False} para decidir si se  ejecutar o no el bloque de código en la declaración {\ttfamily if}. Si la prueba condicional se evalúa como Verdadera, Python ejecuta el bloque de código seguido en la declaración del {\ttfamily if}; si la prueba se evalúa como False, Python ignora el código seguido en la declaración {\ttfamily if}.\\



A continuación se presenta una tabla con los operadores relacionales o de comparación, que se usan para construir las pruebas o los test condicionales, y los operadores lógicos para las expresiones booleanas: 



\begin{center}
	\noindent Tabla \ref{relacinales}. Operadores relacionales o de comparación:\\
	{\footnotesize{{\ttfamily 
				\hspace{-0.7cm}	
				\begin{tabular}{c l l l l}	% white   % \hline linea superior o inferior
					\multicolumn{1}{|>{\columncolor{azul2}} l }{\color{black}\bf Símbolo}& 
					\multicolumn{1}{|>{\columncolor{azul2}} l }{\color{black}\bf Significado}& 
					\multicolumn{1}{|>{\columncolor{azul2}} l }{\color{black}\bf Ejemplo}& 
					\multicolumn{1}{|>{\columncolor{azul2}} l }{\color{black}\bf Resultado} \\ 
					==  &Igual que        & "tres{}" {}== 3        & Falso      \\
					!=  & Distinto que    & rojo != verde & Verdadero    \\
					<   &  Menor que      &  7 < { }15       & Verdadero    \\
					<=  & Menor o igual que & 11 <= 11    & Verdadero    \\   
					>   & Mayor que       &     8 >{ }14    & Falso        \\
					>=  & Mayor o igual que &  13 >= 10    & Verdadero     \\          
					\arrayrulecolor{vino}\hline    
				\end{tabular}\label{relacinales} %	\hfill 	
				%  	\vspace{0.2cm} 
	}}} \\
\end{center}


Tenga en cuenta que se utilizan dos signos de igualdad ({\ttfamily ==}) para comparar igualdad; y una sola  ({\ttfamily =})  para asignar un valor a una variable.


\subsection{Expresiones booleanas}\label{logico}

Un valor booleano tiene dos posibles valores: Verdadero o Falso ({\ttfamily True} o {\ttfamily False}); un condicional es como  una expresión booleana, con dos posibles resultados {\ttfamily True} o {\ttfamily False}. Los valores booleanos se utilizan a menudo para realizar un seguimiento de ciertas condiciones, como en un juego que está en  ejecución o  un usuario que puede editar cierto contenido en una página web:

\begin{lstlisting}
	juego\_active = Verdadero
	puede\_edit = Falso
\end{lstlisting}Los valores booleanos proporcionan una manera eficiente  para rastrear el estado de un programa,
o una condición particular que sea importante en un programa.\\



Si necesita hacer varias comparaciones en un mismo condicional, se puede utilizar los operadores booleanos, {\ttfamily and} o {\ttfamily \&}, {\ttfamily or}, {\ttfamily xor}, {\ttfamily not}, efectuar la operación y luego determinar el resultado final,  considerado {\ttfamily 1 = true} y {\ttfamily 0 = False} .



\begin{center}
	\noindent Tabla \ref{logicos}. Operadores lógicos:\\
	{\footnotesize{{\ttfamily 
				\hspace{-0.7cm}	
				\begin{tabular}{c c l c l} % \arrayrulecolor{vino}\hline % white   % \hline linea superior o inferior
					\multicolumn{1}{|>{\columncolor{azul2}} l }{\color{black}\bf Operador}& 
					\multicolumn{1}{|>{\columncolor{azul2}} l }{\color{black}\bf Significado}& 
					\multicolumn{1}{|>{\columncolor{azul2}} l }{\color{black}\bf Ejemplo}& 
					\multicolumn{1}{|>{\columncolor{azul2}} l }{\color{black}\bf Operación}& 
					\multicolumn{1}{|>{\columncolor{azul2}} l }{\color{black}\bf Resultado} \\ 
					%	\arrayrulecolor{vino}\hline 
					and &  y  & 5 == 7 and 7 < { }12  & 0 y 1  & Falso   \\
					\&	&     & 9 < { }12  and 12 > { }7 & 1 y 1  & Verdadero   \\
					&     & 9 < { }12 and 12 > { }15 & 1 y 0  & Falso    \\
					\arrayrulecolor{vino}\hline   
					or  & 0   & 12 == 12 or 15 <{ } 7 &  1 o 0 &  Verdadero  \\   
					|   &     & 8 > { }5 or 4 < { }11    & 1 o 1 & Verdadero      \\
					\arrayrulecolor{vino}\hline    
					xor  & 0 excluyente   &  4 == 4 xor 9 > { }3 & 1 o 1 & Falso  \\ 
					&  & 4 == 4 xor 9 < { }3 & 1 o 0 &Verdadero  \\ 
					-  not & negación  & not 4 < { }20 & 0 & False   \\  
					in & está en  & item en lista &  &  False o True  \\     
					\arrayrulecolor{vino}\hline 				
				\end{tabular}\label{logicos} %	\hfill 	
				%  	\vspace{0.2cm} 
	}}} \\
\end{center}


% in, is  Membership 


\subsubsection{Recomendación para  hábitos en los test if}

En cada ejemplo de esta capítulo, has visto buenos hábitos de estilo. La única recomendación que PEP-8 proporciona para las pruebas condicionales de estilo, es utilizar un sólo espacio alrededor de los operadores de comparación, como {\ttfamily ==, >=, !=} y {\ttfamily <=}.

Por ejemplo: el test {\ttfamily if edad < 4},  es mejor que {\ttfamily if edad<4}.  Esta recomendación  no afecta la forma en que Python interpreta código, es sólo que el código se hace más fácil de leer y entender tanto para Usted como para otros programadores que en el futuro lo leerán.




A continuación se introducen algunos ejemplos para ilustrar el uso de estos operadores.   


\subsubsection{Test por igualdad {\ttfamily ==}}

La mayoría de las pruebas condicionales comparan el valor de una variable con un valor de interés específico. La prueba condicional más simple comprueba si el valor de un variable es igual al valor de interés, ejemplo:

\begin{lstlisting}
	In[]:
	>>> # Asigna valor
	>>> prueba = 'Verdades'
	>>> # y compara por igualdad  
	>>> prueba  == 'Verdades'	
	True
\end{lstlisting}

En la primera línea de código se asigna el valor a la variable {\ttfamily prueba} igual a {\ttfamily {}"{}Verdades{}"{}} (usando un sólo  signo igual); en la siguiente línea de código se compara si el valor de la variable es igual al valor  {}"{}Verdades{}"{}  mediante el uso de un doble signo igual ({\ttfamily==}). Este operador de igualdad devuelve un valor booleano {\ttfamily True},  si los valores en el lado izquierdo y derecho del operador coinciden, y   {\ttfamily False} si no coinciden.

Los valores en este ejemplo coinciden, por lo que Python devuelve el valor  {\ttfamily True}. Cuando el valor de la variable prueba no coincide, la prueba devuelve el valor {\ttfamily False}:

\begin{lstlisting}
	In[]:	
	>>> prueba = 'Mentiras'
	>>> prueba == 'Verdades'	
	False
\end{lstlisting}

Un único signo igual es en realidad una declaración; podrías leer la primera línea código como: se configura el valor de prueba igual a  "Verdades". Por otra lado, un doble signo igual plantea una pregunta: "¿El valor de {\ttfamily prueba}  es igual a {}"{}Verdades{}"{}?. 
La mayoría de los lenguajes de programación utilizan los signos iguales de esta forma.

\subsubsection{Las pruebas por igualdad son sensibles a mayúsculas y minúsculas}

Las pruebas de igualdad distinguen entre mayúsculas y minúsculas en Python. Por ejemplo,
los siguientes valores se consideran diferentes: 

\begin{lstlisting}
	In[]:
	>>> prueba = 'Verdades'
	>>> prueba == 'verdades'
	False
\end{lstlisting}

Si las mayúsculas y minúsculas son importante, este comportamiento es ventajoso; pero si no es importante
y en lugar de eso sólo quieres probar el valor de una variable, puedes convertir el valor de la variable en minúsculas antes de hacer la comparación:

\begin{lstlisting}
	In[]:
	>>> prueba = 'Verdades'
	>>> prueba.lower() == 'verdades'
	True
	>>> prueba
	'Verdades'
\end{lstlisting}

Esta prueba devuelve {\ttfamily True}, debido a que el método {\ttfamily .lower()}  devuelve una nueva cadena en minúsculas. Recuerde que las cadenas son  estructuras inmutables; el método {\ttfamily lower()} no cambia el valor original de la variable,  como se puede ver en la última linea de código.\\


Los sitios web imponen ciertas reglas de similitud  a los datos ingresados por los usuarios. Por ejemplo, un sitio podría utilizar una prueba condicional como esta para asegurarse de que cada usuario tenga un nombre de usuario realmente único, no sólo una variación de la mayúscula del nombre de usuario de otra persona. Cuando alguien presenta un nuevo nombre de usuario, ese nuevo nombre de usuario se convierte a minúsculas y se compara a las versiones en minúsculas de todos los nombres de usuario existentes. Durante esta verificación, un nombre de usuario como {}"{}Carlos21{}"{} se rechazarán si ya se utiliza alguna variación de {}"{}carlos21{}"{}.



\subsubsection{Test por desigualdad}

Cuando sea necesario determinar si dos valores son diferentes, se puede utilizar el operador desigualdad ({\ttfamily !=}). Veamos un ejemplo con el condicional {\ttfamily if} para visualizar la aplicación del operador desigualdad. \\

Consideremos un ejemplo en una pizzería. Supongamos un aderezo para pizza es solicitado  y registrado en una variable y luego se imprimir un mensaje si la persona no ordenó anchoas (la orden es diferente de anchoas ):

\hspace{-2.5cm}{\ttfamily \scriptsize toppings.py}  \vspace{-0.5cm}
\begin{lstlisting}	
	# La persona ordena champiñones	
	ordeno_topping = 'champiñones'
	if ordeno_topping != 'anchoas':
	print("¡Guarde las anchoas!")
\end{lstlisting}

Este código compara el valor de la variable {\ttfamily ordeno\_topping} con el valor {\ttfamily 'anchoas'}.
Si estos  valores no coinciden, Python devuelve {\ttfamily True} y ejecuta el código 
seguido en la declaración {\ttfamily if}. Si los dos valores coinciden, Python devuelve {\ttfamily False} y
no ejecuta el código seguido en la declaración {\ttfamily if}. 

En este ejemplo, como el valor de {\ttfamily ordeno\_topping} es diferente de {\ttfamily 'anchoas'}, se ejecuta la función {\ttfamily  print("¡Guarde las anchoas!")}:\vspace{0.08cm}

{\ttfamily\scriptsize\scriptsize \noindent Out[]:\vspace{-0.3cm}	
	{\color{blue} 	
		\begin{verbatim}
			¡Guarde las anchoas!			
		\end{verbatim}
}}


La mayoría de las expresiones condicionales que se escriben probarán la igualdad, pero 
algunas  veces le resulta más eficaz comprobar la desigualdad.

\subsubsection{Comparaciones numéricas}

Probar valores numéricos es bastante sencillo. Por ejemplo, el siguiente
El código comprueba si una persona tiene 18 años:

\begin{lstlisting}
	In[]:
	>>> edad = 18
	>>> edad == 18
	True
\end{lstlisting}

También puedes probar para verificar si dos números no son iguales. Por ejemplo, el
siguiente código imprime un mensaje si la respuesta dada no es correcta:

\hspace{-2.5cm}{\ttfamily \scriptsize concurso.py}  \vspace{-0.62cm}
\begin{lstlisting}		
	respuesta = 17
	if respuesta != 42:
	print("Esa no es la respuesta correcta. ¡Inténtelo de nuevo!")	
\end{lstlisting}	


El condicional pasa la prueba, porque el valor de la respuesta 17 diferente de 42. El resultado del condicional es  {\ttfamily True}, se ejecuta código: 

{\ttfamily\scriptsize\scriptsize \noindent Out:\vspace{-0.3cm}	
	{\color{blue} 	
		\begin{verbatim}
			Esa no es la respuesta correcta. ¡Inténtelo de nuevo!		
		\end{verbatim}
}}

Puedes incluir varias comparaciones matemáticas en la declaración del condicional, 
con los operadores  menor que, menor  o igual que, mayor que, mayor o igual que.
Veamos algunos ejemplos de comparaciones o test matemáticos, el resultado de la comparación es un boolenao, True o False:

\begin{lstlisting}
	In[]:	
	>>> edad = 19
	>>> edad < 21
	True
	>>> edad <= 21
	True
	>>> edad > 21
	False
	>>> edad >= 21
	False
\end{lstlisting}

Cada comparación matemática se puede incluir como parte de una declaración {\ttfamily if}, lo
que puede facilitar la precisión de la expresión de interés.

\subsection{Comprobación de múltiples condiciones}

Es posible que se necesite comprobar varias condiciones al mismo tiempo. Por ejemplo,
dos condiciones que sean verdaderas para realizar una acción, o tal vez se necesita que una sola condición de las dos condiciones sea Verdadera. En estas situaciones los operadores lógicos  {\ttfamily and} y {\ttfamily or}  son los ideales para resolver el problema.

\subsubsection{Múltiples comparaciones con {\ttfamily and}}


Para comprobar si dos condiciones son verdaderas simultáneamente, se puede hacer uso de los operadores lógicos  {\tt and} y {\ttfamily or} para combinar los dos condicionales; se evalúa cada test, y luego se evalúa el resultado de la expresión combinada (general). Si ambas son verdaderas, la expresión es verdadera; si cualquiera de las pruebas falla o si ambas fallan, la expresión se evalúa como False.

Por ejemplo, puedes comprobar si dos personas tienen más de 21 años:
\begin{lstlisting}
	In[]:
	>>> edad_0 = 22
	>>> edad_1 = 18
	>>> edad_0 >= 21 and edad_1 >= 21
	False
\end{lstlisting}

Primero, definimos las dos edades, {\ttfamily edad\_0} y {\ttfamily edad\_1}.  { }Luego comprobamos si ambas
edades son mayores o igual a 21. La prueba de la izquierda resulta verdadera (True),  la prueba de la derecha
resulta falsa (False), por lo que la expresión condicional general Falsa (False).\\



Ahora  hacemos {\ttfamily edad\_1 = 22}, de esta forma el valor de {\ttfamily edad\_1} es mayor que 21, por lo que ambos las pruebas son verdadras, lo que hace que la expresión condicional general se evalúe como Verdadera.

\begin{lstlisting}
	In[]:
	>>> edad_1 = 22
	>>> edad_0 >= 21 and edad_1 >= 21
	True
\end{lstlisting}

Para mayor legibilidad, puede utilizar paréntesis alrededor del cada test:
\begin{lstlisting}
	In[]:	
	>>> (edad\_0 >= 21) and (edad\_1 >= 21)
\end{lstlisting}
Python no es estricto, pero esta practica evita confusión, y facilitan la comprensión.



\subsubsection{Múltiples comparaciones con {\ttfamily or}}


El operador lógico  {\ttfamily or}  permite verificar múltiples condiciones,  es verdadera 
cuando al menos una es verdaderas. Y falsa cuando todas son falsas.

Consideremos nuevamente las dos edades, pero esta vez usaremos el operador {\ttfamily or}:


\begin{lstlisting}
	In[]:
	>>> edad_0 = 22
	>>> edad_1 = 20
	>>> edad_0 >= 21 or edad_1 >= 21
	True
\end{lstlisting}

La prueba de la izquierda es Verdadera ({\ttfamily  edad\_0 >= 21}), la prueba de la derecha y ({\ttfamily edad\_1 >= 21}) es Falsa, entonces la expresión general es Verdadera.\\ 


Ahora le asignamos a {\ttfamily  edad\_0} el valor de  18. De esta forma, las  dos pruebas son  Falsa, 
en consecuencia la expresión general es Falsa:
%   ({\ttfamily edad\_0 >= 21}) es falsa, y la prueba de la derecha {\ttfamily edad\_1 >= 21} también es
\begin{lstlisting}
	In[]:
	>>> edad_0 = 18
	>>> edad_1 = 20
	>>> edad_0 >= 21 or edad_1 >= 21
	False
\end{lstlisting}


\subsection{Comprobar si un valor está en una lista}

En algunas oportunidades es importante comprobar si un valor determinado está en una lista, para realizar una acción. Por ejemplo,
es posible que desee comprobar si el nombre de usuario ya existe en una lista de nombres de usuario  de algún en un sitio web,  antes de  registrarlo en la lista de usuarios.


En un proyecto de mapeo, es posible que quiera comprobar si una ubicación enviada ya existe en una lista de ubicaciones conocidas. Para saber si un valor en particular ya está en una lista, puede usar el operador
{\tt in}. \\

Consideremos en una pizzería, un ejemplo con una lista de ingredientes para pizza.  Suponga que el cliente ha solicitado para una pizza,
y luego verifique si ciertos ingredientes están en la lista.

\begin{lstlisting}
	In[]:
	>>> ingredientes = ['champiñones','hongos', 'cebollas', 'piña']
	>>> 'hongos' in  ingredientes
	True
	>>> 'pepperoni' in ingredientes
	False 
\end{lstlisting}

En este caso la lógica de la instrucción verifica la existencia de {}"{}hongos{}"{} 
y {}"{}pepperoni{}"{} en la lista de ingredientes.  El resultado es booleano.\\ 


Esta técnica es muy poderosa,  puedes crear una lista de valores esenciales y luego comprobar fácilmente si un determinado valor está la lista.

\subsubsection{Comprobar si un valor no está en una lista}

Otras veces, es importante saber si un valor no aparece en una lista.  Para este caso se puede utilizar el operador {\ttfamily not}. Por ejemplo, considere una lista de usuarios que tienen prohibido comentar en una foro. Puedes comprobar si un usuario ha sido prohibido antes de permitirle enviar un comentario:

\hspace{-2.8cm}{\ttfamily \scriptsize prohibidos.py}\vspace{-0.5cm}
\begin{lstlisting} 
	prohibidos = ['carlos', 'carolina', 'david', 'josé']
	usuario = 'maría'
	if usuario not in prohibidos:
	print(f"{usuario.title()}, puedes publicar una respuesta si lo deseas.")
\end{lstlisting}	

Aquí la declaración {\ttfamily if}  se lee con bastante claridad. Si el valor del usuario no está en
la lista prohibidos, Python devuelve {\ttfamily True} y ejecuta el bloque de código seguido del {\ttfamily if}.

La usuaria {}"{}maría{}"{ } no está en la lista prohibidos, por eso ve un mensaje invitándola a publicar una respuesta:

{\ttfamily\scriptsize\scriptsize \noindent Out[]:\vspace{-0.3cm}	
	{\color{blue} 	
		\begin{verbatim}
			María, puedes publicar una respuesta si lo deseas.	
		\end{verbatim}
}}


\section{Sentencia {\ttfamily if} en su forma simple}


En Python la forma más simple de una sentencia o declaración condicional es con un {\ttfamily if} seguido por la condición o test a evaluar y dos puntos (:) al final; luego, en la siguiente línea y con una {\bf indentación con cuatro espacios}, y el código  (o bloque de código) a ejecutar en caso de que se cumpla la condición.

Consideremos el siguiente ejemplo con señales de colores:

\hspace{-2.3cm}{\ttfamily \scriptsize colores.py}\vspace{-0.5cm}
\begin{lstlisting}
	entrada = "verde"
	if entrada == "verde":
	print("Continúe, paso libre") 
\end{lstlisting}


\noindent En este ejemplo el condicional o  prueba condicional es una comparación que usa el operador igualdad, si la entrada es igual a  "{}verde{}"{} ({\ttfamily entrada == {}"{}verde{}"{}}). Se evalúa la prueba, si resulta verdadero ({\ttfamily True})  se ejecuta la acción, que en éste caso es imprimir el mensaje: {}"{}Continúe, paso libre"; si la condición resulta falsa no se ejecuta la acción y Python sale del condicional.
\noindent En este caso, el test es verdadero ({\ttfamily True}), se imprime el mensaje:{\vspace{0.08cm}
	
	{\ttfamily\scriptsize\scriptsize \noindent In[]:\vspace{-0.3cm}	
		{\color{blue} 	
			\begin{verbatim}
				Continúe, paso libre
			\end{verbatim}
	}}
} 
Observe que el la doble igualdad ({\ttfamily ==}) es el operador de comparación (compara si es exactamente igual),  si colocas simplemente la igualdad ({\ttfamily =}) en el {\ttfamily if}, le generará un error, porque la igualdad  es el operador de asignación, asigna un valor a la variable.

Y es muy importante la indentación del código, juega el mismo papel en las declaraciones {\ttfamily if} que en los bucles {\ttfamily for}.
Todo  el bloque de líneas indentadas después de una declaración {\ttfamily if} se ejecutarán si la prueba es verdadera, y  se ignorará si la prueba es falsa.\\


Consideremos un segundo ejemplo, esta vez una variable que representa la edad de una persona, y supongamos que estamos interesados en conocer si esa persona tiene suficiente edad  para votar. El siguiente código prueba si
la persona puede votar, con el test de prueba ({\ttfamily edad >= 18}):

\hspace{-2.5cm}{\ttfamily \scriptsize votacion.py}\vspace{-0.5cm}
\begin{lstlisting}	
	edad = 19
	if edad >= 18:
	print("\n¡Tienes edad suficiente para votar!")
\end{lstlisting}Python verifica si el valor de edad es mayor o igual a 18; en caso de se verdad, ejecuta la función  {\ttfamily print({}"{}¡Tienes edad suficiente para votar!{}"{})}, que se encuentra en el bloque indentado.\\

Puedes poner cualquier test condicional en la primera línea y cualquier acción en la siguiente línea después del test y los dos puntos, respetando los cuatro espacios de sangría. \\

Agreguemos otra línea de código al ejemplo anterior, suponiendo que  la persona tiene suficiente edad para votar, debe estar registrado para ejercer el voto, así que podemos preguntarle ¿Ya estás registrado para votar?:

\hspace{-2.5cm}{\ttfamily \scriptsize votacion.py}\vspace{-0.5cm}
\begin{lstlisting}			
	edad = 19
	if edad >= 18:	
	print("¡Tienes edad suficiente para votar!")
	print("¿Ya te has registrado para votar?")
\end{lstlisting}
{\ttfamily\scriptsize\noindent Out[]:\vspace{-0.15cm}	
	{\color{blue} 
		\begin{verbatim}
			¡Tienes edad suficiente para votar!
			¿Ya te registraste para votar?
		\end{verbatim}
}}
En este caso el condicional resulta verdadero, se ejecuta bloque de código dentro de indentación.
Si el valor de la edad es inferior a 18 años, este programa no producirá ningún resultado.\\


Generalmente cuando se presenta un condicional, siempre hay una acción por realizar, si resulta verdadero ejecuta una acción, y si resulta falso ejecuta otra acción. La sintaxis {\ttfamily if-else} de Python permite este tipo de condicional.

\section{Sentencia {\ttfamily if-else}}
%  https://datascientest.com/es/python-if-else 

A la sentencia {\ttfamily if} se le puede añadir opcionalmente la sentencia {\ttfamily else}, para formar el bloque {\ttfamily if-else} que es similar a una declaración {\ttfamily if}, pero la declaración {\ttfamily else}  permite definir una acción o un bloque  de acciones que se ejecutan cuando la prueba falla; es decir,  contiene el código a ejecutar en caso de que {\bf no se cumpla la condición} del  {\ttfamily if}. 

La sintaxis de esta estructura se expresa como se muestra a continuación:

\begin{lstlisting}
	if condición:
	ación A
	else:
	ación B
\end{lstlisting}

Para ilustrar estas declaraciones consideremos un ejemplo con la sentencia {\ttfamily else}, y definimos la condición como un test de comparación con el operador {\ttfamily ==}, que genera un valor booleano, {\ttfamily True} o {\ttfamily False}, veamos: \vspace{0.15cm}


\hspace{-2.5cm}{\ttfamily \scriptsize dos\_}

\vspace{-0.2cm}
\hspace{-2.5cm}{\ttfamily \scriptsize opciones.py}\vspace{-0.84cm}
\begin{lstlisting}
	prompt = input("¿Python es fácil? s/n ")
	if prompt == 's':
	print("¡Si, es muy fácil y agradable trabajar con Python!")
	else:
	print("Si le prestas más atención, seguro le ¡interesará!")
\end{lstlisting}

La condición (o test de comparación) en este ejemplo compara  dos cadenas, una es el texto ingresado por el usuario y la otra es el definido en el programa como "s".  El resultado  {\ttfamily True} o  {\ttfamily False}.

Observe, nuevamente la indentación es muy importante, después de la condición del {\ttfamily if}, y del {\ttfamily else}, luego de los dos puntos, y en la próxima línea van cuatro espacios de indentación, si no la respeta generará  error o un resultado no  esperado.

Siempre utilice cuatro espacios para sangría por cada bloque de código (aunque puedes usar cualquier sangría que desees, pero Python espera que seas coherente). Las líneas deben tener la misma sangría, y alinearse a la izquierda. Y tenga presente que el estilo recomendado en PEP-8, son  cuatro espacios. No uses tabulaciones, no mezcle tabulaciones con espacios; puede arruinar el recuento de sangrías.\\


Para continuar ilustrando el uso de este condicional, consideremos nuevamente el ejemplo de las votaciones. Ahora le agregaremos un mensaje en caso de que no tenga la edad suficiente para votar:


Si el condicional ({\ttfamily edad >= 18}) pasa la prueba  del {\ttfamily if}, se ejecuta el primer bloque indentado; en caso de que  la prueba resulta falsa, se ejecuta el bloque del {\ttfamily else}:


\hspace{-2.5cm}{\ttfamily \scriptsize votacion.py}\vspace{-0.53cm}
\begin{lstlisting}	 
	edad = 16
	if edad >= 18:
	print("¡Tienes edad suficiente para votar!")
	print("¿Ya te has registrado para votar?")
	else:
	print("Lo siento, eres demasiado joven para votar.")
	print ("¡Regístrese para votar tan pronto como cumpla 18 años!")
\end{lstlisting} 

Debido a que esta vez la edad es menor a 18 años, la prueba condicional falla; y se ejecuta el  bloque de código
del {\ttfamily else}:\vspace{0.1cm}

{\ttfamily\scriptsize \noindent Out[]:\vspace{-0.26cm}	
	{\color{blue} 	
		\begin{verbatim}			
			Lo siento, eres demasiado joven para votar.	
			¡Regístrese para votar tan pronto como cumpla 18 años!
		\end{verbatim}
}}

Esta estructura {\ttfamily if-else} funciona  bien en situaciones en las que Python tiene que ejecutar 
{\bf una de las dos} acciones posibles.\\

Sin embargo, hay situaciones que se  necesitan evaluar más de dos comparaciones,  en estos casos Python tiene la instrucción  {\ttfamily elif}. 

\vspace{0.2cm}
\begin{minipage}[t]{0.2\textwidth}
	\centering\raisebox{\dimexpr\topskip-\height}{%
		
\includegraphics[scale=0.20]{../Logos_de_Python/Busquedas_R}}  %[width=\textwidth]	
\end{minipage}%\hfill
\begin{minipage}[t]{0.70\textwidth} 
	{\scriptsize 	La sentencia {\ttfamily if-else} se utiliza para ejecutar uno de los dos bloques. Si la condición es verdadera, se ejecuta el bloque de código {\ttfamily if}; en caso contrario,  se ejecuta el bloque de código {\ttfamily else}. 
	}
\end{minipage}
\vspace{0.2cm}


\section{La sentencia {\ttfamily if-elif-else}}

A la sentencia {\ttfamily if-else}, se pueden añadir un número indefinido de pruebas condicionales. Estas pruebas  se agregan mediante la instrucción {\ttfamily elif}, de acuerdo a la siguiente la sintaxis:

\begin{lstlisting}
	if condicion A1:
	Acción A1
	elif condicion A2:
	Acción A2
	.
	.
	.     
	elif condicion An:
	Acción An
	else:
	Acción B
\end{lstlisting}

Esta instrucción corre en orden, probando cada test condicional hasta que aparece la primera prueba que resulta verdadera.  Cuando el primer condicional es verdadero, el bloque de código que sigue a esa prueba se ejecuta, y Python salta el resto de las pruebas. En otras palabras Python ejecuta sólo un bloque en la cadena de código {\ttfamily if-elif-else}.\\


Muchas casos de programación real se presenta el problema de evaluar más de dos condiciones, esta situación se puede resolver con una cadena de código {\ttfamily if-elif-else}.  Por ejemplo, considere un parque de diversiones donde las tarifas de entrada se cobra de acuerdo con la edad: La entrada para menores de 4 años es libre; la entrada para personas entre 4 y 18 años cuesta \$25; y la entrada para personas mayores de 18 años cuesta \$40.

¿Cómo podemos utilizar una declaración {\ttfamily if} para determinar el costo de la entrada de una persona?

El siguiente código prueba la edad de una persona y luego imprime un mensaje de precio de entrada:


\hspace{-2.4cm}{\ttfamily \scriptsize parque1.py}\vspace{-0.52cm}
\begin{lstlisting}		
	edad = 12
	if edad < 4:
	print("Su costo de entrada es 0\$.")
	elif edad < 18:
	print("Su costo de entrada es 25\$.")
	else:
	print("Su costo de entrada es 40\$.")
\end{lstlisting}	

El test del {\ttfamily if} compara si la persona tiene menos de 4 años. Cuando el test pasa la prueba,  imprime un mensaje apropiado y Python se salta el resto de pruebas. La línea del {\ttfamily elif} es en realidad otra prueba {\ttfamily if}, que se ejecuta sólo si la prueba anterior falla. En este punto de la cadena, sabemos que la persona es mayor de  4 años porque falló la primera prueba. Si la persona es menor de 18 años, se imprime un mensaje adecuado y Python omite el bloque del {\ttfamily else}. Si tanto el  {\ttfamily if} como el  {\ttfamily elif} fallan, Python ejecuta el código en el bloque {\ttfamily else}. En este ejemplo,  la prueba de comparación del {\ttfamily elif}  resulta verdad (12 es menor que 18), por lo que se ejecuta el código. Y el resultado es una frase que informa al usuario el costo de la entrada:\vspace{0.08cm}

{\ttfamily\scriptsize\noindent Out[]: \vspace{-0.3cm}	
	{\color{blue} 	
		\begin{verbatim}
			Su costo de entrada es de 25\$.			
		\end{verbatim}
}} 

\noindent Cualquier edad mayor a 17 años hace que las dos primeras pruebas fallen. En estas
condiciones, el bloque {\ttfamily else} se ejecutará y el precio de entrada será 40\$.\\

Este código se puede simplificar, en lugar de imprimir el precio de entrada dentro del bloque {\ttfamily  if-elif-else}, es más conciso establecer solo el precio dentro de la cadena {\ttfamily  if-elif-else} y hacer una única llamada {\ttfamily print()} que se ejecute después de que la condición haya sido evaluada:

\hspace{-2.2cm}{\ttfamily \scriptsize parque.py}\vspace{-0.52cm}
\begin{lstlisting}		
	edad = 12
	if edad < 4:
	precio = 0
	elif edad < 18:
	precio = 25
	else:
	precio = 40
	print(f"El costo de tu entrada es {precio}\$.")
\end{lstlisting} 

{\ttfamily\scriptsize\noindent Out[]: \vspace{-0.3cm}	
	{\color{blue} 	
		\begin{verbatim}
			El costo de tu entrada es \$25.			
		\end{verbatim}
}} 

Las líneas indentadas establecen el valor del precio según la edad de la persona,
como en el ejemplo anterior.  Después de que la cadena  {\ttfamily  if-elif-else} fija el precio, se la llamada la función {\ttfamily print()} fuera del bloque de la cadena {\ttfamily  if-elif-else}, y se  usa el valor del  precio para mostrar el reporte del precio en el mensajes.

Este código produce el mismo resultado que el ejemplo anterior, pero el
el propósito de la cadena {\ttfamily  if-elif-else} es más corto. En lugar de determinar un precio
y mostrar un mensaje, simplemente determina el precio de la entrada. Además de ser más eficiente, este código  es más fácil de modificar que el enfoque original. Por ejemplo, para cambiar el texto del mensaje de salida, solo será necesario cambiar una vez el {\ttfamily print()}, en lugar de tres como seria en el caso anterior.


\vspace{0.2cm}
\begin{minipage}[t]{0.2\textwidth}
	\centering\raisebox{\dimexpr\topskip-\height}{%
		
\includegraphics[scale=0.20]{../Logos_de_Python/Busquedas_R}}  %[width=\textwidth]	
\end{minipage}%\hfill
\begin{minipage}[t]{0.70\textwidth} 
	{\scriptsize La cadena de código {\ttfamily if-elif-else} ejecutará sólo  una condición,  la primera condición que se evalúe como {\ttfamily True}. Aunque haya otra prueba que sea verdadera ({\ttfamily True}), el bloque  {\ttfamily if-elif-else} termina después de la primera condición verdadera. En caso de que todas las prueba de la cadena {\ttfamily if-elif} sean falsas, se ejecuta el bloque {\ttfamily else}. Esto significa que el bloque de la sentencia {\ttfamily else} se ejecuta sólo en el caso de que ninguna  prueba de la cadena {\ttfamily if-elif}  sea verdadera.}
\end{minipage}
\vspace{0.2cm}

\subsection{Usando múltiples bloques {\ttfamily elif}}


Puedes usar tantos bloques {\ttfamily elif} en el código como sean necesarios. Por ejemplo, si el
parque de atracciones implementara un descuento para personas mayores, se podría agregar
un test condicional más al código para determinar si alguien califica para el descuento de personas mayores. Digamos que cualquier persona de 60 años o más paga un precio especial de 20\$:

\hspace{-2.2cm}{\ttfamily \scriptsize parque.py}\vspace{-0.5cm}
\begin{lstlisting}	 
	edad = 12
	if edad < 4:
	precio = 0
	elif edad < 18:
	precio = 25
	elif  edad < 60:
	precio = 40
	else:
	precio = 20
	
	print(f"El costo de tu entrada es {precio}\$.")   
\end{lstlisting}    


La mayor parte de este código no ha cambiado. El segundo bloque {\ttfamily elif} ahora asegura que una persona tenga menos de 60 años antes de asignarle la entrada completa de 40\$. Observe que el valor asignado en el bloque  {\ttfamily else} debe cambiarse a 20\$, porque las únicas edades que llegan a este bloque son para personas de 60 años o más.


\subsubsection{Omitiendo el bloque {\ttfamily else}}

Python no requiere estrictamente la instrucción  {\ttfamily else} al final de una cadena {\ttfamily  if-elif}.
Algunas  veces, un bloque {\ttfamily else}  es útil, otras veces es mejor utilizar una 
declaración {\ttfamily elif}  que capta la condición específica de interés. En el ejemplo del parque de diversiones se presenta una situación ejemplar:

\hspace{-2.2cm}{\ttfamily \scriptsize parque.py}\vspace{-0.5cm}
\begin{lstlisting}
	edad = 12
	if edad < 4:
	precio = 0
	elif edad < 18:
	precio = 25
	elif  edad < 60:
	precio = 40
	elif  edad >= 60:
	precio = 20
	
	print(f"El costo de tu entrada es {precio}\$.")    
\end{lstlisting}

Al final del bloque  {\ttfamily elif}  se asigna el precio de \$20 cuando la persona tiene 60 o más edad, 
es un poco más claro que el bloque general con la sentencia  {\ttfamily else}.  




\subsubsection{Probando Múltiples Condicionales}

La cadena  {\ttfamily if-elif-else} es poderosa, y es apropiada para situaciones donde sólo una de las prueba debe ser verdadera. 
Las demás pruebas Python las ignora, luego de encontrar  la primera prueba que resulta verdadera y sale.   Este comportamiento es beneficioso porque es eficiente y le permite realizar pruebas para una condición específica. Sin embargo, si se necesita comprobar todas las condiciones de interés, este instrucción no es útil.

En este caso, debe utilizar una serie de declaraciones {\ttfamily if} simples sin los bloques del  {\ttfamily elif}  y 
del {\ttfamily else}.  Esta técnica tiene sentido cuando más de una condición debe ser probada y ejecutar alguna acción en cada uno de los test que resulte verdadero.\\


Consideremos nuevamente el ejemplo de la pizzería. Si alguien ordena  dos ingredientes para la 
pizza, deberás asegurarte de incluir ambos ingredientes en su pizza: 

\hspace{-2.1cm}{\ttfamily \scriptsize  pizza.py}\vspace{-0.52cm} 
\begin{lstlisting}
	ingredientes = ['champiñones', 'queso extra']
	if 'champiñones' in  ingredientes:
	print("Agregando champiñones.")     
	if 'pepperoni' in ingredientes:
	print("Agregando pepperoni.")     
	if 'queso extra'in ingredientes:
	print("Agregando queso extra.")
	
	print("\n¡Terminaste de hacer tu pizza!")
\end{lstlisting}



Comenzamos con una lista que contiene los ingredientes (toppings) ordenados de la pizza. En el primer  {\ttfamily if},  la prueba verifica si la persona ordeno  champiñones (hongos) en su pizza. Si es así, se imprime un mensaje confirmando ese aderezo.   El segundo  {\ttfamily if}, el test del pepperoni, es otra declaración  simple, no una declaración  {\ttfamily elif}  o  {\ttfamily else}, por lo que  esta prueba se ejecuta independientemente de la prueba anterior. La última declaración  {\ttfamily if} verifica si se ordeno queso adicional, sin considerar   los resultados de las dos primeras pruebas.  Estas tres pruebas independientes se ejecutan cada vez que se ejecuta este programa. 

Debido a que se evalúan todas las condiciones en este ejemplo,  champiñones  y queso extra se agrega  a la pizza:\vspace{0.09cm}

{\ttfamily\scriptsize\noindent Out[]: \vspace{-0.3cm}	
	{\color{blue} 	
		\begin{verbatim}
			Agregando champiñones.
			Agregando queso extra.
			
			¡Terminaste de hacer tu pizza!
		\end{verbatim}
}} 

Este código no funcionaría correctamente si usáramos un bloque {\ttfamily if-elif-else},
porque el código dejaría de ejecutarse después de pasar una prueba. Aquí está
cómo se vería eso:

\hspace{-2.2cm}{\ttfamily \scriptsize  pizza.py}\vspace{-0.5cm}  
\begin{lstlisting}
	ingredientes = ['champiñones', 'queso extra'] 
	if 'champiñones' in  ingredientes:
	print("Agregando champiñones.")  
	elif 'pepperoni' in ingredientes:
	print("Agregando pepperoni.")  
	elif 'queso extra'in ingredientes:
	print("Agregando queso extra.")
	
	print("\n¡Terminaste de hacer tu pizza!")
\end{lstlisting}

La prueba de  'champiñones' es la primera que se pasa, por eso se añaden champiñones
a la pizza. Sin embargo, los valores "queso extra" y "pepperoni" son ignorados, nunca se
ejecutan, porque Python no ejecuta ninguna prueba más allá de la primera prueba que
pasa en una cadena {\ttfamily if-elif-else}. Se agregará el primer ingrediente del cliente,
pero se perderán todos sus demás ingredientes:\vspace{0.09cm}


{\ttfamily\scriptsize\noindent Out[]:\vspace{-0.292cm}	
	{\color{blue} 	
		\begin{verbatim}
			Agregando champiñones.
			
			¡Terminaste de hacer tu pizza!
		\end{verbatim}
}} 


En resumen, si desea que sólo se ejecute un único bloque de código, puede usar la cadena {\ttfamily if-elif-else}; pero si necesita ejecutar más de un bloque de código, debe usar una serie de bloques independientes {\ttfamily if}.


\section{Uso de {\ttfamily if-else} sobre listas}

Se pueden hacer programas interesantes combinando listas y declaraciones {\ttfamily if}. Se pueden obtener algunos valores especiales que necesitan para ser tratados de manera diferente a los otros valores de la lista. Se puede gestionar eficientemente las condiciones dinámicas; como por ejemplo, la disponibilidad de ciertos item en un restaurante durante un turno.  


\subsection{Comprobando item  con {\ttfamily for} y el {\ttfamily if}}

Este capítulo comenzó con un ejemplo simple que mostró cómo manejar un especial
valor como 'prueba', que debía imprimirse en un formato diferente al
otros valores de la lista. Ahora que tienes un conocimiento básico de los  condicionales
pruebas y declaraciones {\ttfamily if}, echemos un vistazo más de cerca a cómo puede ver
para valores especiales en una lista y manejar esos valores apropiadamente.\\

Sigamos con el ejemplo de la pizzería. En este ejemplo se muestra un mensaje 
cada vez que se agrega un aderezo a la pizza, mientras se prepara. El código
para esta acción se puede escribir de manera muy eficiente haciendo una lista de ingredientes que
el cliente ha ordenado, y utilizando un bucle para anunciar cada ingrediente a medida que se
agregan a la pizza: 

\hspace{-2.5cm}{\ttfamily \scriptsize toppings.py}\vspace{-0.52cm}
\begin{lstlisting}		
	ingredientes = ['champiñones', 'pimientos verdes', 'queso extra']	
	for ingredente in ingredientes:
	print(f"Agregando {ingrediente}.")
	print("\n¡Terminaste de hacer tu pizza!")	
\end{lstlisting}	

El resultado es sencillo porque este código es simplemente un bucle {\ttfamily for} simple:\vspace{0.09cm}


{\ttfamily\scriptsize\noindent Out[]:\vspace{-0.292cm}	
	{\color{blue} 	 	
		\begin{verbatim}
			Agregando champiñones.
			Agregando pimientos verdes.
			Agregando queso extra.
			¡Terminaste de hacer tu pizza!		
		\end{verbatim}
}}

Ahora supongamos que la pizzería se queda sin pimientos verdes, y debemos dar respuesta al cliente. Esta situación se puede manejar con una declaración del {\ttfamily if} dentro del bucle {\ttfamily for}, veamos:


\hspace{-2.5cm}{\ttfamily \scriptsize toppings.py}  \vspace{-0.5cm}
\begin{lstlisting}		
	ingredientes = ['champiñones', 'pimientos verdes', 'queso extra']	
	for ingrediente in ingredientes:
	if ingrediente == 'pimientos verdes':   
	print("Lo sentimos, se nos acabaron los pimientos verdes en este momento.")
	else:
	print(f"Agregando {ingrediente}.")
	print("\n¡Terminaste de hacer tu pizza!")	
\end{lstlisting}	

Esta vez, revisamos cada elemento ordenado antes de agregarlo a la pizza. La declaración  {\ttfamily if} verifica si la persona solicitó pimientos verdes. En ese caso,  se  muestra un mensaje informándoles que no hay pimientos verdes.  El bloque {\ttfamily else} garantiza que todos los demás ingredientes se agregarán a la pizza.  El resultado muestra los ingrediente agregados a la pizza, que han sido seleccionados por el usuario.\vspace{0.09cm}


{\ttfamily\scriptsize\noindent Out[]:\vspace{-0.292cm}	
	{\color{blue} 		
		\begin{verbatim}
			Lo sentimos, se nos acabaron los pimientos verdes en este momento.
			Agregando queso extra.
			
			¡Terminaste de hacer tu pizza!			
		\end{verbatim}
}}


\subsection{Comprobar que una lista no esté vacía}  
% Hemos hecho una suposición simple sobre cada lista con la que hemos trabajado hasta ahora:

En los ejemplos anteriores hemos asumimos que cada lista tiene al menos un elemento. Sin embargo, al permitir que los usuarios
agreguen de los ingredientes para su pizza, no podremos esperar que lista contenga siempre los mismos elementos 
cada vez que se ejecuta un bucle. En este caso, es útil comprobar si la lista está vacía antes de ejecutar un bucle {\ttfamily for}.

Como ejemplo, verifiquemos si la lista de ingredientes  esta  vacía antes de comenzar la pizza. Si la lista está vacía, le indicaremos al usuario  si realmente quiere una pizza sencilla. Si la lista no está vacía, 
hacemos la pizza tal como lo hicimos en los ejemplos anteriores:

\hspace{-2.5cm}{\ttfamily \scriptsize sencilla.py} \vspace{-0.5cm}
\begin{lstlisting}	
	ingredientes = []
	if ingredientes:
	for ingrediente in ingredientes:
	print(f"Agregando {ingrediente}.")
	print("\n¡Terminaste de hacer tu pizza!")
	else:
	print("¿Estás seguro de que quieres una pizza sencilla?")
\end{lstlisting}

Esta vez comenzamos con una lista vacía de ingredientes solicitados. En lugar de ir directamente al bucle {\ttfamily for}, primero hacemos una chequeo rápido, sobre el pedido. 

Cuando una lista se encuentra en declaración {\ttfamily if} como un condicional, Python  evalúa y devuelve {\ttfamily True} si la lista contiene al menos  un elemento, y {\ttfamily True} si la lista está vacía.  Si el ingrediente está en la lista, pasa el test condicional, y se  ejecuta el mismo bucle {\ttfamily for} que usamos en la prueba del ejemplo anterior.  Si la condición de la prueba falla, imprimimos un mensaje preguntando al cliente quiere una pizza sencilla, sin aderezos.

En este ejemplo la lista está vacía, por lo que la salida es la pregunta al usuario si realmente quiere
una pizza sencilla:\vspace{0.09cm}


{\ttfamily\scriptsize\noindent Out[]:\vspace{-0.292cm}	
	{\color{blue} 		
		\begin{verbatim}
			¿Estás seguro de que quieres una pizza sencilla?			
		\end{verbatim}
}}
Si la lista no está vacía, el resultado mostrará cada ingrediente solicitado que se agrega a la pizza. 

\subsubsection{Usar múltiples listas}

La gente pedirá casi cualquier cosa, especialmente cuando se trata de aderezos para pizza.
Por ejemplo, un cliente  puede pedir papas fritas en su pizza u otro ingrediente no común;  para tratar esta situación se pueden usar listas y declaraciones {\ttfamily if} para asegurarse de que la entrada tenga sentido antes iniciar el proceso. Se verifican las solicitudes de ingredientes inusuales antes empezar la  preparación de la pizza.

El siguiente ejemplo define dos listas. La primera es una lista de ingredientes disponibles en 
la pizzería, y la segunda es la lista de ingredientes que el usuario ordena.
Esta vez, cada ingrediente ordenado se compara con la lista de ingredientes disponibles antes de agregarlos a la pizza:

\hspace{-2.7cm}{\ttfamily \scriptsize inusuales.py}\vspace{-0.55cm}
\begin{lstlisting}
	disponible = ['champiñones', 'aceitunas', 'pimientos verdes',
	'pepperoni', 'piña', 'queso extra']
	ordenados  =  ['champiñones', 'papas fritas', 'queso extra']
	
	for ordenado in ordenados:
	if ordenado in disponible:
	print(f"\nAgregando {ordenado}.")
	else:
	print(f""Lo siento, no tenemos {ordenado}.")
	
	print("\n¡Terminaste de hacer tu pizza!") 
\end{lstlisting}

Primero, definimos una lista de ingredientes disponibles en la pizzería. Observe,  puede ser una tupla si la pizzería tiene una selección estable de ingredientes. Hacemos una lista con los ingredientes que un cliente ha ordenó. 
Hay un pedido inusual como aderezo por ejemplo: 'papas fritas'. 

Hacemos un bucle {\ttfamily for} sobre la lista de ingredientes ordenados. Dentro del bucle, comprobamos si cada ingrediente ordenado está en la lista de ingredientes disponibles. Si es así, agregamos ese aderezo a la pizza. Si el ingrediente ordenado no está en la lista de ingredientes disponibles, el bloque {\ttfamily else} se ejecutará, y se imprime un mensaje para decirle al usuario ese ingrediente no están disponibles: "Lo siento, no tenemos papas fritas". 


Esta sintaxis de código produce una salida limpia e informativa:\vspace{0.08cm}

{\ttfamily\scriptsize\noindent Out[]:
	\vspace{-0.3cm}
	{\color{blue}		
		\begin{verbatim}
			Agregando champiñones.
			"Lo siento, no tenemos papas fritas"
			Agregando queso extra.
			
			¡Terminaste de hacer tu pizza!
		\end{verbatim}
	}
}

En sólo con unas pocas líneas de código, hemos manejado una situación del mundo real efectiva.

\subsection{Ejercicios}

Ejercicios 

\begin{enumerate}[6.5.1]
	
	\item   Pruebas o test condicionales: Escriba una serie de pruebas condicionales. Imprimir una declaración
	describiendo cada prueba y su predicción para los resultados de cada prueba. Tu código debería verse algo como esto:
	
	\begin{lstlisting}		
		coche = 'Spark'
		print("¿Es el auto == 'Spark'? Predigo que es Verdadero.")
		print(coche == 'Spark')
		print("\n¿Es auto == 'Toyota'? Predigo Falso.")
		print(coche == 'Toyota')
	\end{lstlisting}
	
	Revise bien los resultados y asegúrese de comprender por qué cada línea se evalúa como Verdadero o Falso.
	
	Cree al menos 10 pruebas: Haga que al menos 5 pruebas se evalúen como Verdadero y otra 5 pruebas se evalúan como Falso.
	
	
	\item  Más pruebas condicionales: No tiene que  limitarse a 10 pruebas, como mínimo 10. Si quieres probar más comparaciones, escribe más pruebas y agrégalas  al archivo {\ttfamily{tests.py}}. Tener al menos un resultado Verdadero y uno Falso para cada uno de la siguiente:
	
	Pruebas de igualdad y desigualdad con cadenas.
	
	Pruebas usando el método {\ttfamily{lower()}}.
	
	Pruebas numéricas que involucran igualdad y desigualdad, mayor que y menor 	que, mayor o igual que, y menor o igual que.
	
	Pruebas utilizando la palabra clave {\ttfamily{and}} y la palabra clave {\ttfamily{or}}.
	
	Probar si un elemento está en una lista.
	
	Probar si un elemento no está en una lista.	
	
	
	\item Alienígenas de Colores 1: Imagina que acaban de derribar a un extraterrestre en un juego. Crear un
	variable llamada {\ttfamily{alien\_color}} y asígnele un valor de 'verde', 'amarillo' o 'rojo'.	
	Escribe una declaración {\ttfamily{if}} para comprobar si el color del extraterrestre es verde. Si es así, imprima
	un mensaje de que el jugador acaba de ganar 5 puntos.	
	Escriba una versión de este programa que pase la prueba {\ttfamily{if}} y otra que
	falla. (La versión que falla no tendrá resultados).
	
	\item Alienígenas de Colores 2: Elige un color para un alienígena como lo hiciste en el ejercicio anterior,
	y escribe una cadena {\ttfamily{if-else}}.	
	Si el color del alienígena es verde, imprima una declaración de que el jugador acaba de ganar 5
	puntos por dispararle al extraterrestre.	
	Si el color del alienígena no es verde, imprima una declaración de que el jugador acaba de ganar
	10 puntos.	
	Escriba una versión de este programa que ejecute el bloque {\ttfamily{if}} y otra que
	ejecuta el bloque {\ttfamily{else}}.
	
	
	\item  Alienígenas de Colores 3: Convierte la cadena {\ttfamily{if-else}} del Ejercicio anterior en una cadena  {\ttfamily{if-elif else}}.	
	Si el alienígena es verde, imprime un mensaje de que el jugador obtuvo 5 puntos.	
	Si el alienígena es amarillo, imprime un mensaje de que el jugador obtuvo 10 puntos.	
	Si el alienígena es rojo, imprime un mensaje de que el jugador obtuvo 15 puntos.	
	Escriba tres versiones de este programa, asegurándose de que cada mensaje esté impreso para el color alienígena apropiado.
	
	\item Etapas de la vida: Escriba una cadena {\ttfamily{if-elif-else}} que determine la etapa de vida de una persona.
	Establezca un valor para la variable edad y luego:	
	Si la persona tiene menos de 2 años, imprima un mensaje de que la persona es un bebé.	
	Si la persona tiene al menos 2 años pero menos de 4, imprima un mensaje indicando la persona es un niño pequeño.	
	Si la persona tiene al menos 4 años pero menos de 13, imprima un mensaje que la persona es un niño.	
	Si la persona tiene al menos 13 años pero menos de 20, imprima un mensaje que la persona es un adolescente.	
	Si la persona tiene al menos 20 años pero menos de 65, imprima un mensaje que la persona es un adulto.	
	Si la persona tiene 65 años o más, imprima un mensaje que indique que la persona es un adulto mayor.	
	
	\item  Fruta favorita: Haga una lista de tus frutas favoritas y luego escriba una serie de
	declaraciones {\ttfamily{if}} independientes que verifican ciertas frutas en su lista.
	
	Haga una lista con tres frutas favoritas y llámela {\ttfamily{frutas\_favoritas}}.
	
	Escribe cinco declaraciones {\ttfamily{if}}. 
	Cada uno debería comprobar si un determinado tipo de fruta 	está en tu lista. 
	Si la fruta está en su lista, el bloque {\ttfamily{if}} debería imprimir una declaración,
	como ¡Te gustan mucho las Bananas!
	
	
	\item Hola administrador: Haga una lista de cinco o más nombres de usuario, incluido el nombre del 'administrador'. Imagine que está escribiendo un código que imprimirá un saludo para cada usuario, después de iniciar sesión en un sitio web. Recorra la lista e imprima un saludo para cada usuario.
	
	Si el nombre de usuario es 'administrador', imprima un saludo especial, como Hola 'Administrador',
	¿Le gustaría ver un informe de estado?
	
	De lo contrario, imprima un saludo genérico, como Hola Carlos, gracias por iniciar nuevamente sesión.
	
	\item NO  usuario: Agregue una prueba {\ttfamily{if}} a {\ttfamily{hola\_admin.py}} para asegurarse de que la lista de usuarios no este vacía.
	
	Si la lista está vacía, imprima el mensaje ¡Necesitamos encontrar algunos usuarios!
	
	Elimine todos los nombres de usuario de su lista y asegúrese de que aparezca el mensaje correcto que está impreso.
	
	
	\item Chequeando nombres de usuario: Haga lo siguiente para crear un programa que simule
	cómo los sitios web garantizan que todos tengan un nombre de usuario único.
	
	Haga una lista con cinco o más nombres de usuarios {\ttfamily{usuarios\_actuales}}, y haga otra lista de cinco nombres de usuario, y llamela  {\ttfamily{nuevos\_usuarios}}. Asegúrese de que uno o dos de los nombres nuevos de usuario también están en la lista {\ttfamily{usuarios\_actuales}}.
	
	Recorra la lista {\ttfamily{nuevos\_usuarios}} para verificar si ya está en uso (ha sido usado). 
	Si es así, imprima un mensaje indicando que la persona deberá ingresar un nuevo nombre de usuario. Si un nombre de usuario no se ha  utilizado, imprima un mensaje indicando qué  nombre de usuario está disponible.
	
	Asegúrese de que su comparación no distinga entre mayúsculas y minúsculas. Si se ha utilizado 'John',
	'JOHN' no debe aceptarse. (Para hacer esto, necesitará hacer una copia de  {\ttfamily{usuarios\_actuales}} que 
	contiene las versiones en minúsculas de todos los usuarios existentes).
	
	\item  Números ordinales: Los números ordinales indican su posición en una lista, como $1\textsuperscript{\underline{o}}$, $2\textsuperscript{\underline{o}}$,  etc.
	
	Almacene los números del 1 al 9 en una lista.
	
	Recorra la lista (bucle).
	
	Utilice una cadena {\ttfamily{if-elif-else}} dentro del bucle para imprimir el final ordinal adecuado
	para cada número. Su salida debería decir  "$1\textsuperscript{\underline{o}}$, $2\textsuperscript{\underline{o}}$, 
	$3\textsuperscript{\underline{o}}$, $4\textsuperscript{\underline{o}}$, $5\textsuperscript{\underline{o}}$, $6\textsuperscript{\underline{o}}$, $7\textsuperscript{\underline{o}}$, $8 \textsuperscript{\underline{o}}$ y   $9\textsuperscript{\underline{o}}$", y cada resultado debe estar en una línea separada.
	
	
	\item  Diseñar declaraciones {\ttfamily{if}}: Revise los programas que escribió en este capítulo y
	asegúrese de haber diseñado sus pruebas condicionales de manera adecuada.
	
	\item Sus ideas: En este punto, eres un programador más capaz que  cuando comenzó  este libro. Ahora que tienes una mejor idea de cómo son las situaciones del mundo real, como  se modelan en programas, es posible que esté pensando en algunas
	problemas que podrías resolver con tus propios programas. Registre cualquier idea nueva que tenga sobre algún problema que  quieras resolver, y a medida que sus habilidades de programación avancen puedas mejorar. Considere los juegos que tal vez quiera escribir, los conjuntos de datos que tal vez quiera explorar o escribir, y las aplicaciones web que desea crear.
	
	
	\item  Sándwich: Haga una lista llamada {\ttfamily{ordenes\_sándwich}} y rellénela con los  varios nombres de
	sándwiches. Luego haga una lista vacía llamada {\ttfamily{sándwich\_terminados}}. 
	Recorra la lista  de pedidos de sándwiches en bucle, e imprimir un mensaje para cada pedido, 
	como el que hizo el  sándwich de atún. A medida que se prepara cada sándwich, muévalo a la
	lista {\ttfamily{sándwich\_terminados}}. Una vez hechos todos los sándwiches, imprima una lista de mensajes para cada sándwich que se hizo.
	
	\item  Sin cebolla: Utilizando la lista {\ttfamily{ordenes\_sándwich}} del ejercicio 7.8, asegúrese
	el sándwich con 'cebolla' aparezca  en la lista al menos tres veces. Agregue un código cerca del comienzo de su programa
	para imprimir un mensaje que diga que el cafetín de Sandwich  no tiene cebollas, y luego usar un bucle {\ttfamily{while}} para eliminar todas las apariciones de 'cebolla' en la lista  {\ttfamily{ordenes\_sándwich}}. Asegúrate de que no hayan los sándwich con cebolla	en  {\ttfamily{sándwich\_terminados}}.
	
	\item Vacaciones de ensueño: Escriba un programa que encueste a los usuarios sobre las vacaciones de sus sueños.
	
	Escribe un mensaje similar a: Si pudieras visitar un lugar en el mundo, ¿ a donde irías? 
	Incluya un bloque de código que imprima los resultados de la encuesta.
	
	
\end{enumerate}	




%%%%%%%%%%%%%%%%%%%%%%%%%%%%%%%%%   Comprensiones  %%%%%%%%%%%%%%%%%%%%%%%%%%%%%%%%%%%%



%%%%%%%%%%%%%%%%%%%%%%%%%%%%%%%%%   Comprensiones  %%%%%%%%%%%%%%%%%%%%%%%%%%%%%%%%%%%%




\section{Comprensiones} %Comprensión de generadores

Una comprensión es una forma compacta de crear una estructura de datos de Python  a partir de uno o más
iteradores, como listas, cadenas y tuplas. Las comprensiones  permiten combinar bucles y condicionales con una sintaxis
menos detallada y más concisa, como indica su nombre por comprensión.  Pueden ser útiles cuando necesite crear nuevas listas a partir de listas existentes o cualquier objeto iterable. Por ejemplo, puedes usar comprensiones de listas para crear una lista de cuadrados a partir de una lista de números. 

Python permite una construcción  llamada de comprensión de listas, que nos permite crear listas a partir de otras listas, tuplas y cadenas o cualquier objeto iterable. También hay otras comprensiones,como  comprensiones de diccionarios, conjuntos y un tipo de objeto generador.\\ 

Las listas por comprensión no siempre se presentan en cursos para principiantes, pero se han incluido aquí porque lo más probable es que se encuentre con estas listas tan pronto como empieces a ver el código de otras personas.

El uso de una comprensión a veces se toma como una señal, que conoces Python a un nivel superior al de principiante. En otras palabras, es más Pythónico.\\

En las próximas secciones se estudiaran las compresiones de lista, diccionarios, conjuntos y el tipo generador, bien empecemos con comprensión de Listas. %(Bill pag 83)\\


% 




% La sintaxis general un a comprensión de  diccionario es semejante a listas: {\ttfamily\{}clave:valor {\ttfamily for} expresión {\ttfamily in} iterable{\ttfamily \}}, observe los elementos son pares clave-valor y encerrados en llaves. En la sección \ref{compresionD} continuamos estudiando este tipo de comprensión. En la próxima sección nos enfocamos en la comprensión de listas. \noindent{\ttfamily \{key\_expresión\! :\!\!\! value\_expression for expression in iterable\}}, 




\subsection{Comprensión de Listas}
%  (y por extensión también otras colecciones que veremos más tarde \ref{extension})

La comprensión de listas en Python es un método sintáctico para crear listas a partir de los elementos de otras listas, cadenas, tuplas o cualquir objeto iterable,  con una forma rápida de escribir, muy legible y funcionalmente eficiente. Su sintaxis le permite crear listas avanzadas en forma más concisa, y en una sola línea de código. La forma más simple de comprensión de listas es:\\

{\ttfamily[} expresión {\ttfamily for} item  {\ttfamily in} iterable{\ttfamily ]}\\


Una lista por comprensión combina el bucle {\ttfamily for ... in} crea de nuevos elementos en una línea, añadiendo  automáticamente  cada nuevo elemento, de acuerdo a la expresión indicada. \\

El siguiente ejemplo crea la misma lista de números al cuadrado que desarrollamos, con una comprensión de lista:

\hspace{-2.5cm}{\ttfamily \scriptsize cuadrado.py}\vspace{-0.5cm}
\begin{lstlisting}
	cuadrados = [valor**2 for valor in range(1, 6)]
	print(cuadrados)
\end{lstlisting}

La sintaxis para la comprensión, comienza con un nombre descriptivo para la lista, como cuadrados. 
Y luego un conjunto de corchetes y se define la expresión para los valores que desea almacenar en la nueva lista. En este ejemplo la expresión es {\ttfamily valor**2}, que eleva el {\ttfamily valor} al exponente 2. Luego, escribe un  bucle {\ttfamily for} 
para generar los números que desea introducir en la expresión y cierra los corchetes. 

El  bucle {\ttfamily for}  en este ejemplo, es para extraer el valor del iterable {\ttfamily range(1, 6)},
que introduce los valores del 1 al 5 en la expresión {\ttfamily valor**2}. Observe no se utilizaron los dos puntos al final de la declaración {\ttfamily for}. El resultado es la misma lista de números cuadrados:

\vspace{0.09cm}
{\ttfamily\scriptsize\scriptsize 
	\noindent Out[]\vspace{-0.3cm}	
	{\color{blue} 	
		\begin{verbatim}
			[1, 4, 9, 16, 25]
		\end{verbatim}
}}

Se necesita un poco práctica para escribir comprensiones. Cuando escribas un código con tres o cuatro líneas  para generar listas y  comienza a sentirse repetitivo, considere escribir su propia comprensión de listas.

\subsection{Comprensión de Listas y el {\ttfamily if} }

% list comprehensions o comprensión de listas 

Una comprensión de listas también puede incluir una expresión condicional, después del objeto iterable, siguiendo la sintaxis:\\

{\ttfamily [} expresión {\ttfamily for} item  {\ttfamily in} iterable  {\ttfamily if} condición {\ttfamily ] }\\

\noindent Ilustremos con un ejemplo. Hagamos una  comprensión  de lista  con  números impares entre 1 y 10  (la expresión {\ttfamily  (numero \% 2 == 1) } es verdadero para números impares y falso para números pares)\footnote{Parte entera:  al dividir por 2, el resto es cero (par) o uno (impar)}.

\begin{lstlisting}
	In[]:
	>>> comprension = [numero for numero in  range(1,11) if  (numero % 2 == 1)]  
	>>> comprension
	[1, 3, 5, 7, 9]
\end{lstlisting}

\noindent La forma tradicional para hacer este código puede ser:

\hspace{-3cm}{\ttfamily \scriptsize tradicional.py}\vspace{-0.5cm}
\begin{lstlisting}
	lista = []
	for numero in range(1,11):
	if numero%2 == 1:
	lista.append(numero)
	print(lista)         
\end{lstlisting}

{\ttfamily\scriptsize\scriptsize\noindent Out[]:\vspace{-0.3cm}	                   %   \hspace{-0.65cm}
	{\color{blue}	
		\begin{verbatim}
			[1, 3, 5, 7, 9]
		\end{verbatim}
}}

Finalmente, así como puede hacer bucles anidados, también se pueden hacer más de un conjunto 
cláusulas de {\ttfamily for ... in}  en la comprensión correspondiente. Podemos ilustrar con un ejemplo típico, 
bucle anidado e imprimir los resultados:
% lista.append((fila,columna))	

\hspace{-2.3cm}{\ttfamily \scriptsize tipico.py}  \vspace{-0.5cm}
\begin{lstlisting}
	filas = range(1,4)
	columnas = range(1,3)			
	for fila in filas:
	for columna in columnas:
	print(fila, columna)
\end{lstlisting}
{\ttfamily\scriptsize\scriptsize 
	Out[]: \vspace{-0.3cm}	            %   \hspace{-0.65cm}
	{\color{blue} 	
		\begin{verbatim}
			1 1
			1 2
			2 1
			2 2
			3 1
			3 2		
		\end{verbatim}
}}

Construimos una comprensión de lista con tuplas;  listas de tuplas, con fila y columna como elementos (fila,columna), luego iterando en filas y columnas. Ahora usamos la compresión  de lista que llamaremos celdas, para imprimir las tuplas:

\hspace{-2.3cm}{\ttfamily \scriptsize celdas.py}  \vspace{-0.5cm}
\begin{lstlisting}		
	filas = range(1,4)
	columnas = range(1,3)
	celdas = [(fila, columna) for fila in filas for columna in columnas]
	# celdas es una lista de tuplas
	for celda in celdas:
	print(celda)
\end{lstlisting}
{\ttfamily\scriptsize\scriptsize 
	Out[]: \vspace{-0.3cm}	  %   \hspace{-0.65cm}
	{\color{blue} 	
		\begin{verbatim}
			(1, 1)
			(1, 2)
			(2, 1)
			(2, 2)
			(3, 1)
			(3, 2)	
		\end{verbatim}
}}
\noindent En este ejemplo hemos extraído las tuplas de la comprensión de lista (celdas) con un {\ttfamily for ...in}.\\


Si sólo se necesita los valores y no las tuplas; éstos se pueden ser extraer iterando sobre celdas (comprensión de lista), e indicando  los valores fila y columna como expresión:

%hspace{-2.3cm}{\ttfamily \scriptsize valores.py}  \vspace{-0.5cm}

\begin{lstlisting}
	In[]:	
	>>> for fila,columna in celdas:
	print(fila, columna)
\end{lstlisting}  
{\ttfamily\scriptsize\scriptsize 
	\noindent Out[]: \vspace{-0.3cm}	
	{\color{blue} 	
		\begin{verbatim}
			1 1
			1 2
			2 1
			2 2
			3 1
			3 2
		\end{verbatim}
}} En este caso sólo se presentan los valores, y no las tuplas.\\



\vspace{0.2cm}
\begin{minipage}[t]{0.2\textwidth}
	\centering\raisebox{\dimexpr\topskip-\height}{%
		
\includegraphics[scale=0.20]{../Logos_de_Python/Advertencia}}  %[width=\textwidth]
	
\end{minipage}%\hfill
\begin{minipage}[t]{0.70\textwidth} 
	{\scriptsize Debe tener presente que no puedes escribir cada bucle {\ttfamily for} como una comprensión de lista.\\ 
		
		El nombre {}"{}comprensión de {}"{}lista{}"{} puede ser un poco confuso porque parece sugerir que las comprensiones solo son para trabajar con listas. En realidad, la palabra {}"{}lista{}"{} en comprensiones de lista, es usada para indicar que puedes hacer iterar sobre cualquier objeto iterable en Python y el producto final será una lista.}
\end{minipage}
\vspace{0.2cm}


\subsection{Ejercicios}

\begin{enumerate}[6.6.3.1]
	
	\item Comprensión de cubos: Use una lista de comprensión para generar una lista de los primeros 10 cubos.
	
	\item Multiplos: Construya una comprensión de lista para generar los multiplos de tres entre uno y treinta. Y luego imprime la lista. 	
	%  multiplos = [i for i in range(30) if i \% 3 == 0] % resto es cero print(multiplos)
\end{enumerate}









\subsection{Comprensión de Diccionarios}\label{compresionD} 
% Comprensiones   De Bill  pag 81 .
% https://www.freecodecamp.org/espanol/news/compresion-de-diccionario-en-python-explicado-con-ejemplos/
Los diccionarios al igual que las listas también tienen comprensiones. La comprensión de diccionarios (dict comprehension) puede ser muy útil en crear nuevos diccionarios basados en diccionarios existentes, u objetos iterables, en caso dinámico, como entrada de usuario en una web.

La sintaxis general de la comprensión de diccionario se muestra a continuación:
\begin{center}
	{\ttfamily \{}clave:valor {\ttfamily for} expresión {\ttfamily in} iterable {\ttfamily\}}  
\end{center}\noindent Observe, los elementos son pares, por supuesto los {\ttfamily items} en los diccionarios son pares (clave-valor).  Así, para cada elemento del iterable se  genera un par (clave-valor) del nuevo diccionario.\\

A continuación exponemos algunos ejemplo de comprensión de diccionarios, comenzamos con un diccionario que nos indica el número veces que aparece una letra en una cadena. En el diccionario las claves serán las letras y el valor su cuenta:
\begin{lstlisting}
	In[]:
	>>> cadena = 'caracas'
	>>> cuenta = {letras:cadena.count(letras) for letras in cadena}
	>>> cuenta
	{'c': 2, 'a': 3, 'r': 1, 's': 1}
\end{lstlisting}

El método {\ttfamily count()} recibe una letra  {\ttfamily "{}x{}"{}} como parámetro de transferencia,  y cuenta cuántas veces aparece  {\ttfamily "{}x{}"{}} en la cadena. \\


Analizando el proceso de recorrido del {\ttfamily for .. in},  podemos observa que al contar la letra {\ttfamily c}
lo hace dos veces, y con respecto a la letra  {\ttfamily a} lo hace tres veces, es un costo innecesario, una forma de eliminar este trabajo es usar los conjuntos,  en un proceso mas limpio y más pythónico:

\begin{lstlisting}
	In[]:	
	>>> cadena = 'caracas'
	>>> cuenta = {letras:cadena.count(letras) for letras in set(cadena)}
	>>> cuenta
	{'r': 1, 's': 1, 'c': 2, 'a': 3}
\end{lstlisting}

Las claves del diccionario están en un orden diferente al del ejemplo anterior, porque al iterar
{\ttfamily set(cadena)} devuelve letras en un orden diferente al de iterar la palabra de la cadena.\\



Ahora consideraremos un segundo ejemplo. Esta  vez suponga dos listas, una lista de usuarios y la otra con las puntuaciones de los usuarios:
\begin{lstlisting}
	In[]:		
	>>> usuarios  = ["Alexis", "Alberto", "Carolina", "María", "Luisa"]
	>>> desempeño = [12, 15, 13, 14,  18]
\end{lstlisting}

Para construir un diccionario con los nombres de los usuarios como las claves y las puntuaciones como los valores, se puede plantear de varias formas, pero para hacerlo como una comprensión, podemos organizar primero las dos listas, en una listas de tuplas, cada tupla con dos elementos, usuario y puntuación. 
Para hacer una comprensión de diccionario, se usan los elementos de las tuplas como clave-valor respectivamente en le nuevo diccionario; luego iteramos sobre la lista de tuplas: 

\hspace{-2.5cm}{\ttfamily \scriptsize usuario.py}  \vspace{-0.5cm}
\begin{lstlisting}	
	# lista de tupas, usuario y su desempeño:
	lista = [('Alexis', 12), ('Alberto', 15), ('Carolina', 13), ('María', 14), 
	('Luisa', 18)]	
	
	# Creando un diccionario a partir de dos listas
	diccionario = {usuario:desempeño for usuario,desempeño in lista}
	print(diccionario)
\end{lstlisting}	

Aquí se obtiene el nuevo diccionario a partir de las dos listas. 

{\ttfamily\scriptsize\scriptsize 
	\hspace{-0.65cm} \vspace{-0.3cm}	
	{\color{blue} 	
		\begin{verbatim}
			{'Alexis': 12, 'Alberto': 15, 'Carolina': 13, 'María': 14, 'Luisa': 18}	
		\end{verbatim}
}}
\noindent Una comprensión de diccionario a partir de una lista de tuplas con dos elementos.\\  

% Traducido y adaptado por andres-torres
Sin embargo,  se puede hacer utilizando la función {\ttfamily zip()}, que permite construir un objeto iterable, a partir de las dos listas. Esta función comprime las dos listas en un objeto iterable tipo {\ttfamily <{}class '{}zip{}'{}>{}}, su sintaxis es simple: {\ttfamily zip(lista1,lista2)}. 

Ilustremos la comprensión de diccionario usando función  {\ttfamily zip()} con los datos del ejemplo anterior:

\begin{lstlisting}
	In[]:	
	>>> diccionario = {usuario:puntuación for usuario,puntuación in  
		zip(usuario,desempeño )}
	>>> diccionario
	{'Alexis': 12, 'Alberto': 15, 'Carolina': 13, 'María': 14, 'Luisa': 18} 
\end{lstlisting}	

En este caso la comprensión es muy Pythónico, con menos líneas de código. \\

Ahora construimos una comprensión de diccionarios a partir de un diccionario existente.
Por ejemplo: Crear un nuevo diccionario que contenga solo los elementos de un diccionario existente, que satisfagan una condición particular. La sintaxis general para esta comprensión es la siguiente:\\

{\ttfamily \{}clave:valor {\ttfamily  for}  clave:valor {\ttfamily in diccionario.items() if} condición{\ttfamily \}}\\

El método {\ttfamily items()} se usa para obtener los pares clave-valor del diccionario existente. Accedemos al primer ítem del diccionario, si cumple la condición ({\ttfamily True}).
Si la condición es True, se añade el primer ítem al nuevo  diccionario.\\

Supongamos que los usuarios en el ejemplo anterior, han solicitado una consideración para los que obtuvieron una puntuación menor a 15. Atendiendo esta solicitud, queremos construir un diccionario sólo con los usuarios sujetos a esta condición. Para sintetizar el código podemos  usar el diccionario construido con la comprensión anterior: 

\begin{lstlisting}
	In[]:	
	>>> considerados = {clave:valor for clave,valor in diccionario.items()
		if valor<15}
	>>> considerados
	{'Alexis': 12, 'Carolina': 13, 'María': 14}
\end{lstlisting}

En este caso se obtiene una comprensión de diccionario a partir de un diccionario existente, construido con la condición que los valores sean menor que 15, y se obtienen tres usuarios. Esta es una comprensión muy Pythónico y elegante. 





%  https://es.khanacademy.org/computing/computer-programming/sql
%conjuntos comprensiones
\subsection{Compresión de conjuntos}\label{conjuntosC} 

También hay comprensiones de conjuntos. La versión más simple es semejante a las comprensiones de listas  y de diccionario, y siguen la siguiente sintaxis:

\begin{center}
	\{ expresión{ } {\ttfamily for} { } expresión { } {\ttfamily in} { } iterable \} 
\end{center}

Por ejemplo, la siguiente comprensión de conjunto, genera un conjunto de números del 1 al 5: 
\begin{lstlisting}
	In[]:	
	>>> conjunto = {número for número in range(1,6)}
	>>> conjunto
	{1, 2, 3, 4, 5}
\end{lstlisting}



\subsection*{Comprensión de conjunto y el {\ttfamily if} }
En las comprensiones de conjuntos también son válidas los condicionales  {\ttfamily if} y los bucles de  {\ttfamily for}, como en listas y diccionarios.

En el siguiente ejemplo,  se muestra una comprensión de conjunto, que genera un conjunto con los múltiplos de tres, entre los números del 1 al nueve 9:  
\begin{lstlisting}
	In[]:	
	>>> conjunto = {número for número in range(1,10) if número % 3 == 0}
	>>> conjunto
	{9, 3, 6}
\end{lstlisting}



Y por último, tal vez por intuición pensaste que también haya comprensión de tuplas (las tuplas son inmutables), pero no acertaste, las tuplas no tienen comprensiones.

Sin embargo, se puede generar un objeto iterable tipo {\ttfamily  <class {}'{}generator{}'{}>{}}.



\subsection{Compresión de generador}\label{compresionG} 
% https://pywombat.com/articles/generadores-python

¡Las tuplas no tienen comprensiones!\\ 

Quizás pensaste que al cambiar los corchetes (o las llaves)  en una lista o diccionario de comprensión por  paréntesis lograría crear un comprensión de tupla, y al parecer funcionar porque no hay excepción, veamos que sucede:

\begin{lstlisting}
	In[]:	
	>>> prueba = (número for número in range(1, 6))
\end{lstlisting}

Lo que está entre paréntesis es una comprensión de generador y devuelve un objeto objeto tipo {\ttfamily  <{}class {}'{}generator{}'{}>{}}.

\begin{lstlisting}
	In[]:
	>>> type(prueba)
	<clase 'generador'>
\end{lstlisting}

Un generador es una forma para proporcionar datos a un iterador. (Es un tipo de función especial que retorna un objeto iterable). La principal ventaja de un objeto generator recae en el uso de memoria, ya que se obtienen solo los valores que necesitemos en el momento, no se reserva espacio en memoria con datos que no usaremos. Esto nos viene perfecto al trabajar con colecciones de miles, cientos de miles o inclusive millones de registros. Pudiendo  programar con un mucho mayor eficiencia.

En nuestro ejemplo podemos obtener los valores del gererador {\ttfamily prueba} con un {\ttfamily for}, veamos:


\begin{lstlisting}
	In[]:	
	>>> for x in prueba:
	print("x = ",x)
	
	x =  1
	x =  2
	x =  3
	x =  4
	x =  5
\end{lstlisting}

En este caso ya hemos usado los valores del generador, y en este momento no tiene elemento disponibles.\\

También se puede usar la función {\ttfamily list()} sobre la comprensión de generador para que funcione como
una lista de comprensión:

\begin{lstlisting}
	In[]:	
	# Nuevamente se contruye objeto prueba
	>>> prueba = (número for número in range(1, 6))
	>>> lista_números = list(prueba)
	>>> lista_números
	[1, 2, 3, 4, 5]
\end{lstlisting}

Observe, hemos calculado nuevamente el objeto prueba; si no lo calculamos de nuevo, el resultado  será una lista vacía, 
esto se debe a que {\bf un generador sólo puede funcionar una vez}. Los objetos como Listas, conjuntos, cadenas y diccionarios existen en la memoria, pero un generador crea sus valores sobre la marcha y los entrega a través de un iterador una sola vez; 
este objeto  no recuerda, por lo que no puede reiniciar ni hacer una copia de seguridad de un generador.
Si intenta repetir este generador, encontrará que está agotado:

\begin{lstlisting}
	In[]:
	>>> intentar_otra vez = lista(prueba)
	>>> intentar_otra vez
	[]
\end{lstlisting}

Puedes crear un generador a partir de una comprensión del generador, como hicimos aquí, pero también se puede crear como una función generadora. Más adelante en la sección \ref{generadorF} continuamos con este tema de gereadores,  un  especial de funciones generadoras.\\

En las comprensiones de generadores también son válidas los condicionales  {\ttfamily if} y los bucles de  {\ttfamily for}, como en listas y diccionarios.

\begin{lstlisting}
	prueba = (número*2 for número in range(1, 11) if número %2 == 0 )
	
	prueba = list(número*2 for número in range(1, 11) if número %2 == 0 )
\end{lstlisting}

Esta comprensión  genera los cuadrados de los números pares del 1 al 10. Pruébelo. Y usando la función {\ttfamily list()}, 
se construye una lista con el tipo  generador, que viene a ser una comprensión de lista. Pruébelo y construya un test.  


% Cada vez que itera a través de un generador, realiza un seguimiento de dónde estaba la última vez fue llamado y devuelve el siguiente valor. Esto es diferente de una función normal, que no tiene memoria de llamadas anteriores y siempre comienza en su primera línea con el mismo estado. Si desea crear una secuencia potencialmente grande y el código es demasiado grande para un generador comprensión, escriba una función generadora. Es una función normal, pero devuelve su valor. con una declaración de rendimiento en lugar de retorno. Escribamos nuestra propia versión de range():

\subsection{Ejercicos}

%  yield  rendimiento 
%   retorno


\section{Anidación de estructuras}
% Pendientes Diccionarios, listas,  sentencia declaración   {\ttfamily if}

Hasta ahora hemos estudiado datos desde los valores más simples como booleanos, números y cadenas hasta estructuras más complejas como tuplas, listas, diccionarios y conjuntos; con esta nociones podemos combinar estas estructuras de datos y construir nuestras propias estructuras  más grandes y complejas. Por ejemplo, si   necesitamos  almacenar varios diccionarios en una lista (una lista de  diccionarios), o una lista de elementos dentro de un diccionario, o incluso un diccionario dentro otro diccionario, y por supuesto las lista de listas o lista de tupla. A este procedimiento se le llama anidación, en este caso es una anidación de estructuras de datos; pero también hay    
anidación de estructuras de control de flujo, y por supuestos con importantes aplicaciones.  Comencemos con anidación de estructuras
de datos.


\subsection{Anidación de estructuras de datos}

En esta sección estudiaremos anidación de estructura de datos que involucran diccionarios, lista, conjuntos,  y tupla. 
Si se trata de anidar diccionarios es importante que tengan una estructura idéntica, de tal forma que facilite trabajar con cada elemento del diccionario y recorrer lista cuando sea necesario. Una aplicación de este proceso de anidación se puede encontrar en un juego de extraterrestres, donde todas las características cada alienígena o extraterrestre se describe en un diccionario, digamos color,  velocidad, tamaño, puntuación,  etc;  luego se pueden almacenados en una lista.\\  


Generalmente es conveniente almacenar varios diccionarios en una lista, cuando cada diccionario contiene varias características (información) sobre un objeto. Por ejemplo, para un sitio web, se puede crear un diccionario para cada usuario que  almacene  su información. Luego almacenar estos diccionarios en una lista llamada usuarios. En este caso es natural elegir la clave como el nombre del usuario que debe ser único en el sitio web.\\



Construyamos algunas estructuras de datos para fijar ideas. Consideremos tres listas con la información de los tres estudiantes, y comencemos con una tuplas de listas:

\hspace{-3cm}{\ttfamily \scriptsize estudiantes.py}  \vspace{-0.5cm}
\begin{lstlisting}
	Antonio_Sarco = ['29.231.253', 'Mérida', 'Arquitectura', 'Colombia']
	María_Jiménez = ['19.153.654', 'Zulia', 'Medicina', 'Venzuela']
	Carlos_Dávila = ['22.562.148', 'Sucre', 'Estadística', 'España']
	
	# Tupla de listas
	tupla_de_listas = (Antonio_Sarco, María_Jiménez, Carlos_Dávila)
	print(tupla_de_listas)
\end{lstlisting}	
{\ttfamily\scriptsize\scriptsize 
	\noindent Out[]: \vspace{-0.3cm}	
	{\color{blue} 	
		\begin{verbatim}
			(['29.231.253', 'Mérida', 'Arquitectura', 'Colombia'], 
			['19.153.654', 'Zulia', 'Medicina', 'Venzuela'],
			['22.562.148', 'Sucre', 'Estadística', 'España'])
		\end{verbatim}
}}
Este objeto es una tupla, pero una tupla de listas, que contiene la información los estudiantes.\\

Ahora construimos una lista que contenga listas, las tres listas contenidas en una lista:
\begin{lstlisting}
	In[]:
	>>> lista_de_listas = [Antonio_Sarco, María_Jiménez, Carlos_Dávila ]
	>>>lista_de_listas
	
	[['29.231.253', 'Mérida', 'Arquitectura', 'Colombia'], 
	['19.153.654', 'Zulia', 'Medicina', 'Venzuela'],
	['22.562.148', 'Sucre', 'Estadística', 'España']]
\end{lstlisting}

Una lista de listas.\\


Seguimos un diccionario de listas. En este ejemplo, usaremos el nombre del estudiante como clave y la lista como valor:

\begin{lstlisting}
	In[]:
	>>> dicci_de_listas = {"Antonio_Sarco": Antonio_Sarco,
		"María_Jiménez": María_Jiménez,
		"Carlos_Dávila":Carlos_Dávila }
	>>> dicci_de_listas
	
	{'Antonio_Sarco': ['29.231.253', 'Mérida', 'Arquitectura', 'Colombia'], 
		'María_Jiménez': ['19.153.654', 'Zulia', 'Medicina', 'Venzuela'],
		'Carlos_Dávila': ['22.562.148', 'Sucre', 'Estadística', 'España']}
\end{lstlisting}



Y finalmente,  un diccionario de conjuntos. Para este ejemplo consideremos un diccionario llamado bebidas, donde  el nombre de la bebida será la clave del diccionario, y el valor correspondiente un conjunto con los ingredientes de cada bebida:
% queso, jamón, huevo, lechuga, cebolla, tomate, salsa, mostaza, tomate, mayonesa,

%\hspace{-2.5cm}{\ttfamily \scriptsize bebidas.py}  \vspace{-0.5cm}
\begin{lstlisting}
	In[]:
	>>> bebidas = {'martini': {'vodka', 'vermouth'},
		'black russian': {'vodka', 'kahlua'},
		'white russian': {'crema', 'kahlua', 'vodka'},
		'manhattan': {'rye', 'vermouth', 'bitters'},
		"screwdriverr": {'jugo de naranja', 'vodka'} }
\end{lstlisting}

Ambas estructuras están encerrados entre llaves (\{  \}), pero el conjunto es sólo una secuencia de valores,
mientras que el diccionario son uno o varios pares clave-valor.\\

Con un {\ttfamily for ... in}, y aplicado el método {\ttfamily item()} al diccionario, podemos extraer los pares; llamando a las claves nombre y a los valores ingredientes:

\begin{lstlisting}
	>>> for nombre, ingredientes in bebidas.items():
	print(ingredientes)
	
	{'vermouth', 'vodka'}
	{'kahlua', 'vodka'}
	{'crema', 'kahlua', 'vodka'}
	{'vermouth', 'bitters', 'rye'}
	{'jugo de naranja', 'vodka'}			
\end{lstlisting}

\noindent Este código muestra los valores del diccionario bebidas, que son conjuntos.\\


Las únicas limitaciones para construir estructuras de datos como éstas, serán los propios tipos de datos. Por ejemplo, en el diccionario las claves deben ser inmutables, por lo que una lista, un diccionario o un conjunto no pueden ser claves para otro diccionario.
Una  tupla se puede ser clave, por ejemplo, podrías indexar sitios de interés mediante coordenadas GPS, latitud, longitud y altitud:

\begin{lstlisting}
	>>> casas = {(-77.1434, 8.5952, 2650):'Mi Casa',
		(38.89, -77.03, 13):'La Casa Blanca'}
\end{lstlisting}

Ya puedes construir sus propias estructuras de datos, crearlas las estructuras necesarias  para incluirlas en sus programas.

\subsection{Anidación de estructuras de flujo}

En proceso....\\

En lenguaje Python, la toma de decisiones se basa en sentencias condicionales.
Existen cuatro formas: {\ttfamily if}, {\ttfamily if-else},   {\ttfamily if-elif-else} y anidadas.



Estructuras  {\ttfamily if} anidadas:

Cuando una sentencia {\ttfamily if} o {\ttfamily if-else} está presente dentro de otro bloque de sentencia {\ttfamily if} o {\ttfamily if-else}, se trata de una sentencia  {\ttfamily if} anidada. Esta situación se produce cuando es necesario filtrar una variable varias veces y comprobar varias condiciones.

En una sentencia  {\ttfamily if} anidada, la indentación que permite definir el alcance de cada sentencia es esencial. El número de anidaciones posibles es ilimitado, pero cada una de ellas reduce la optimización del programa y lo hace más complejo de leer y entender. Por lo tanto, es mejor reducir al mínimo el número de anidaciones.



Pueden construirse pruebas anidadas con {\ttfamily if}, unas dentro de otras, con tantos niveles de profundidad como sea necesario. En la sección \ref{prueba} se exponen los ejemplo con pruebas anidadas. Por ahora nos enfocamos en algunas test o pruebas comparativa útiles en los condicionales.  


% (https://pypi.python.org/pypi)

\section{Algunas Aplicaciones Prácticas}\label{aplicaciones}
%   :listas, diccionarios, condicionales, bucles
%   De  Eric
Ya estamos preparados para escribir pequeños programas que involucren las herramientas que hemos tratado, como estructuras de datos, sentencias condicionales y bucles. El juego de extraterrestre es una buena ilustración para aplicar estas herramientas. 

En este juego los personajes son alienígenas, que se describen mediante un diccionario, que contiene toda la información sobre el mismo.  Sin embargo, no hay forma para almacenar la información sobre un segundo extraterrestre, y mucho menos sobre un pantalla llena de extraterrestres.  Aquí es donde entran en acción las listas de extraterrestres y diccionarios.

Otras aplicaciones que introducimos son las bebidas, hamburguesas o pizzas y xxxx.




\subsection{Un Juego de Extraterrestres}

En el juego de extraterrestre o alienígenas,  cada alienígena se describe mediante un diccionario {\ttfamily alien\_0i}, que contiene toda la información sobre el mismo, pero no tiene espacio para almacenar información sobre un segundo extraterrestre, y mucho menos sobre una pantalla llena de extraterrestres. 

Para gestionar una flota de extraterrestres, consideraremos una lista de extraterrestres ({\bf listas de diccionario}),  cada extraterrestre es un diccionario con  la información del alienígena.

Para fijar ideas comencemos con una lista con tres extraterrestres, tres diccionarios; pero antes debemos  crear los tres diccionarios, cada uno de los cuales representa un extraterrestre diferente y los almacenamos en una lista llamada {\ttfamily aliens}:

\hspace{-2.5cm}{\ttfamily \scriptsize aliens.py}  \vspace{-0.5cm}
\begin{lstlisting}
	alien_00 = {'color': 'verde', 'puntos': 5}
	alien_01 = {'color': 'amarillo', 'puntos': 10}
	alien_02 = {'color': 'rojo', 'puntos': 15}
	
	# Lista de diccionarios
	aliens = [alien_00, alien_01, alien_02]
	for alien in aliens:
	print(alien)
\end{lstlisting}

Finalmente, con un bucle {\ttfamily for} imprimimos cada  elemento de la lista:

{\ttfamily\scriptsize\scriptsize \noindent Out[]: \vspace{-0.3cm}	
	{\color{blue} 	
		\begin{verbatim}
			{'color': 'verde', 'puntos': 5}
			{'color': 'amarillo', 'puntos': 10}
			{'color': 'rojo', 'puntos': 15}			
		\end{verbatim}
}}

Un ejemplo más realista involucra más de tres extraterrestres con código que genera automáticamente cada extraterrestre. En el siguiente ejemplo, usamos la función {\ttfamily range()} para crear una flota de 30 alienígenas, observe:

\hspace{-2.5cm}{\ttfamily \scriptsize aliens.py}  \vspace{-0.5cm}
\begin{lstlisting}	
	# lista vacía para almacenar extraterrestres.
	extraterrestres = []
	# 30 alienígenas verdes.
	for nuevo_alien in range(30):
	nuevo_alien = {'color': 'verde', 'puntos': 5, 'velocidad': 'lento'}
	extraterrestres.append(nuevo_alien)
	# Muestra los primeros 5 extraterrestres.
	for extraterrestre in extraterrestres[:5]:
	print(extraterrestre)
\end{lstlisting}	

{\ttfamily\scriptsize\scriptsize \noindent Out[]: \vspace{-0.3cm}	
	{\color{blue}	
		\begin{verbatim}
			{'color': 'verde', 'puntos': 5, 'velocidad': 'lento'}
			{'color': 'verde', 'puntos': 5, 'velocidad': 'lento'}
			{'color': 'verde', 'puntos': 5, 'velocidad': 'lento'}
			{'color': 'verde', 'puntos': 5, 'velocidad': 'lento'}
			{'color': 'verde', 'puntos': 5, 'velocidad': 'lento'}
		\end{verbatim}
}}

Este ejemplo comienza con una lista vacía para contener todos los extraterrestres que se crearan. La función {\ttfamily range(k)} devuelve una serie de números de {\ttfamily 0} a {\ttfamily k}, que en este caso la usamos para indicarle a Python cuántas veces queremos que se repita el ciclo. Cada vez que el bucle se ejecuta, creamos un nuevo alienígena  y lo agregamos a la lista extraterrestres. Usamos un corte para imprimir los primeros cinco extraterrestres, y finalmente imprimimos la longitud de la lista para demostrar que realmente hemos generado una flota completa de 30 extraterrestres:

% # Número extraterrestres que se han creado.
\hspace{-2.5cm}{\ttfamily \scriptsize }  \vspace{-0.5cm}
\begin{lstlisting}	
	In[]:
	>>> print(f"Número total de extraterrestres: ",len(extraterrestres))
\end{lstlisting}	
{\ttfamily\scriptsize\scriptsize \noindent Out[]: \vspace{-0.3cm}	
	{\color{blue} 	
		\begin{verbatim}
			Número total de extraterrestres: 30
		\end{verbatim}
}}


Todos estos alienígenas tienen las mismas características, pero Python considera cada uno
como un objeto separado, lo que nos permite modificar cada alienígena individualmente.

Ahora supongamos que los extraterrestres cambian de aspecto, digamos que cambian a color, y se mueven más rápido 
a medida que avanzan, y por supuesto también sube o baja la puntuación. Cuando llega el momento de hacer los  cambios, podemos usar un bucle {\ttfamily for} y una declaración {\ttfamily if} para cambiar el color de los extraterrestres. Por ejemplo, para cambiar los tres primeros alienígenas de verde a amarillo, con 10 puntos y con velocidad media cada uno,
podríamos hacerlo como sigue:

\hspace{-2.5cm}{\ttfamily \scriptsize aliens.py}  \vspace{-0.5cm}
\begin{lstlisting}	
	#  lista vacía para almacenar extraterrestres.
	extraterrestres = []
	# 30 extraterrestre verdes.
	for extraterrestre in range(30):
	nuevo_extraterrestre = {'color': 'verde', 'puntos': 5, 'velocidad': 'lento'}
	extraterrestres.append(nuevo_extraterrestre)
	for extraterrestre in extraterrestres[:3]:
	if extraterrestre['color'] == 'verde':
	extraterrestre['color'] = 'amarillo' # cambia color
	extraterrestre['velocidad'] = 'media'
	extraterrestre['puntos'] = 10
	
	# Muestra los primeros 5 extraterrestres.
	for extraterrestre in extraterrestres[:5]:
	print(extraterrestre)
\end{lstlisting}	

Aquí se modificaron los primeros tres extraterrestre, recorremos una porción de la lista, que incluye solo los primeros tres elementos. Todos los extraterrestre son verdes al inicio.  Pero eso no siempre será el caso, por lo que escribimos una declaración {\ttfamily if} para asegurarnos sólo estamos modificando extraterrestres verdes. Si el alienígena es verde, le cambiamos el color a 'amarillo', la velocidad a 'media' y el valor de puntos a 10, como se muestra en la salida:

{\ttfamily\scriptsize\scriptsize \noindent Out[]:\vspace{-0.3cm}	
	{\color{blue} 	
		\begin{verbatim}
			{'color': 'amarillo', 'puntos': 10, 'velocidad': 'media'}
			{'color': 'amarillo', 'puntos': 10, 'velocidad': 'media'}
			{'color': 'amarillo', 'puntos': 10, 'velocidad': 'media'}
			{'color': 'verde', 'puntos': 5, 'velocidad': 'lento'}
			{'color': 'verde', 'puntos': 5, 'velocidad': 'lento'}
		\end{verbatim}
}}

Podrías expandir este bucle agregando un bloque {\ttfamily elif} de extraterrestre  rojos se  vuelvan amarillo, 
y que se muevan rápido con un valor 15 puntos cada uno. Nuevamente sin mostrar el resultado completo, ese bucle se vería así:

\hspace{-2.5cm}{\ttfamily \scriptsize aliens.py} \vspace{-0.5cm}
\begin{lstlisting}
	for extraterrestre in extraterrestres[:3]:
	if  extraterrestre['color'] == 'verde':
	extraterrestre['color'] = 'amarillo'
	extraterrestre['velocidad'] = 'media'
	extraterrestre['puntos'] = 10
	elif extraterrestre['color'] == 'amarillo':
	extraterrestre['color'] = 'rojo'
	extraterrestre['velocidad'] = 'rápido'
	extraterrestre['puntos'] = 15
	
	# Muestra los primeros 5 extraterrestres.
	for extraterrestre in extraterrestres[:5]:
	print(extraterrestre)
\end{lstlisting}	

{\ttfamily\scriptsize\scriptsize \noindent Out[]: \vspace{-0.3cm}	
	{\color{blue} 	
		\begin{verbatim}
			{'color': 'amarillo', 'puntos': 10, 'velocidad': 'media'}
			{'color': 'amarillo', 'puntos': 10, 'velocidad': 'media'}
			{'color': 'amarillo', 'puntos': 10, 'velocidad': 'media'}
			{'color': 'verde', 'puntos': 5, 'velocidad': 'lento'}
			{'color': 'verde', 'puntos': 5, 'velocidad': 'lento'}
			...	
		\end{verbatim}
}}

Claro se supone que hay tres rojos, y cambian a amarillo.\\





Es común almacenar varios diccionarios en una lista cuando cada diccionario contiene muchos tipos de
información sobre un objeto. Por ejemplo, puedes crear un diccionario para cada usuario en un sitio web, 
que almacene los diccionarios individuales en una lista llamada usuarios. 

Todo de los diccionarios en la lista deben tener una estructura idéntica, para poder recorrer la lista y trabaje con cada objeto del diccionario de la misma manera.\\

%  \subsection{Conjuntos listas tuplas diccionarios, in , for ... in , if}  


En el siguiente ejemplo, se usa un diccionario cuyos valores son conjuntos de licores, y las claves es el nombre de la bebida preparada con esos tipos de licores.  

\subsection{Ejemplos de Bebidas}

% Este ejemplo se puede hacer con pizza o con hamburguesa,  agencia de viajes, ubicacion de sitio para construir viviendas, etec...  Recomendaciones de acuerdo con sus exigencias 

% (Tenga en cuenta que este tema del uso de {\ttfamily for},  {\ttfamily if},  {\ttfamily and}, y  {\ttfamily or} se estudiara  en el próximo capítulo \ref{condicionales}  redactar mejorar ...).
En este ejemplo presentamos un diccionario cuyos valores son conjuntos y aplicaremos la herramientas sobre conjuntos.

Consideremos nuevamente el ejemplo de las bebidas, y suponga que estamos interesados en saber: ¿Qué bebidas contienen vodka? 

\hspace{-2.5cm}{\ttfamily \scriptsize bebidas.py}  \vspace{-0.5cm}
\begin{lstlisting}
	bebidas = {'martini': {'vodka', 'vermouth'},
		'black russian': {'vodka', 'kahlua'},
		'white russian': {'crema', 'kahlua', 'vodka'},
		'manhattan': {'rye', 'vermouth', 'bitters'},
		"screwdriverr": {'jugo de naranja', 'vodka'} }
	
	for nombre, ingredientes in bebidas.items():
	if 'vodka' in ingredientes:
	print(nombre)
\end{lstlisting}

{\ttfamily\scriptsize\scriptsize \noindent Out[]: \vspace{-0.3cm}	
	{\color{blue} 	
		\begin{verbatim}
			martini
			black russian
			white russian
			screwdriverr
		\end{verbatim}
}}

\noindent Como resultado, cuatro de la bebidas contienen Vodka, con excepción de manhattan.  


Ahora, supongamos que se quiere combinar sabores o ingredientes. Por ejemplo, supongamos que quieres encontrar alguna bebida que tenga jugo de naranja o vermouth.

\begin{lstlisting}
	In[]:
	for nombre, ingredientes in bebidas.items():
	if ingredientes & {'vermouth', 'jugo de naranja'}:
	print(nombre)
\end{lstlisting}

{\ttfamily\scriptsize\scriptsize \noindent Out[]: \vspace{-0.3cm}	
	{\color{blue} 	
		\begin{verbatim}
			martini
			manhattan
			screwdriverr
		\end{verbatim}
}}

En este ejemplo el operador intersección \& de conjuntos se usa como una expresión booleana (test o prueba condicional), y se toman las bebidas que contienen al menos un ingrediente de los exigidos. Si ninguno de esos ingredientes están disponibles, la operación \& devuelve un conjunto vacío, que se considera Falso.\\


Otra situación interesante, suponga se desea una  bebida con vodka, pero sin crema y sin vermouth; está exigencia se escribe como una expresión booleana de ésta forma:  {\ttfamily  ('vodka' in ingredientes and not (ingredientes \& \{'vermouth', 'crema'\}))},  y se codifica: 

\hspace{-2.5cm}{\ttfamily \scriptsize bebidas.py}  \vspace{-0.5cm}
\begin{lstlisting}
	# Continuando el código anterior
	for nombre, ingredientes in bebidas.items():
	if 'vodka' in ingredientes and not (ingredientes & {'vermouth', 'crema'}):
	print(nombre)
\end{lstlisting}

{\ttfamily\scriptsize\scriptsize \noindent Out[]: \vspace{-0.3cm}	
	{\color{blue} 	
		\begin{verbatim}
			black russian
			screwdriverr
		\end{verbatim}
}}

Se obtienen dos bebidas con estas exigencias. En este ejemplo se aplicó operaciones con conjuntos y las expresiones booleana,  el operador intersección  {\ttfamily \& }, y la negación {\ttfamily and not}.\\



A Continuación ilustramos algunas  operaciones con conjuntos aplicadas en este ejemplo de bebidas,  que considero  
son importante conocer su utilidad. 

Primero extraemos los conjuntos de ingredientes de las dos bebidas  black russian y  white russian,  en dos variables, digamos:

\begin{lstlisting}
	In[]:
	>>> black_r = bebidas['black russian']
	>>> white_r = bebidas['white russian']
\end{lstlisting}


Ahora con los objetos   {\ttfamily  black\_r}  y  {\ttfamily white\_r} que son dos conjuntos de ingredientes de las bebidas; calculamos los ingrediente comunes en ambos: 

\begin{lstlisting}
	In[]:
	>>> conjunto = black_r &  white_r
	>>> conjunto
	{'kahlua', 'vodka'}
\end{lstlisting}
Se obtiene los elementos comunes con la intersección.\\


y calculamos la unión, que se puede obtener en dos formas: 

\begin{lstlisting}
	In[]:
	>>> black_r |  white_r      # Primera forma
	{'kahlua', 'crema', 'vodka'}
	
	>>> black_r.union(white_r)  # Segunda forma
	{'kahlua', 'crema', 'vodka'}
\end{lstlisting}


Y apliquemos la diferencia y diferencia simétrica:

\begin{lstlisting}
	>>> white_r - black_r
	{'crema'}
	>>> white_r.difference(black_r)
	{'crema'}
\end{lstlisting}


Y  aplicado para encontrar los ingredientes exclusivos en las dos bebidas rusas:
\begin{lstlisting}
	>>> black_r ^ white_r
	{'crema'}
	>>> black_r.symmetric_difference(white_r)
	{'crema'}	
\end{lstlisting}

En este caso se obtiene crema.\\



En el siguiente ejemplo ilustraremos una lista en un diccionario, se trata de ordenar una pizza,  donde interviene el tamaño de pizza y los ingredientes. 


\subsection{Ordenar pizzas}

En el ejemplo del juego de extraterrestres usamos una lista de diccionarios. Ahora ilustraremos una lista en un diccionario, con un ejemplo que pide una pizza por tamaño y selecciona los ingredientes.  

Planteamos el problema  de cómo podría describir una pizza que alguien está ordenando. Supongamos que tiene posibilidad de elegir el tamaño de la pizza, y los ingredientes disponibles que le ofrecen.  En este caso podemos construir una lista con los tamaños de pizza y una segunda lista con los ingredientes; y construimos un  diccionario que almacenará dos tipos de información para cada pizza: el tamaño de pizza y una lista con ingredientes, ambos seleccionados por el cliente. Aquí tenemos un ejemplo de {\bf listas en diccionario}, en este caso las listas son los valores.


\begin{lstlisting}
	tamaño = ['pequeña','mediana', 'grande']
	ingredientes'= ['jamón','queso','champiñones','anchoas','pimientos verdes', 
	'cebollas']
	pizza = {'tamaño':['pequeña','mediana', 'grande'],
		'ingredientes':['jamón','queso','champiñones',
		'anchoas','pimientos verdes', 'cebollas'],} 
\end{lstlisting}

Además, le permitimos hacer varios pedidos, y le mostramos al final el pedido.
Usamos un {\ttfamily while} para mantener abierto los pedidos de pizzas, y seleccionar el tamaño, además con un {\ttfamily for} agregamos los ingrediente que el cliente desea. Iniciamos con un diccionario vacío  y dos listas vacías, para llenar con lo que el cliente desee. El {\ttfamily orden} un diccionario para almacena el pedido, con clave como tamaño se seccionado por el cliente, y el valor correspondiente a los ingredientes que también es selecciona el cliente.  Aquí tenemos un ejemplo de  {\bf listas en un diccionario}.   Y las ordenes de pizza se almacenan en la lista {\ttfamily ordenes}, ejemplo de {\bf lista de diccionarios}.

\hspace{-2.5cm}{\ttfamily \scriptsize pizza.py}  \vspace{-0.5cm}
\begin{lstlisting}
	orden ={}
	ingrediente = []
	ordenes = []
	activo = True
	pizza = {'tamaño':['pequeña','mediana', 'grande'],
		'ingredientes':['jamón','queso','champiñones',
		'anchoas','pimientos verdes', 'cebollas'],}
	while activo:
	i = 0
	tamaño = input(f"Qué tamaño desea {pizza['tamaño']} ")
	for ingredient in pizza['ingredientes']:
	ingrediente_orde = input(f"Ingredientes que desea \
	{pizza['ingredientes'][i]}, si/no ")
	if ingrediente_orde == "si":
	ingrediente.append(pizza['ingredientes'][i])
	i += 1
	orden = {"tamaño":tamaño.title(),"ingrediente":ingrediente}
	activo = input('\nDesea otra pizza si/no ' )
	if activo == "si":
	activo = True
	else:
	activo = False
	ordenes.append(orden)
	ingrediente = []
	
	print("\n","-*"*10,"Lista de su pedido","-*"*10,"\n")
	for orden in ordenes:
	print(orden)    
\end{lstlisting}

Al ejecutar el programa, le solicita el tamaño de la pizza que desea, y los ingredientes que prefiere, y antes de salir se le pregunta si quiere ordenar otra pizza; en caso de ser afirmativo se inicia nuevamente el proceso de seleccionar la pizza, en caso de ser negativo, sale y le muestra su pedido:


{\ttfamily\scriptsize\scriptsize \noindent Out[]: \vspace{-0.3cm}	
	{\color{blue} 	
		\begin{verbatim}
			Qué tamaño desea ['pequeña', 'mediana', 'grande'] Pequeño
			Ingredientes que desea jamón, si/no si
			Ingredientes que desea queso, si/no si
			Ingredientes que desea champiñones, si/no si
			Ingredientes que desea anchoas, si/no no
			Ingredientes que desea pimientos verdes, si/no no
			Ingredientes que desea cebollas, si/no no
			
			Desea otra pizza si/no si
			Qué tamaño desea ['pequeña', 'mediana', 'grande'] Mediana
			Ingredientes que desea jamón, si/no si
			Ingredientes que desea queso, si/no si
			Ingredientes que desea champiñones, si/no no
			Ingredientes que desea anchoas, si/no no
			Ingredientes que desea pimientos verdes, si/no si
			Ingredientes que desea cebollas, si/no no
			
			Desea otra pizza si/no no
			
			-*-*-*-*-*-*-*-*-*-* Lista de ordenes -*-*-*-*-*-*-*-*-*-* 
			
			{'tamaño': 'Pequeño', 'ingrediente': ['jamón', 'queso', 'champiñones']}
			{'tamaño': 'Mediana', 'ingrediente': ['jamón', 'queso', 'pimientos verdes']}		
		\end{verbatim}
}}

Es este ejemplo el cliente solicito dos pizzas diferente, una pequeña y mediana con sus ingredientes, y se muestra el pedido en dos diccionarios.\\


En éste caso un programa interactivo, donde el usuario puede responder a sus gustos.




\vspace{0.2cm}
\begin{minipage}[t]{0.2\textwidth}
	\centering\raisebox{\dimexpr\topskip-\height}{%
		
\includegraphics[scale=0.20]{../Logos_de_Python/Busquedas_R}}  %[width=\textwidth]
\end{minipage}%\hfill
\begin{minipage}[t]{0.70\textwidth} 
	{\scriptsize 
		No debe anidar listas y diccionarios demasiado profundas. Si estás anidando artículos mucho
		más profundo que lo que ves en los ejemplos anteriores, o si estás trabajando con algún
		código de otra persona con niveles significativos de anidamiento, lo más probable es que exista una forma más sencilla de resolver
		el problema.}
\end{minipage}
\vspace{0.2cm}



\subsection{Ejemplo de bodega}

Este ejemplo consiste la presentación de una bodega con tres departamentos de artículos disponibles, víveres, frutería y quincallería, cada uno con sus artículos y su presentación.  
En este caso las claves del diccionario son los nombres de los departamentos y los valores son diccionarios con los productos y su presentación, los nombres de los  productos son las claves del diccionario interno y los valores son listas de presentación de cada articulo, este es un ejemplo de {\bf diccionarios en diccionario}:

\begin{lstlisting}
	bodega = {"vívere":{"aceite":["1/2 l","1 l", "1 caja"],
			"arroz":["1/2 kg","1 kg", "1 bulto"],
			"pasta":["1/2 kg","1 Kg", "1 caja"]}, 
		"frutas":{"cambur":["verde", "maduro"],
			"lechosa":["mediana", "grande"],
			"naranja":["criolla", "california"] },
		"quincalla":{"tijeras":["punta roma","punta fina"],
			"cintas":["delgada", "mediana", "gruesa"],
			"bombillos":["5w","9w","15w"],  }  }
\end{lstlisting}

Podemos extraer la información de cualquier departamento, conociendo el nombre del departamento;  o con un bucle  {\ttfamily for} recorremos el diccionario configurado con el método {\ttfamily item()}, y extraemos los artículos, luego con un segundo bucle {\ttfamily for} recorremos el diccionario interno igualmente configurado con el método {\ttfamily item()} y  extraemos los artículos y sus  presentaciones, que son las listas:

\hspace{-2.5cm}{\ttfamily \scriptsize bodega.py}  \vspace{-0.5cm}
\begin{lstlisting}
	for depart , valor in bodega.items():
	print("\nDepartamento de :", depart.title(), "Artículos:")
	for cla, val in valor.items():
	print(cla.title())
	print(val)
	print("*-"*20)    
\end{lstlisting}
Al ejecutar obtenemos por departamento los artículos y sus presentaciones:

{\ttfamily\scriptsize\scriptsize \noindent Out[]: \vspace{-0.3cm}	
	{\color{blue} 	
		\begin{verbatim}
			Departamento de : Vívere Articulos:
			Aceite
			['1/2 l', '1 l', '1 caja']
			Arroz
			['1/2 kg', '1 kg', '1 bulto']
			Pasta
			['1/2 kg', '1 Kg', '1 caja']
			
			Departamento de : Frutas Articulos:
			Cambur
			['verde', 'maduro']
			Lechosa
			['mediana', 'grande']
			Naranja
			['criolla', 'california']
			
			Departamento de : Quincalla Articulos:
			Tijeras
			['punta roma', 'punta fina']
			Cintas
			['delgada', 'mediana', 'gruesa']
			Bombillos
			['5w', '9w', '15w']
		\end{verbatim}
}}



Observe que la estructura del diccionario de cada artículo es semejante.  Sin embrego Python no lo requiere, pero esta organización hace que los diccionarios anidados sean más fáciles de trabajarlos.\\



Hasta ahora, todos los ejemplos de código Python que hemos desarrollado han sido bloque o líneas de código para tareas específicas. Sin embargo si necesitamos hacer una misma tarea varias veces durante la ejecución del programa deberíamos escribir el código las veces que se necesite, esto se convierte en  una situación incomoda; Python  tiene  una  forma para organizar los bloques código en partes manejables, llamada funciones. Ésta estructura se estudia en el siguiente capítulo.



